\documentclass[12pt,a4paper]{report}
%DIF LATEXDIFF DIFFERENCE FILE
%DIF DEL /tmp/tmp.3qsVXPFBeX   Wed Mar 18 23:15:25 2026
%DIF ADD /tmp/tmp.MoDP70m5Sb   Wed Mar 18 23:15:26 2026

\usepackage[utf8]{inputenc}
\usepackage[T1]{fontenc}
\usepackage{amsmath}
\usepackage{amssymb}
\usepackage{amsthm}
\usepackage{mathtools}
\usepackage{physics}
\usepackage{hyperref}
\usepackage{geometry}
\usepackage{graphicx}
\usepackage{cite}
\usepackage{amscd}             \usepackage{ccicons}          \usepackage{titlesec}         \usepackage{tikz}             

\geometry{margin=1in}

\titleclass{\chapter}{straight}
\titleformat{\chapter}[hang]{\normalsize\bfseries}{\thechapter.}{0.5em}{}
\titlespacing*{\chapter}{0pt}{1.5ex plus 0.5ex minus 0.2ex}{0.8ex plus 0.2ex}

\titleformat{\section}[hang]{\normalsize\bfseries}{\thesection}{0.5em}{}
\titlespacing*{\section}{0pt}{1.0ex plus 0.3ex minus 0.2ex}{0.5ex plus 0.1ex}

\titleformat{\subsection}[hang]{\normalsize\bfseries}{\thesubsection}{0.5em}{}
\titlespacing*{\subsection}{0pt}{0.8ex plus 0.3ex minus 0.1ex}{0.4ex plus 0.1ex}

\newtheorem{theorem}{Theorem}[chapter]
\newtheorem{lemma}[theorem]{Lemma}
\newtheorem{proposition}[theorem]{Proposition}
\newtheorem{corollary}[theorem]{Corollary}
\newtheorem{definition}[theorem]{Definition}
\newtheorem{example}[theorem]{Example}
\newtheorem{remark}[theorem]{Remark}
\newtheorem{conjecture}[theorem]{Conjecture}

\newcommand{\pp}{\mathfrak{p}}    

\title{Quantum Observable Universe}
\author{Frederick Carl Leitner\thanks{World Health Organization}}
\date{\today}

\hypersetup{
  pdftitle  = {Quantum Observable Universe},
  pdfauthor = {Frederick Carl Leitner},
  colorlinks = true,
  linkcolor  = blue,
  citecolor  = blue,
  urlcolor   = blue,
}
%DIF PREAMBLE EXTENSION ADDED BY LATEXDIFF
%DIF CFONT PREAMBLE %DIF PREAMBLE
\RequirePackage{color}\definecolor{RED}{rgb}{1,0,0}\definecolor{BLUE}{rgb}{0,0.5,0} %DIF PREAMBLE
definecolor{DIFgreen}{rgb}{0,0.5,0} %DIF PREAMBLE
\DeclareOldFontCommand{\sf}{\normalfont\sffamily}{\mathsf} %DIF PREAMBLE
\providecommand{\DIFaddtex}[1]{{\protect\color{DIFgreen} \sf #1}} %DIF PREAMBLE
\providecommand{\DIFdeltex}[1]{{\protect\color{red} \scriptsize #1}} %DIF PREAMBLE
%DIF SAFE PREAMBLE %DIF PREAMBLE
\providecommand{\DIFaddbegin}{} %DIF PREAMBLE
\providecommand{\DIFaddend}{} %DIF PREAMBLE
\providecommand{\DIFdelbegin}{} %DIF PREAMBLE
\providecommand{\DIFdelend}{} %DIF PREAMBLE
\providecommand{\DIFmodbegin}{} %DIF PREAMBLE
\providecommand{\DIFmodend}{} %DIF PREAMBLE
%DIF FLOATSAFE PREAMBLE %DIF PREAMBLE
\providecommand{\DIFaddFL}[1]{\DIFadd{#1}} %DIF PREAMBLE
\providecommand{\DIFdelFL}[1]{\DIFdel{#1}} %DIF PREAMBLE
\providecommand{\DIFaddbeginFL}{} %DIF PREAMBLE
\providecommand{\DIFaddendFL}{} %DIF PREAMBLE
\providecommand{\DIFdelbeginFL}{} %DIF PREAMBLE
\providecommand{\DIFdelendFL}{} %DIF PREAMBLE
%DIF HYPERREF PREAMBLE %DIF PREAMBLE
\providecommand{\DIFadd}[1]{\texorpdfstring{\DIFaddtex{#1}}{#1}} %DIF PREAMBLE
\providecommand{\DIFdel}[1]{\texorpdfstring{\DIFdeltex{#1}}{}} %DIF PREAMBLE
\newcommand{\DIFscaledelfig}{0.5}
%DIF HIGHLIGHTGRAPHICS PREAMBLE %DIF PREAMBLE
\RequirePackage{settobox} %DIF PREAMBLE
\RequirePackage{letltxmacro} %DIF PREAMBLE
\newsavebox{\DIFdelgraphicsbox} %DIF PREAMBLE
\newlength{\DIFdelgraphicswidth} %DIF PREAMBLE
\newlength{\DIFdelgraphicsheight} %DIF PREAMBLE
% store original definition of \includegraphics %DIF PREAMBLE
\LetLtxMacro{\DIFOincludegraphics}{\includegraphics} %DIF PREAMBLE
\newcommand{\DIFaddincludegraphics}[2][]{{\color{DIFgreen}\fbox{\DIFOincludegraphics[#1]{#2}}}} %DIF PREAMBLE
\newcommand{\DIFdelincludegraphics}[2][]{% %DIF PREAMBLE
\sbox{\DIFdelgraphicsbox}{\DIFOincludegraphics[#1]{#2}}% %DIF PREAMBLE
\settoboxwidth{\DIFdelgraphicswidth}{\DIFdelgraphicsbox} %DIF PREAMBLE
\settoboxtotalheight{\DIFdelgraphicsheight}{\DIFdelgraphicsbox} %DIF PREAMBLE
\scalebox{\DIFscaledelfig}{% %DIF PREAMBLE
\parbox[b]{\DIFdelgraphicswidth}{\usebox{\DIFdelgraphicsbox}\\[-\baselineskip] \rule{\DIFdelgraphicswidth}{0em}}\llap{\resizebox{\DIFdelgraphicswidth}{\DIFdelgraphicsheight}{% %DIF PREAMBLE
\setlength{\unitlength}{\DIFdelgraphicswidth}% %DIF PREAMBLE
\begin{picture}(1,1)% %DIF PREAMBLE
\thicklines\linethickness{2pt} %DIF PREAMBLE
{\color[rgb]{1,0,0}\put(0,0){\framebox(1,1){}}}% %DIF PREAMBLE
{\color[rgb]{1,0,0}\put(0,0){\line( 1,1){1}}}% %DIF PREAMBLE
{\color[rgb]{1,0,0}\put(0,1){\line(1,-1){1}}}% %DIF PREAMBLE
\end{picture}% %DIF PREAMBLE
}\hspace*{3pt}}} %DIF PREAMBLE
} %DIF PREAMBLE
\LetLtxMacro{\DIFOaddbegin}{\DIFaddbegin} %DIF PREAMBLE
\LetLtxMacro{\DIFOaddend}{\DIFaddend} %DIF PREAMBLE
\LetLtxMacro{\DIFOdelbegin}{\DIFdelbegin} %DIF PREAMBLE
\LetLtxMacro{\DIFOdelend}{\DIFdelend} %DIF PREAMBLE
\DeclareRobustCommand{\DIFaddbegin}{\DIFOaddbegin \let\includegraphics\DIFaddincludegraphics} %DIF PREAMBLE
\DeclareRobustCommand{\DIFaddend}{\DIFOaddend \let\includegraphics\DIFOincludegraphics} %DIF PREAMBLE
\DeclareRobustCommand{\DIFdelbegin}{\DIFOdelbegin \let\includegraphics\DIFdelincludegraphics} %DIF PREAMBLE
\DeclareRobustCommand{\DIFdelend}{\DIFOaddend \let\includegraphics\DIFOincludegraphics} %DIF PREAMBLE
\LetLtxMacro{\DIFOaddbeginFL}{\DIFaddbeginFL} %DIF PREAMBLE
\LetLtxMacro{\DIFOaddendFL}{\DIFaddendFL} %DIF PREAMBLE
\LetLtxMacro{\DIFOdelbeginFL}{\DIFdelbeginFL} %DIF PREAMBLE
\LetLtxMacro{\DIFOdelendFL}{\DIFdelendFL} %DIF PREAMBLE
\DeclareRobustCommand{\DIFaddbeginFL}{\DIFOaddbeginFL \let\includegraphics\DIFaddincludegraphics} %DIF PREAMBLE
\DeclareRobustCommand{\DIFaddendFL}{\DIFOaddendFL \let\includegraphics\DIFOincludegraphics} %DIF PREAMBLE
\DeclareRobustCommand{\DIFdelbeginFL}{\DIFOdelbeginFL \let\includegraphics\DIFdelincludegraphics} %DIF PREAMBLE
\DeclareRobustCommand{\DIFdelendFL}{\DIFOaddendFL \let\includegraphics\DIFOincludegraphics} %DIF PREAMBLE
%DIF END PREAMBLE EXTENSION ADDED BY LATEXDIFF

\begin{document}

\maketitle

\thispagestyle{empty}
\vspace*{\fill}
\begin{center}
  \ccby\\[0.5em]
  This work is licensed under a
  \href{https://creativecommons.org/licenses/by/4.0/}{Creative Commons Attribution 4.0 International License}.\\[0.3em]
  You are free to share and adapt this material for any purpose,\\
  provided appropriate credit is given, a link to the license is provided,\\
  and any changes made are indicated.
\end{center}
\vspace*{\fill}

\tableofcontents

\begin{abstract}
This paper presents a mathematical framework for the Quantum Observable
Universe (QOU).  We introduce the notion of quantum torsion in this setting,
study its lifting to higher categorical structures, and examine the resulting
interpretational and experimental consequences.  Speculative extensions are
discussed in the final chapter.
\end{abstract}

\chapter{Quantum Observable Universe}
\label{chap:qou}



\begin{definition}[Quantum Universe]
\label{def:quantum-universe}
Let $k$ be a field.  A \emph{quantum universe over~$k$} is a quadruple
$(\mathcal{C}, A, \omega, q)$ where:
\begin{enumerate}
  \item $(\mathcal{C}, \otimes, \mathbf{1})$ is a $k$-linear rigid
        monoidal category;
  \item $A$ is a \emph{Frobenius--Hopf object} in $\mathcal{C}$: an object
        equipped with a Hopf structure
        $(\mu, \eta, \Delta, \varepsilon, S)$ and a Frobenius form
        $\lambda \colon A \to \mathbf{1}$ such that the pairing
        $A \otimes A \xrightarrow{\mu} A \xrightarrow{\lambda} \mathbf{1}$
        is non-degenerate and the Frobenius relation
        $(\mu \otimes \mathrm{id}) \circ
        (\mathrm{id} \otimes \Delta_\mu)
        = \Delta_\mu \circ \mu$ holds, where $\Delta_\mu$ is the
        comultiplication dual to $\mu$ under the pairing;
  \item $\omega \colon \mathcal{C} \to \mathrm{Vect}_k$ is a fibre functor
        (exact faithful $k$-linear monoidal functor);
  \item $q \in k^\times$ is a deformation parameter (invertibility of~$q$
        is required because the Hopf-algebra relations defining~$A$
        involve~$q^{-1}$; at $q = 0$ the antipode~$S$ degenerates and
        the Hopf structure collapses); the Hopf structure on
        $A$ depends on~$q$, and $q = 1$ recovers the classical
        (cocommutative) case.
\end{enumerate}
The \emph{quantum gauge group} is the affine $k$-group scheme
$G_q = \mathrm{Aut}^{\otimes}(\omega)$, reconstructed via
Tannaka--Krein duality.
\end{definition}

\begin{definition}[Quantum Observable Universe]
\label{def:quantum-observable-universe}
A quantum universe $(\mathcal{C}, A, \omega, q)$ over~$k$
(Definition~\ref{def:quantum-universe}) is \emph{observable} if:
\begin{enumerate}
  \item the Frobenius form $\lambda$ on $A$ is non-degenerate; and
  \item $\mathrm{Hom}_{\mathcal{C}}(\mathbf{1},\, \mathrm{Ad}(A)) = 0$,
        where $\mathrm{Ad}(A)$ is the adjoint object
        ($A$ with the conjugation coaction
        $a \mapsto \textstyle\sum a_{(1)}\, b\, S(a_{(2)})$).
\end{enumerate}
\end{definition}

\begin{remark}[Base-field independence]
\label{rem:base-field-independence}
Definitions~\ref{def:quantum-universe} and~\ref{def:quantum-observable-universe}
are formulated over an arbitrary field~$k$.  The Frobenius--Hopf axioms
replace the analytic vanishing conditions of the $\mathbb{C}$ realisation
(Example~\ref{def:c-quantum-observable-universe}) with a single algebraic
non-degeneracy condition on the Frobenius form.  Concrete realisations
include:
\begin{itemize}
  \item $k = \mathbb{C}$: the classical setting of CR geometry and contact
        manifolds (Example~\ref{def:c-quantum-universe});
  \item $k = \mathbb{Q}_\pp$ (the $\pp$-adic field): the gauge datum becomes
        an $F$-isocrystal with Frobenius lift, and integrability is controlled
        by the $\pp$-adic Frobenius endomorphism
        (Example~\ref{ex:p-adic-realisation},
        Chapter~\ref{chap:p-adic});
  \item $k = \mathbb{F}_\pp$ (a finite field): the monoidal category is
        the category of finite-dimensional representations of a quantum group
        over~$\mathbb{F}_\pp$
        (Example~\ref{ex:finite-field-realisation}).
\end{itemize}
The fibre functor $\omega$ ensures that $G_q = \mathrm{Aut}^{\otimes}(\omega)$
is well-defined over any~$k$.
\end{remark}



\begin{example}[$\mathbb{C}$ realisation: quantum universe]
\label{def:c-quantum-universe}
Over $k = \mathbb{C}$, a quantum universe
(Definition~\ref{def:quantum-universe}) is realised by a $5$-tuple
$(M, \theta, G, P, E)$ where:
\begin{enumerate}
  \item $M$ is a CR manifold of hypersurface type
        (real dimension $2n+1$, $n \geq 1$) with non-degenerate Levi form
        of signature $(r,s)$, $r + s = n$;
  \item $\theta$ is a pseudo-Hermitian structure (contact form) on $M$,
        determining the Levi distribution $H = \ker\theta$, the Reeb vector
        field $\xi$, the Webster metric $g_\theta$ (of signature $(2r,2s)$
        on $H$), and the Levi form $L_\theta$;
  \item $G = \mathrm{SU}(2)$;
  \item $P \to M$ is a principal $G$-bundle whose holonomy reduces to
        $\mathrm{SO}(3) \subset \mathrm{SU}(2)/\{\pm I\}$;
  \item $E = P \times_G \mathbb{CP}^1 \to M$ is the associated
        \emph{state bundle}, where $\mathrm{SU}(2)$ acts on
        $\mathbb{CP}^1$ by M\"obius transformations.
\end{enumerate}
The monoidal category is $\mathcal{C} = \mathrm{Rep}(\mathrm{SU}_{q(p)}(2))$
(the representation category of the quantum group at the substrate
parameter~$q(p)$, Definition~\ref{def:substrate-parameter}), the
Frobenius--Hopf object is
$A = \mathcal{O}_q(\mathrm{SU}(2))$ with Frobenius form given by the
Haar state, and the fibre functor $\omega$ is the forgetful functor
to $\mathrm{Vect}_{\mathbb{C}}$.
\end{example}

\begin{remark}[Notation and special cases]
\label{rem:notation-special-cases}
When $\theta$ is understood from context, we abbreviate the $5$-tuple to
$(M, G, P, E)$.  The \emph{strictly pseudoconvex} case corresponds to
signature $(n,0)$, where $L_\theta$ is positive definite.
\end{remark}

\begin{example}[$\mathbb{C}$ realisation: quantum observable universe]
\label{def:c-quantum-observable-universe}
A \emph{$\mathbb{C}$ quantum observable universe} is a $\mathbb{C}$ quantum
universe $(M, \theta, G, P, E)$
(Example~\ref{def:c-quantum-universe}) that satisfies
the vanishing conditions
\[
  H^0(M, \mathrm{Ad}(P)) = 0
  \qquad\text{and}\qquad
  \Gamma_c(M, E) = 0.
\]
These analytic conditions are the $\mathbb{C}$ realisation of the
categorical observability conditions of
Definition~\ref{def:quantum-observable-universe}: the Frobenius
non-degeneracy of the Haar state on $\mathcal{O}_q(\mathrm{SU}(2))$,
combined with irreducibility of the instanton connection
(Lemma~\ref{lem:irreducible-connection}), forces
$H^0(M, \mathrm{Ad}(P)) = 0$; the section-vanishing
$\Gamma_c(M, E) = 0$ follows from the non-existence of
$A$-equivariant sections of the associated $\mathbb{CP}^1$-bundle.
\end{example}





\begin{definition}[Instanton connection]
\label{def:instanton}
Let $(M, \theta, G, P, E)$ be a $\mathbb{C}$ quantum universe with
$\dim_{\mathbb{R}} H = 4$ ($n = 2$).  An \emph{instanton connection} on $P$
is a connection $A$ whose curvature $F_A$, restricted to the horizontal
distribution $H = \ker\theta$, is self-dual:
$F_A|_H = {*_H}\,F_A|_H$, where $*_H$ is the Hodge star on
$\Lambda^2 H^*$ determined by the Webster metric $g_\theta$.  More generally,
when $\dim_{\mathbb{R}} H \neq 4$, an instanton connection is a
Yang--Mills connection that minimises the functional
$\int_M \lVert F_A|_H \rVert^2\,\theta \wedge (d\theta)^n$.
\end{definition}

\begin{lemma}[Irreducibility of the connection]
\label{lem:irreducible-connection}
Let $(M, G, P, E)$ be a $\mathbb{C}$ quantum observable universe.  Then the
vanishing condition $H^0(M, \mathrm{Ad}(P)) = 0$ implies that every
instanton connection $A$ on $P$ (Definition~\ref{def:instanton}) is
irreducible.
\end{lemma}

\begin{proof}
Suppose for contradiction that an instanton connection $A$ on $P$ is
reducible.  Then the stabiliser of $A$ under the gauge group contains a
non-trivial element $s \in \Gamma(M,\, \mathrm{Ad}(P))$ satisfying
$D_A s = 0$.  Such $s$ is a global parallel section of $\mathrm{Ad}(P)$,
and hence a non-zero element of $H^0(M, \mathrm{Ad}(P))$, contradicting
$H^0(M, \mathrm{Ad}(P)) = 0$.  Hence $A$ is irreducible.
\end{proof}

\begin{proposition}[Self-duality of the curvature]
\label{prop:self-dual-curvature}
Let $(M, \theta, G, P, E)$ be a $\mathbb{C}$ quantum observable universe
with $\dim_{\mathbb{R}} M = 5$ (i.e.\ the CR dimension $n = 2$, so that the
Levi distribution $H = \ker\theta$ has $\dim_{\mathbb{R}} H = 4$).  Let $A$
be an instanton connection on $P$.  Then the curvature $F_A$, restricted to
$H$, is self-dual with respect to the Hodge star $*_H$ determined by the
Webster metric $g_\theta$ and the complex structure $J$:
\[
  F_A\big|_H \;=\; {*_H}\, F_A\big|_H.
\]
\end{proposition}

\begin{proof}
The contact form $\theta$ (Example~\ref{def:c-quantum-universe}) splits
the tangent bundle as $TM = H \oplus \mathbb{R}\xi$, where
$H = \ker\theta$ has real dimension $4$.  The curvature $F_A$ is an
$\mathrm{Ad}(P)$-valued $2$-form on $M$.  By the Reeb decomposition
(Definition~\ref{def:reeb-decomposition}), $F_A$ splits into horizontal and
vertical parts; we consider $F_A|_H \in \Gamma(\Lambda^2 H^* \otimes
\mathrm{Ad}(P))$.

Since $\dim_{\mathbb{R}} H = 4$, the operator $*_H$ maps $2$-forms to
$2$-forms and satisfies $*_H^2 = \mathrm{id}$ on $\Lambda^2 H^*$.  This
decomposes $\Lambda^2 H^* = \Lambda^+\! H^* \oplus \Lambda^-\! H^*$ into
self-dual and anti-self-dual parts.

An instanton connection $A$ minimises the Yang--Mills functional.  By the
Bogomolny argument, the integrated norm satisfies
\[
  \int_M \lVert F_A|_H \rVert^2 \,\theta \wedge (d\theta)^2
  \;=\; 2\!\int_M \lVert F_A^-|_H \rVert^2 \,\theta \wedge (d\theta)^2
        + \int_M \mathrm{tr}(F_A|_H \wedge F_A|_H)
          \wedge \theta,
\]
where the topological term
$\int_M \mathrm{tr}(F_A|_H \wedge F_A|_H) \wedge \theta$ is fixed by the
characteristic class of $P$.  Minimising the functional forces
$F_A^-|_H = 0$, which is precisely the self-duality condition
$F_A|_H = {*_H}\,F_A|_H$.
\end{proof}







\begin{definition}[Substrate parameter]
\label{def:substrate-parameter}
Let $(\mathcal{C}, A, \omega, q)$ be a quantum universe over a field~$k$
(Definition~\ref{def:quantum-universe}).  Let
$\mathcal{U}^{\mathbb{Z}}_q(\mathfrak{su}(2))$ denote the
\emph{Lusztig integral form}~\cite{lusztig1993}: the
$\mathbb{Z}[q,q^{-1}]$-subalgebra of $\mathcal{U}_q(\mathfrak{su}(2))$
generated by $K^{\pm 1}$ and the divided powers
$E^{(n)} = E^n/[n]_q!$, $F^{(n)} = F^n/[n]_q!$.
The \emph{substrate parameter} is a ring homomorphism
\[
  \iota \;\colon\; \mathbb{Z}[q, q^{-1}] \;\longrightarrow\; k,
  \qquad q \;\longmapsto\; q_0 \in k^{\times},
\]
and the base-changed algebra
$\mathcal{U}^{\mathbb{Z}}_q(\mathfrak{su}(2))
  \otimes_{\mathbb{Z}[q,q^{-1}]} k$
is the quantised enveloping algebra at the specialised value~$q_0$.
When the quantum universe is fibred over a space~$M$ (as in the
$\mathbb{C}$ realisation, Example~\ref{def:c-quantum-universe}), $q_0$
may vary from point to point, yielding a \emph{substrate field}
$q \colon M \to k^{\times}$.
\end{definition}

\begin{remark}[$\mathbb{C}$ substrate via the Levi form]
\label{rem:c-substrate-via-levi}
In the $\mathbb{C}$ quantum observable universe
(Example~\ref{def:c-quantum-observable-universe}), let $A$ be an
instanton connection on~$P$
(Definition~\ref{def:instanton}), $F_A^+$ its self-dual curvature
(Proposition~\ref{prop:self-dual-curvature}), and
$L_\theta \colon H \times H \to \mathbb{C}$ the Levi form on
$H = \ker\theta$, defined by
$L_\theta(V, W) = -i\,d\theta(V, \bar{W})$.
The specialisation map of
Definition~\ref{def:substrate-parameter} sends
\[
  q \;\longmapsto\;
  q(p) \;=\; 1 + \lvert\langle L_\theta,\, F_A^+ \rangle\rvert(p),
  \qquad p \in M,
\]
where $\langle L_\theta, F_A^+\rangle$ is the natural contraction
(Remark~\ref{rem:substrate-contraction}) and $\lvert\cdot\rvert$ denotes
the absolute value.  This defines a smooth field
$q \colon M \to [1,\infty) \subset \mathbb{C}^{\times}$ with $q(p) = 1$
at points where $F_A^+ = 0$ (classicality) and $q(p) \to \infty$ in the
deep-quantum regime.

\textbf{Relation to the exponential.}
Over~$\mathbb{Q}$ (or any subfield of~$\mathbb{C}$), the classical
exponential $\exp(\lvert\langle L_\theta, F_A^+\rangle\rvert)$ is an
equally valid specialisation map, and the two are related by the Taylor
expansion $\exp(x) = 1 + x + \tfrac{1}{2}x^2 + \cdots$:
the formula $q(p) = 1 + \lvert\langle L_\theta, F_A^+\rangle\rvert(p)$
is the first-order approximation of
$\exp\!\bigl(\lvert\langle L_\theta, F_A^+\rangle\rvert(p)\bigr)$.
Both define smooth fields $M \to [1,\infty)$ with
the same boundary behaviour ($q = 1$ at classicality, $q \to \infty$
deep-quantum), and all downstream formulas
($\hbar_q = 1 - q^{-1}$, the Heisenberg commutator, the Schrödinger
equation, the Yang--Mills functional) are algebraic in~$q$ and therefore
insensitive to the choice between $1+x$ and $\exp(x)$.  The algebraic
form is adopted in Definition~\ref{def:substrate-parameter} because it
is meaningful over every base field~$k$, whereas $\exp(x)$ requires an
archimedean topology.

The non-archimedean analogues are developed in
Chapter~\ref{chap:p-adic}: over~$\mathbb{Q}_\pp$ the Levi-form
contraction is replaced by the Frobenius endomorphism and
$q = \pp^{\,-v_\pp(\det\varphi)}$
(Example~\ref{ex:p-adic-realisation}); over~$\mathbb{F}_\pp$, $q$
is a primitive root of unity
(Example~\ref{ex:finite-field-realisation}).
\end{remark}

\begin{remark}[Natural contraction and field character of $q$]
\label{rem:substrate-contraction}
The Levi form $L_\theta$ is a section of $\Lambda^{1,1} H^*$
(the $(1,1)$-component of $\Lambda^2 H^*$ with respect to the complex
structure $J$).  The self-dual curvature $F_A^+$ is a section of
$\Lambda^+ H^* \otimes \mathrm{Ad}(P)$.  On a $4$-dimensional oriented
inner-product space, $\Lambda^{1,1} \subset \Lambda^+$, so $L_\theta$ and
$F_A^+$ live in the same subbundle.  The Webster metric $g_\theta$
provides a fibrewise inner product on $\Lambda^+ H^*$; composing this
with the Killing form on $\mathrm{Ad}(P)$ yields the natural contraction
\[
  \langle L_\theta,\, F_A^+ \rangle_p
  \;\in\; \mathbb{R}, \qquad p \in M.
\]
The substrate parameter is not a
fixed constant but a field on $M$: its value varies from point to point
according to the local coupling between the Levi form and the self-dual
curvature.
\end{remark}

\begin{definition}[Quantum Gauge Group]
\label{def:quantum-gauge-group}
Let $(M, G, P, E)$ be a $\mathbb{C}$ quantum observable universe with substrate
parameter $q$ (Definition~\ref{def:substrate-parameter}).  The
\emph{quantum gauge group} is the field of quantum groups
$\{\mathrm{SU}_{q(p)}(2)\}_{p \in M}$ obtained by deforming the structure
group $G = \mathrm{SU}(2)$ pointwise at the parameter $q(p)$.
\end{definition}

\begin{definition}[Quantum Planck Constant]
\label{def:quantum-planck-constant}
Let $(M, G, P, E)$ be a $\mathbb{C}$ quantum observable universe with substrate
parameter $q$ (Definition~\ref{def:substrate-parameter}).  The
\emph{quantum Planck constant} is the smooth field
$\hbar_q \colon M \to [0,1)$ defined by
\[
  \hbar_q(p) \;=\; 1 - q(p)^{-1}, \qquad p \in M.
\]
\end{definition}

\begin{remark}[Limiting regimes of $\hbar_q$]
\label{rem:planck-limits}
Like $q$ itself, $\hbar_q$ is not a universal constant but a field on $M$.
As $q(p) \to 1$ (classicality), $\hbar_q(p) \to 0$ and quantum effects
vanish; as $q(p) \to \infty$ (deep-quantum regime), $\hbar_q(p) \to 1$.
\end{remark}

\begin{remark}[Heisenberg uncertainty as commutator]
\label{rem:heisenberg-commutator}
In the $\mathbb{C}$ quantum observable universe the Heisenberg uncertainty
principle is expressed through the $q$-deformed commutator.  Let $\hat{X}$ and
$\hat{P}$ denote position and momentum observables on the fibre
$\mathbb{CP}^1$.  Let $J \colon H \to H$ be the complex structure operator
on the Levi distribution ($J^2 = -\mathrm{Id}$); throughout this paper every
occurrence of the imaginary unit~$i$ in a physical equation is identified with
the geometric action of~$J$ on horizontal planes.  We define their
$q$-dependent commutator at $p \in M$ by
\[
  [\hat{X},\, \hat{P}]_q \;:=\; J\,\hbar_q(p),
\]
where $\hbar_q(p) = 1 - q(p)^{-1}$
(Definition~\ref{def:quantum-planck-constant}).  In this remark the subscript
$q$ serves only to record the dependence of the commutator on the substrate
field $q$ and should not be confused with any $q$-deformed Lie bracket
$[\cdot,\cdot]_q$ on vector fields that may appear elsewhere.  The commutator
encodes the local uncertainty: when $q(p) \to 1$ (classicality),
$[\hat{X}, \hat{P}]_q \to 0$ and observables commute; when $q(p) \to \infty$
(deep-quantum regime), $[\hat{X}, \hat{P}]_q \to J$ and maximal uncertainty
prevails.

\medskip
\noindent\textbf{Analogy with the Lie bracket.}\;
The position--momentum commutator $[\hat{X},\hat{P}]_q$ is the
quantum-mechanical counterpart of the Lie bracket
$[e_a,e_b]=\epsilon_{abc}\,e_c$ in $\mathfrak{su}(2)$
(Definition~\ref{def:su2-basis}).  Both commutators measure the
first-order non-commutativity of their respective algebras
(Remark~\ref{rem:first-order-taylor}).  In the braid-group setting,
this non-commutativity appears as the defining relation
$\sigma_1\sigma_2\sigma_1=\sigma_2\sigma_1\sigma_2$
(Definition~\ref{def:braid-group}); on the Bring surface, it is
realised by the $S_3$ permutation action on the spectral resolvents
(Remark~\ref{rem:braid-group-action}).

\medskip
\noindent\textbf{Commutator as Reeb vector.}\;
The commutator $[\hat{X},\hat{P}]_q$ may be identified with the value
of the $q$-deformed Reeb vector field~$\xi_q$
(Definition~\ref{def:q-reeb-vector-field}) evaluated on the contact
form.  Concretely, the Reeb field is the unique direction along which
the symplectic pairing $d\theta$ degenerates, so $\theta(\xi_q)$
encodes precisely the obstruction to commutativity of canonical
coordinates.  On the exceptional divisor $\mathbb{CP}^1$ produced by
blowing up a spacetime defect on the Bring surface
(Remark~\ref{rem:brings-smoothness}), the Reeb vector field generates
the $U(1)$ fibration of the Hopf bundle $S^3\to\mathbb{CP}^1$, and
its ``strength'' is set by~$\hbar_q(p)$.
\end{remark}



\begin{definition}[$q$-Harmonic form (Particle)]
\label{def:q-harmonic-form}
Let $(M, \theta, G, P, E)$ be a $\mathbb{C}$ quantum observable universe with
substrate parameter $q$ (Definition~\ref{def:substrate-parameter}).  Let $d_q$
denote the $q$-deformed exterior derivative
(Definition~\ref{def:q-ext-deriv} below) and let $\delta_q$ denote its formal
adjoint with respect to the inner product on
$\Gamma(M,\, \Lambda^k T^*M \otimes J^1 E)$ induced by the Webster metric
$g_\theta$.  A \emph{particle} is a section
$\alpha \in \Gamma(M,\, \Lambda^k T^*M \otimes J^1 E)$ that is
$q$-harmonic, i.e.\ lies in the kernel of the $q$-Laplacian:
\[
  \Delta_q \alpha
  \;=\; (d_q\,\delta_q + \delta_q\,d_q)\,\alpha
  \;=\; 0.
\]
\end{definition}

\begin{definition}[Reeb decomposition]
\label{def:reeb-decomposition}
Let $(M, \theta)$ be a CR manifold with non-degenerate Levi form
(Example~\ref{def:c-quantum-universe}) and let $\xi$ denote the
associated \emph{Reeb vector field}, characterised by
$\iota_\xi\,d\theta = 0$ and $\theta(\xi) = 1$.  Every differential
$k$-form $\alpha$ ($k \geq 1$) on $M$ decomposes uniquely as
\[
  \alpha \;=\; \alpha_H + \theta \wedge \iota_\xi\,\alpha,
\]
where $\alpha_H := \alpha - \theta \wedge \iota_\xi\,\alpha$ is the
\emph{horizontal component} (satisfying $\iota_\xi\,\alpha_H = 0$) and
$\alpha_V := \theta \wedge \iota_\xi\,\alpha$ is the \emph{vertical
component}.  A form is \emph{purely horizontal} if and only if
$\iota_\xi\,\alpha = 0$.
\end{definition}

\begin{definition}[Fermion]
\label{def:fermion}
A \emph{fermion} is a $q$-harmonic form
(Definition~\ref{def:q-harmonic-form}) that has a non-trivial vertical
component in the Reeb decomposition
(Definition~\ref{def:reeb-decomposition}):
\[
  \Delta_q\,\alpha = 0
  \qquad\text{and}\qquad
  \iota_\xi\,\alpha_V \neq 0.
\]
\end{definition}

\begin{definition}[Boson]
\label{def:boson}
A \emph{boson} is a $q$-harmonic form
(Definition~\ref{def:q-harmonic-form}) that is purely horizontal in the
Reeb decomposition (Definition~\ref{def:reeb-decomposition}):
\[
  \Delta_q\,\alpha = 0
  \qquad\text{and}\qquad
  \alpha = \alpha_H
  \quad(\text{equivalently, } \iota_\xi\,\alpha = 0
        \text{ by Definition~\ref{def:reeb-decomposition}}).
\]
\end{definition}



\begin{definition}[$q$-Deformed exterior derivative]
\label{def:q-ext-deriv}
Let $(M, \theta, G, P, E)$ be a $\mathbb{C}$ quantum observable universe
with substrate parameter $q$ (Definition~\ref{def:substrate-parameter}).
The \emph{$q$-deformed exterior derivative}
$d_q \colon \Omega^k(M) \to \Omega^{k+1}(M)$ is defined on $k$-forms by
\[
  d_q\,\alpha
  \;=\; q^{-1}\,d\alpha
        + (1 - q^{-1})\,\theta \wedge \mathcal{L}_\xi\,\alpha,
\]
where $d$ is the ordinary exterior derivative, $\theta$ the contact form,
and $\mathcal{L}_\xi$ the Lie derivative along the Reeb field.
\end{definition}

\begin{definition}[$q$-Deformed Lie bracket]
\label{def:q-lie-bracket}
Let $V, W \in \mathfrak{X}(M)$ be smooth vector fields.  The
\emph{$q$-deformed Lie bracket} is
\[
  [V, W]_q
  \;=\; q^{-1}\,[V, W]
        + (1 - q^{-1})\,\bigl(\theta(V)\,\mathcal{L}_\xi W
                              - \theta(W)\,\mathcal{L}_\xi V\bigr),
\]
where $[\,\cdot\,,\cdot\,]$ is the standard Lie bracket on $\mathfrak{X}(M)$.
\end{definition}

\begin{definition}[$q$-Deformed connection]
\label{def:q-connection}
Let $A$ be an instanton connection on $P$
(Lemma~\ref{lem:irreducible-connection}).  The \emph{$q$-deformed
connection} $A_q$ is defined by
\[
  A_q
  \;=\; q^{-1}\,A
        + (1 - q^{-1})\,\iota_\xi\,A\;\theta,
\]
\end{definition}

\begin{definition}[$q$-Deformed Reeb vector field]
\label{def:q-reeb-vector-field}
Let $(M, \theta, G, P, E)$ be a $\mathbb{C}$ quantum observable universe
with substrate parameter $q$ (Definition~\ref{def:substrate-parameter})
and Reeb vector field $\xi$ (Definition~\ref{def:reeb-decomposition}).  The
\emph{$q$-deformed Reeb vector field} is the vector field $\xi_q$ on $M$
defined by its action on smooth functions $f \in C^\infty(M)$:
\[
  \xi_q(f) \;=\; q^{-1}\,\xi(f) + (1 - q^{-1})\,[\xi, X_f]_q,
\]
where $X_f$ is the horizontal lift of $\mathrm{grad}_{g_\theta}\,f$ to
$H = \ker\theta$ (using the Webster metric $g_\theta$) and
$[\,\cdot\,,\cdot\,]_q$ is the $q$-deformed Lie bracket
(Definition~\ref{def:q-lie-bracket}).
\end{definition}

\begin{definition}[$q$-Deformed Reeb flow]
\label{def:q-reeb-flow}
Let $\xi_q$ be the $q$-deformed Reeb vector field
(Definition~\ref{def:q-reeb-vector-field}).  The \emph{$q$-deformed Reeb flow}
is the one-parameter family of diffeomorphisms
$\phi_t^q \colon M \to M$, $t \in \mathbb{R}$, generated by $\xi_q$:
\[
  \frac{d}{dt}\,\phi_t^q \;=\; \xi_q \circ \phi_t^q,
  \qquad \phi_0^q = \mathrm{id}_M.
\]
\end{definition}

\begin{definition}[$q$-Torsion leaf]
\label{def:q-torsion-leaf}
Let $H = \ker\theta \subset TM$ be the horizontal distribution on the
CR manifold $M$.  A \emph{$q$-torsion leaf} is a maximal integral
submanifold $L \hookrightarrow M$ of $H$ that is invariant under the
$q$-deformed Reeb flow (Definition~\ref{def:q-reeb-flow}):
\[
  T_p L \;\subset\; H_p
  \quad\text{for all } p \in L,
  \qquad\text{and}\qquad
  \phi_t^q(L) = L
  \quad\text{for all } t \in \mathbb{R}.
\]
\end{definition}

\begin{remark}[$q$-Torsion leaves: existence and Frobenius obstruction]
\label{rem:q-torsion-existence}
The horizontal distribution $H = \ker\theta$ on a contact manifold is
maximally non-integrable: $d\theta|_{H \times H}$ is non-degenerate.
By the Frobenius theorem, maximal integral submanifolds of $H$ have
dimension at most~$n$ (half the rank of~$H$); these are the
\emph{Legendrian submanifolds}.  Definition~\ref{def:q-torsion-leaf}
therefore requires $L$ to be a Legendrian submanifold invariant under the
$q$-deformed Reeb flow.

Existence is not automatic: it depends on the dynamics of~$\xi_q$.
When $q \to 1$ and $\xi_q \to \xi$, invariant Legendrian submanifolds
exist whenever the Reeb flow is periodic or integrable (e.g.\ on
$S^3/\Gamma$).  For $q > 1$, the deformation generically breaks
periodicity, and one must appeal to KAM-type persistence results.

In the category of Lie algebroids over~$M$, the Frobenius integrability
condition is equivalent to $H$ carrying a Frobenius structure
(Definition~\ref{def:quantum-universe}).  Since the Levi form
$d\theta|_{H \times H}$ is non-degenerate, this Frobenius structure
\emph{fails} --- and it is precisely this failure that forces the
dimension bound on integral submanifolds and gives $q$-torsion leaves
their rigidity.
\end{remark}

\begin{remark}[Physical interpretation of the Reeb decomposition]
\label{rem:reeb-physical}
The Reeb vector field $\xi$ singles out the direction of internal spin on the
CR manifold.  Forms that ``see'' this direction (fermions) carry half-integer
spin; forms confined to the horizontal distribution (bosons) do not couple to
the Reeb flow and carry integer spin.  The decomposition
$\alpha = \alpha_H + \alpha_V$ is therefore the geometric origin of the
spin--statistics distinction in this framework.  Experimental confirmation of
this picture, including photonic quantum Hall measurements and many-body
interference in crystal lattices, is discussed in
Chapter~\ref{chap:experimental}.
\end{remark}



\begin{remark}[General Relativity as the ``Substrate Geometry'']
\label{rem:gr-substrate-geometry}
In this model the metric is effectively replaced by the Levi form
$L_\theta$ and its coupling to the gauge curvature $F_A$.  Gravity is
expressed through the Variational Bicomplex as the back-reaction of the
gauge fields on the CR manifold $M$.

Define the \emph{$q$-weighted Yang--Mills--Levi functional}
\[
  S_q[A, \theta]
  \;=\; \int_M q(p)\,\lVert F_A^+|_H \rVert^2\,
        \theta \wedge (d\theta)^n,
\]
where $q(p) = 1 + \lvert\langle L_\theta, F_A^+\rangle\rvert(p)$
(Remark~\ref{rem:c-substrate-via-levi}).
Variation with respect to the Levi form $L_\theta$ yields field equations
of the schematic form
\[
  \frac{\delta S_q}{\delta L_\theta}
  \;=\; 0
  \qquad\Longrightarrow\qquad
  \mathrm{Ric}_{g_\theta} \;=\; T^{(q)},
\]
where $\mathrm{Ric}_{g_\theta}$ is the Webster Ricci curvature and
$T^{(q)}$ is a $q$-weighted stress-energy $2$-tensor on $H$ (both
coordinate-free objects).  The precise derivation of the field equations
will be given in Chapter~\ref{chap:experimental}; here we note only that
the geometry of $M$ is forced to curve to accommodate the tension of the
$\mathrm{SU}(2)$ instantons.
\end{remark}

\begin{proposition}[Main Commutative Diagram]
\label{prop:main-commutative-diagram}
Let $(M, \theta, G, P, E)$ be a $\mathbb{C}$ quantum observable universe
with instanton connection $A$ (Definition~\ref{def:instanton}) and substrate
parameter $q$ (Definition~\ref{def:substrate-parameter}).  Let
$\mathbf{E}$ denote the Euler--Lagrange operator of the
$q$-weighted Yang--Mills--Levi functional $S_q$
(Remark~\ref{rem:gr-substrate-geometry}), and let $\nabla^A$ denote
parallel transport along $A$.  The following diagram commutes:
\[
\begin{CD}
J^1E @>{\mathbf{E}[S_q]}>>
  \bigl\{ L_\theta : \mathrm{Ric}_{g_\theta} = T^{(q)} \bigr\} \\
@V{\text{jet evaluation}}VV
  @VV{q = 1 + \lvert\langle L_\theta, F_A^+\rangle\rvert}V \\
\Gamma(E_{\mathrm{lin}}) @>{\nabla^A}>>
  \bigl\{ \psi : J\hbar_q\,\iota_{\xi_q}D_A\psi = \hat{H}_q\psi \bigr\}
\end{CD}
\]
where the top row sends a jet (a first-order perturbation of the gauge field)
to the Einstein-type field equation via the Euler--Lagrange operator; the
left column evaluates a jet at a fibre point to produce a linearised section;
the right column is the coupling through the substrate field $q$; and the
bottom row is the parallel-transport equation whose linearisation is the
Schr\"odinger equation.
\end{proposition}

\begin{proof}[Proof sketch]
The top row is the variational equation $\delta S_q / \delta L_\theta = 0$,
which involves $q(p)$ through the functional $S_q$.  The bottom row is the
$q$-deformed parallel-transport equation
$J\hbar_q\,\iota_{\xi_q}D_A\psi = \hat{H}_q\psi$
(Remark~\ref{rem:schrodinger-jet-coaction}), linearised at a chosen section.
The left column is the jet evaluation map
$j^1_p \sigma \mapsto \sigma(p)$ restricted to $E_{\mathrm{lin}}$.
Commutativity amounts to the statement that solving the variational
equation for $L_\theta$ and then substituting $q(p)$ into the transport
equation gives the same solution set $\psi$ as first evaluating the jet
and then transporting along $A$.  This follows from the chain rule:
$\mathbf{E}[S_q]$ determines $L_\theta$ (and hence $q$), which in turn
determines $\hat{H}_q$ and $\xi_q$ appearing in the transport equation.
A full proof requires the variational bicomplex machinery of
Chapter~\ref{chap:experimental} and is given there.
\end{proof}



\section{Preview: Lifting and Descent}
\label{sec:preview-lifting}

The two central structural results of this paper are:
\begin{itemize}
  \item \textbf{Jet Prolongation (Lifting).}  A $\mathbb{C}$ quantum
        observable universe $(M, G, P, E)$ lifts canonically to a new
        quantum observable universe on the total space of the first jet
        bundle $J^1 E$ (Proposition~\ref{prop:jet-prolongation},
        Chapter~\ref{chap:torsion}).
  \item \textbf{Brane Tower Descent.}  Iterating fixed-point reductions
        under the $\mathrm{SU}(2)$ action produces a tower of nested
        quantum observable universes $(M_k)_{k \geq 0}$, terminating at a
        discrete set of algebraic invariants
        (Section~\ref{ssec:descent}, Chapter~\ref{chap:torsion}).
\end{itemize}
The categorical and analytic elaborations needed for these theorems ---
Tannaka--Krein reconstruction, Hopf--Frobenius diagrams, and the
$q$-deformed operator algebra --- are developed in
Chapter~\ref{chap:elaboration}.  The reader wishing to proceed directly
from the definitions above to geometric applications may skip to
Chapter~\ref{chap:brings-surface}.
 \chapter{Elaboration: Categorical and Analytic Foundations}
\label{chap:elaboration}

This chapter collects the categorical, algebraic, and analytic details
that underpin the definitions of Chapter~\ref{chap:qou}.  The reader
who has read Chapter~\ref{chap:qou} may return to specific sections as
needed; the remainder of the paper does not require this chapter in
linear order.

\begin{remark}[TQFT context]
\label{rem:tqft-context}
The axioms of
Definitions~\ref{def:quantum-universe}
and~\ref{def:quantum-observable-universe} are grounded in the algebraic
framework of topological quantum field theory (TQFT).

\medskip\noindent\textit{Frobenius--Hopf objects and 2d TQFT.}\;
A commutative Frobenius algebra over~$k$ is equivalent to a
$2$-dimensional TQFT
$Z \colon \mathbf{2Cob} \to \mathrm{Vect}_k$
\cite{kock2004, atiyah1988}.  The Frobenius--Hopf object~$A$ in
Definition~\ref{def:quantum-universe} carries both the Frobenius structure
(encoding the $2$d cobordism data: the product~$\mu$ corresponds to a
pair-of-pants cobordism, the Frobenius form~$\lambda$ to a cap) and a
Hopf structure (encoding the quantum-group symmetry that extends the theory
to three dimensions).  Concretely, a $2$d TQFT assigns to
the circle~$S^1$ a Frobenius algebra~$A$ and to each cobordism
$\Sigma \colon S^1 \sqcup \cdots \to S^1 \sqcup \cdots$ a linear map
built from~$\mu$, $\Delta_\mu$, $\eta$, and~$\lambda$.  The Hopf
structure~$(\Delta, \varepsilon, S)$ equips~$A$ with the additional data
needed for three-dimensional and extended TQFTs
\cite{turaev2010, witten1988}.

\medskip\noindent\textit{The observability condition as a categorical
contact condition.}\;
In the TQFT setting, the condition
$\mathrm{Hom}_{\mathcal{C}}(\mathbf{1}, \mathrm{Ad}(A)) = 0$
(Definition~\ref{def:quantum-observable-universe}) is the categorical
analogue of the \emph{contact condition} (maximal non-integrability).
A non-zero morphism $\mathbf{1} \to \mathrm{Ad}(A)$ would give a
global invariant section of the adjoint object --- the categorical
counterpart of a nowhere-vanishing section of the adjoint bundle --- and
would force the TQFT to decompose into a direct sum of
lower-dimensional theories.  Requiring
$\mathrm{Hom}(\mathbf{1}, \mathrm{Ad}(A)) = 0$ therefore ensures:
\begin{enumerate}
  \item \emph{irreducibility}: the associated $3$-dimensional TQFT
        (Chern--Simons type, with gauge datum~$A$ and quantum
        parameter~$q$) does not split as a direct sum;
  \item \emph{non-degeneracy across cobordisms}: the TQFT pairing on
        state spaces assigned to surfaces is non-degenerate
        \cite{freed2010};
  \item \emph{full dualisability}: in the sense of the cobordism
        hypothesis \cite{baez1995, lurie2009}, the object~$A$ is
        fully dualisable in the relevant Morita higher category, so the
        TQFT extends to a fully local (extended) field theory.
\end{enumerate}
In the $\mathbb{C}$ realisation
(Example~\ref{def:c-quantum-observable-universe}), this categorical
condition reduces to
$H^0(M, \mathrm{Ad}(P)) = 0$ --- the classical requirement that the
instanton connection admits no infinitesimal automorphisms, which is the
gauge-theoretic manifestation of the contact (maximally non-integrable)
geometry of $(M, \theta)$.
\end{remark}

\begin{remark}[Vanishing conditions: explication]
\label{rem:vanishing-explication}
In Example~\ref{def:c-quantum-observable-universe},
$\mathrm{Ad}(P) = P \times_G \mathfrak{su}(2)$ is the adjoint bundle
(a linear bundle, so sheaf cohomology is defined), and $\Gamma_c(M,E)$
denotes the space of compactly supported continuous sections of the
(non-linear) fibre bundle $E$.  The condition $H^0(M, \mathrm{Ad}(P)) = 0$
asserts the vanishing of \emph{all} global parallel sections of the adjoint
bundle --- not merely compactly supported ones --- eliminating every
infinitesimal gauge symmetry.
\end{remark}

\begin{remark}[Generalisation beyond strict pseudoconvexity]
\label{rem:levi-generalisation}
Example~\ref{def:c-quantum-universe} requires only Levi non-degeneracy
(not strict pseudoconvexity).  This generality is deliberate: the QOU
framework applies whenever the following three cohomological/obstruction
conditions hold.

\begin{enumerate}
  \item \emph{Kohn--Rossi finite-dimensionality.}\;
        The tangential CR complex
        $(\Omega^{0,\bullet}_b(M),\, \bar\partial_b)$ has
        finite-dimensional cohomology
        $H^{0,k}_{\bar\partial_b}(M)$ for all~$k$.  This replaces
        strict pseudoconvexity in ensuring analytic control over
        $q$-harmonic forms (Definition~\ref{def:q-harmonic-form}).
        In the strictly pseudoconvex case, Kohn's sub-elliptic
        estimates guarantee this automatically; for indefinite
        Levi signature, one imposes it as a hypothesis.
  \item \emph{Non-degeneracy of the substrate.}\;
        The Levi form $L_\theta$ is non-degenerate (but not necessarily
        positive definite), so the substrate parameter
        $q(p) = 1 + \lvert\langle L_\theta, F_A^+\rangle\rvert$
        (Remark~\ref{rem:c-substrate-via-levi}) is well-defined and
        the regime $q > 1$ is attainable.  When the Levi form has
        indefinite signature $(r,s)$ with $r,s \geq 1$, the Webster
        metric $g_\theta$ is pseudo-Riemannian; the contraction
        $\langle L_\theta, F_A^+\rangle$ remains well-defined (using
        the pseudo-Riemannian inner product on $\Lambda^+ H^*$), and
        one takes the absolute value to ensure $q \geq 1$.
  \item \emph{QOU vanishing (Example~\ref{def:c-quantum-observable-universe}).}\;
        The conditions $H^0(M, \mathrm{Ad}(P)) = 0$ and
        $\Gamma_c(M,E) = 0$ serve as the gauge-theoretic obstructions.
        These are purely topological/cohomological and independent of the
        Levi signature.
\end{enumerate}

\noindent
Together, conditions~(1)--(3) ensure that all results of this chapter
(irreducibility of the connection, self-duality of the curvature,
substrate dynamics, $q$-deformed operators) extend to CR manifolds of
arbitrary Levi signature.  The strictly pseudoconvex case (signature
$(n,0)$) is the simplest and motivates all examples in this chapter;
the indefinite case arises naturally for real hypersurfaces in
indefinite K\"ahler manifolds and for boundaries of
Grauert tubes over pseudo-Riemannian manifolds.
\end{remark}

\begin{remark}[Jet-bundle formulation]
\label{rem:jet-bundle}
The conditions in Example~\ref{def:c-quantum-observable-universe} admit a
natural reformulation in terms of the first jet bundle $J^1 E$.  The vanishing
$\Gamma_c(M, E) = 0$ ensures that no non-trivial compactly supported section of
$E$ exists, while $H^0(M, \mathrm{Ad}(P)) = 0$ eliminates \emph{all}
infinitesimal gauge symmetries (not merely compactly supported ones).  Together,
these conditions mean that the space of sections of $J^1 E$ injects into the
space of connections on $P$, preventing any local degeneracy of the gauge field.
\end{remark}

\begin{example}[Real hypersurface in a Calabi--Yau $2$-fold]
\label{ex:calabi-yau}
Let $X$ be a Calabi--Yau $2$-fold (K3 surface) with holonomy
$\mathrm{SU}(2)$, and let $M \subset X$ be a smooth, strictly convex real
hypersurface.  Then $M$ is a $3$-dimensional strictly pseudoconvex CR
manifold inheriting its CR structure from the ambient K\"ahler geometry.
Because the holonomy of $X$ is $\mathrm{SU}(2)$, the frame bundle of $X$
restricted to $M$ yields a natural principal $\mathrm{SU}(2)$-bundle
$P \to M$.  The associated state bundle $E = P \times_G \mathbb{CP}^1$
yields a $\mathbb{C}$ quantum observable universe $(M, \theta, G, P, E)$
when the vanishing conditions of
Example~\ref{def:c-quantum-observable-universe} are satisfied.  (For
higher-dimensional Calabi--Yau manifolds with holonomy $\mathrm{SU}(n)$,
$n > 2$, a reduction to $\mathrm{SU}(2)$ must be supplied separately.)
\end{example}

\begin{example}[ALE links]
\label{ex:ale}
Let $\Gamma \subset \mathrm{SU}(2)$ be a finite subgroup and let $X$ be the
minimal resolution of $\mathbb{C}^2 / \Gamma$.  The link
$M = S^3 / \Gamma$ is the boundary of a tubular neighbourhood of the
exceptional set in $X$; it is a compact $3$-dimensional strictly pseudoconvex
CR manifold.  The CR structure of the boundary induces a contact form
$\theta$ on $M$ whose Reeb vector field $\xi$ generates the Hopf-type
fibration $S^3/\Gamma \to S^2$.  (When $\Gamma \neq \{1\}$, the target is
the orbifold $\mathbb{CP}^1/\Gamma$; the bundle structure remains
well-defined on the smooth manifold $M = S^3/\Gamma$.)  Equipped with the
induced principal $\mathrm{SU}(2)$-bundle $P \to M$ and the associated
state bundle $E = P \times_G \mathbb{CP}^1$, the tuple
$(M, \theta, G, P, E)$ provides a $\mathbb{C}$ quantum observable universe.
Note that $\dim_{\mathbb{R}} H = 2$, so the Hodge-star decomposition
$\Lambda^2 H^* = \Lambda^+\! H^* \oplus \Lambda^-\! H^*$ is trivial
($\Lambda^2 H^* \cong \mathbb{R}$); in this dimension the curvature
$F_A|_H$ is automatically ``self-dual'' in the vacuous sense, and the
instanton condition reduces to the Yang--Mills equation
(Definition~\ref{def:instanton}).
\end{example}

\begin{example}[Bring's surface unit circle bundle]
\label{ex:brings-surface}
Let $S$ denote Bring's surface (Definition~\ref{def:brings-surface}), the
genus-$4$ Riemann surface in $\mathbb{P}^4$ cut out by the first three power
sums $\sum x_j^k = 0$, $k=1,2,3$.  Being a Riemann surface (real
dimension~$2$), $S$ itself is not a CR manifold.  The CR geometry lives on the
unit circle bundle $M_S \subset T^*S$: the Liouville $1$-form on $T^*S$
restricts to a contact form $\theta$ on $M_S$, making it a $3$-dimensional
strictly pseudoconvex CR manifold.  The Reeb vector field $\xi$ generates the
$U(1)$ fibre rotation of $M_S \to S$.  The natural projection
$\pi\colon S \to \mathbb{CP}^1$ (a degree-$4$ branched covering with $6$
branch points corresponding to the $6$ defects of Euler
characteristic~$\chi=-6$) lifts to a principal $\mathrm{SU}(2)$-bundle
$P \to M_S$.  Setting $E = P \times_G \mathbb{CP}^1$, the $4$-tuple
$(M_S,\, G,\, P,\, E)$ assembles into a $\mathbb{C}$ quantum observable
universe whose exceptional divisors support the $q$-deformed Reeb flow
(Definition~\ref{def:q-reeb-flow}).
\end{example}

\begin{remark}[$\mathbb{C}$ examples as instances of the generalised framework]
\label{rem:examples-unified}
Each of the preceding $\mathbb{C}$ examples is an instance of the
base-field-independent Definition~\ref{def:quantum-universe}, realised
over $k = \mathbb{C}$ via the Levi non-degenerate framework of
Example~\ref{def:c-quantum-universe} and
Remark~\ref{rem:levi-generalisation}.
\begin{itemize}
  \item \emph{CY hypersurfaces} (Example~\ref{ex:calabi-yau}).  A smooth
        strictly convex real hypersurface $M$ in a K\"ahler manifold inherits
        a strictly pseudoconvex CR structure --- i.e.\ the Levi form has
        definite signature.
        The Kohn--Rossi condition holds automatically by Kohn's sub-elliptic
        estimates.
  \item \emph{ALE links} (Example~\ref{ex:ale}).  The link
        $S^3/\Gamma$ is the CR boundary of the minimal resolution of
        $\mathbb{C}^2/\Gamma$; its Levi form has signature $(1,0)$.
        Kohn--Rossi finite-dimensionality follows from compactness of~$M$.
  \item \emph{Bring's surface unit circle bundle}
        (Example~\ref{ex:brings-surface}).  The unit circle bundle
        $M_S \subset T^*S$ is a compact $3$-dimensional strictly
        pseudoconvex CR manifold with Levi signature~$(1,0)$.
\end{itemize}
All three satisfy the three cohomological conditions of
Remark~\ref{rem:levi-generalisation}: Kohn--Rossi
finite-dimensionality (by strict pseudoconvexity or compactness),
substrate non-degeneracy (positive-definite Levi form ensures
$q > 1$ wherever $F_A^+ \neq 0$), and the QOU vanishing conditions
(by the topology of the respective principal bundles).  More generally,
any real hypersurface in an indefinite K\"ahler manifold, or the boundary
of a Grauert tube over a pseudo-Riemannian manifold, yields a CR manifold
with indefinite Levi signature satisfying these conditions --- extending
the QOU framework beyond the strictly pseudoconvex setting.
\end{remark}

\begin{remark}[Quantum gauge group: category-theoretic perspective]
\label{rem:quantum-gauge-group-category}
The quantum gauge group
(Definition~\ref{def:quantum-gauge-group}) is realised pointwise by
the Hopf algebra $\mathcal{O}_{q(p)}(\mathrm{SU}(2))$ of regular
functions on $\mathrm{SU}_{q(p)}(2)$.  The categorical view replaces
each fibre with its representation category
$\mathrm{Rep}(\mathrm{SU}_{q(p)}(2))$, a braided monoidal category.
Tannaka--Krein duality reconstructs $\mathcal{O}_{q(p)}$ from
$\mathrm{Rep}(\mathrm{SU}_{q(p)}(2))$ via
$\mathrm{End}^{\otimes}(\omega)$, where $\omega$ is the forgetful
fibre functor.  For any two points $p, p' \in M$ connected by a
path~$\gamma$, the passage between fibres is captured by the
commutative square
\begin{equation}\label{eq:tannaka-krein-functoriality}
\begin{CD}
\mathcal{O}_{q(p)}(\mathrm{SU}(2))
  @>{\mathrm{Rep}}>>
  \mathrm{Rep}(\mathrm{SU}_{q(p)}(2)) \\
@V{\text{transition along } \gamma}VV
  @VV{\text{monoidal functor}}V \\
\mathcal{O}_{q(p')}(\mathrm{SU}(2))
  @>>{\mathrm{Rep}}>
  \mathrm{Rep}(\mathrm{SU}_{q(p')}(2))
\end{CD}
\end{equation}
Tannaka--Krein reconstruction inverts the horizontal arrows, so the
Hopf-algebra and representation-category viewpoints are equivalent
fibrewise.
\end{remark}

\begin{remark}[Hopf-algebra structure and non-cocommutativity]
\label{rem:hopf-structure-diagrams}
Write $\mathcal{O}_q = \mathcal{O}_q(\mathrm{SU}(2))$ and let
$\Delta_q$, $\varepsilon_q$, $S_q$ denote the comultiplication,
counit, and antipode.  The Hopf comultiplication satisfies
\emph{coassociativity}:
\begin{equation}\label{eq:coassociativity}
\begin{CD}
\mathcal{O}_q @>{\Delta_q}>>
  \mathcal{O}_q \otimes \mathcal{O}_q \\
@V{\Delta_q}VV
  @VV{\mathrm{id} \otimes \Delta_q}V \\
\mathcal{O}_q \otimes \mathcal{O}_q
  @>>{\Delta_q \otimes \mathrm{id}}>
  \mathcal{O}_q \otimes \mathcal{O}_q \otimes \mathcal{O}_q
\end{CD}
\end{equation}
Let $\sigma \colon \mathcal{O}_q \otimes \mathcal{O}_q \to
\mathcal{O}_q \otimes \mathcal{O}_q$ denote the tensor flip
$(a \otimes b \mapsto b \otimes a)$.  The cocommutativity square
\begin{equation}\label{eq:non-cocommutativity}
\begin{CD}
\mathcal{O}_q @>{\Delta_q}>>
  \mathcal{O}_q \otimes \mathcal{O}_q \\
@V{\Delta_q}VV  @VV{\sigma}V \\
\mathcal{O}_q \otimes \mathcal{O}_q
  @=
  \mathcal{O}_q \otimes \mathcal{O}_q
\end{CD}
\quad\text{commutes}\;\Longleftrightarrow\;q = 1.
\end{equation}
For $q \neq 1$ the failure of~\eqref{eq:non-cocommutativity} is
controlled by the universal $R$-matrix of $U_q(\mathfrak{su}(2))$:
$\sigma \circ \Delta_q = \mathrm{Ad}_{R_q} \circ \Delta_q$,
where $R_1 = 1 \otimes 1$ recovers cocommutativity at $q = 1$.
\end{remark}

\begin{remark}[Self-duality as a short exact sequence]
\label{rem:self-dual-ses}
In the QOU setting with $\dim_{\mathbb{R}} H = 4$, the Hodge star
$*_H$ gives the splitting
$\Lambda^2 H^* = \Lambda^+\!H^* \oplus \Lambda^-\!H^*$ and hence
the short exact sequence of $\mathrm{Ad}(P)$-valued forms
\begin{equation}\label{eq:self-dual-ses}
\begin{CD}
0 @>>>
  \Gamma\!\bigl(\Lambda^-\!H^* \otimes \mathrm{Ad}(P)\bigr)
  @>{\iota}>>
  \Gamma\!\bigl(\Lambda^2 H^* \otimes \mathrm{Ad}(P)\bigr)
  @>{\pi^+}>>
  \Gamma\!\bigl(\Lambda^+\!H^* \otimes \mathrm{Ad}(P)\bigr)
  @>>> 0
\end{CD}
\end{equation}
The instanton condition
$F_A|_H = {*_H}\,F_A|_H$
(Proposition~\ref{prop:self-dual-curvature}) is exactly the condition
$F_A|_H \in \ker(\mathrm{id} - {*_H})$,
i.e.\ $F_A^-|_H = 0$.  In the language of $2$-groups, this has a
clean reinterpretation
(Remark~\ref{rem:self-dual-2group}).
\end{remark}

\begin{remark}[Self-duality as constrained $2$-group gauge theory]
\label{rem:self-dual-2group}
A strict $2$-group is equivalently a \emph{crossed module}: a
group homomorphism $\partial \colon H \to G$ together with an
action of $G$ on $H$ satisfying the Peiffer identity.  The
associated exact sequence
\begin{equation}\label{eq:crossed-module-ses}
\begin{CD}
1 @>>> \ker\partial @>>> H @>{\partial}>> G @>>>
  \mathrm{coker}\,\partial @>>> 1
\end{CD}
\end{equation}
encodes the $2$-group as a homotopy type with
$\pi_0 = \mathrm{coker}\,\partial$ and
$\pi_1 = \ker\partial$.

The \emph{inner-automorphism $2$-group}
$\mathrm{INN}(G) = (G \xrightarrow{\mathrm{id}} G)$
has $\ker\partial = 1$ and $\mathrm{coker}\,\partial = 1$, so it is
connected and simply connected as a $2$-group.  A $2$-connection for
$\mathrm{INN}(G)$ on a principal $2$-bundle over~$M$ consists of a
$\mathfrak{g}$-valued $1$-form~$A$ and a $\mathfrak{g}$-valued
$2$-form~$B$, subject to the \emph{fake-flatness} condition
\begin{equation}\label{eq:fake-flatness}
F_A \;=\; B.
\end{equation}
In the QOU setting, $G = \mathrm{SU}(2)$ and the Levi distribution
$H = \ker\theta$ has $\dim_{\mathbb{R}} H = 4$, so the Hodge
decomposition $\Lambda^2 H^* = \Lambda^+\!H^* \oplus \Lambda^-\!H^*$
is available.  Imposing the \emph{simplicity constraint}
\[
  B \;\in\; \Gamma\!\bigl(\Lambda^+\!H^* \otimes \mathrm{Ad}(P)\bigr)
\]
--- that is, requiring the $2$-form datum to be self-dual ---
and combining with fake flatness~\eqref{eq:fake-flatness} yields
the commutative diagram
\begin{equation}\label{eq:2group-instanton}
\begin{CD}
\Omega^1(M,\,\mathfrak{g})
  @>{d_A}>>
  \Omega^2(M,\,\mathfrak{g})
  @>{\pi^+}>>
  \Omega^{2,+}(M,\,\mathfrak{g}) \\
@.  @AA{=\;B}A  @AA{\text{simplicity}}A \\
@.  F_A @= F_A^+
\end{CD}
\end{equation}
The fake-flatness condition $F_A = B$ combined with the simplicity
constraint $B = F_A^+$ forces $F_A = F_A^+$, hence $F_A^- = 0$.
This is precisely the instanton equation
$F_A|_H = {*_H}\,F_A|_H$
of Definition~\ref{def:instanton}.  Thus, the self-dual formulation
used throughout this paper is a special case of $2$-group gauge theory
for $\mathrm{INN}(\mathrm{SU}(2))$ subject to a simplicity constraint
on the $2$-form datum.

This perspective clarifies the relationship between the instanton
condition in Chapter~1 and the higher-categorical structures of later
chapters: the parallel-transport $2$-functor of the
$\mathrm{INN}(\mathrm{SU}(2))$ $2$-bundle assigns to paths in $M$
not merely group elements in $\mathrm{SU}(2)$ but $2$-morphisms
encoding the self-dual curvature, with the simplicity constraint
selecting the physical (instanton) sector.
\end{remark}

\begin{remark}[Frobenius structure of the $\mathbb{C}$ gauge algebra]
\label{rem:frobenius-simplification}
The Hopf algebra $\mathcal{O}_q(\mathrm{SU}(2))$ carries a natural
\emph{Frobenius structure}: the Haar state
$h_q \colon \mathcal{O}_q \to \mathbb{C}$ induces a non-degenerate
pairing $\langle a, b \rangle = h_q(ab)$, making $\mathcal{O}_q$ into
a Frobenius algebra with comultiplication $\Delta_\mu$ dual to the
multiplication.  This Frobenius comultiplication is compatible with the
Hopf comultiplication~$\Delta_q$ of~\eqref{eq:coassociativity} and
satisfies the \emph{Frobenius relation}:
\begin{equation}\label{eq:frobenius-relation}
\begin{CD}
\mathcal{O}_q \otimes \mathcal{O}_q
  @>{\mu}>>
  \mathcal{O}_q \\
@V{\mathrm{id} \otimes \Delta_\mu}VV
  @VV{\Delta_\mu}V \\
\mathcal{O}_q \otimes \mathcal{O}_q \otimes \mathcal{O}_q
  @>>{\mu \otimes \mathrm{id}}>
  \mathcal{O}_q \otimes \mathcal{O}_q
\end{CD}
\end{equation}
where $\mu$ is the algebra multiplication and $\Delta_\mu$ the
Frobenius comultiplication (dual to $\mu$ under the Haar pairing).
The non-degeneracy of $h_q$ is expressed by requiring the composite
\begin{equation}\label{eq:frobenius-nondegen}
\begin{CD}
\mathcal{O}_q \otimes \mathcal{O}_q
  @>{\mu}>>
  \mathcal{O}_q
  @>{h_q}>>
  \mathbb{C}
\end{CD}
\end{equation}
to be a perfect pairing, i.e.\ the induced map
$\mathcal{O}_q \to \mathcal{O}_q^*$ is an isomorphism.

This Frobenius--Hopf structure on $\mathcal{O}_q(\mathrm{SU}(2))$ is
precisely the datum required by the general
Definition~\ref{def:quantum-universe}: the Haar state $h_q$ realises
the Frobenius form $\lambda \colon A \to \mathbf{1}$, and the
non-degeneracy~\eqref{eq:frobenius-nondegen} is the $\mathbb{C}$
instance of condition~(1) in
Definition~\ref{def:quantum-observable-universe}.
In the $\mathbb{C}$ realisation (Example~\ref{def:c-quantum-universe}),
the $5$-tuple $(M, \theta, G, P, E)$ and the two analytic vanishing
conditions of Example~\ref{def:c-quantum-observable-universe}
reduce to a triple:
\begin{enumerate}
  \item a CR manifold $(M, \theta)$ of hypersurface type with
        non-degenerate Levi form;
  \item a non-degenerate Frobenius--Hopf object~$A$ in the monoidal
        category of vector bundles over~$M$, playing the r\^ole of the
        triple $(G, P, E)$; and
  \item an instanton connection on the underlying $A$-bundle
        (Definition~\ref{def:instanton}).
\end{enumerate}
The vanishing condition $H^0(M, \mathrm{Ad}(P)) = 0$ then follows from
the non-degeneracy of the Frobenius form~$\lambda$ via the diagram
\begin{equation}\label{eq:frobenius-vanishing}
\begin{CD}
\mathrm{Ad}(A) \otimes \mathrm{Ad}(A)
  @>{\mu_{\mathrm{Ad}}}>>
  \mathrm{Ad}(A)
  @>{\lambda_{\mathrm{Ad}}}>>
  \mathcal{O}_M \\
@V{\text{global sections}}VV  @.  @VV{\int_M}V \\
H^0(M,\,\mathrm{Ad}(A)) \otimes H^0(M,\,\mathrm{Ad}(A))
  @.  @>>{\text{pairing}}>
  \mathbb{C}
\end{CD}
\end{equation}
The adjoint representation $\mathrm{Ad}(A)$ inherits a non-degenerate
Frobenius pairing from~$\lambda$; any global parallel section would
pair non-trivially with itself, contradicting the irreducibility of
the instanton connection
(Lemma~\ref{lem:irreducible-connection}).  Hence the Frobenius
non-degeneracy, combined with irreducibility, forces
$H^0(M, \mathrm{Ad}(A)) = 0$.  The section-vanishing condition
$\Gamma_c(M, E) = 0$ follows similarly from the non-existence of
$A$-equivariant sections of the associated $\mathbb{CP}^1$-bundle.

Note that the horizontal distribution $H = \ker\theta$ itself
\emph{fails} to carry a Frobenius structure (as discussed in
Remark~\ref{rem:q-torsion-existence}); it is the gauge algebra, not the
distribution, that is Frobenius.
\end{remark}

\begin{example}[Cosmological limits of $q$ on $S^3$]
\label{ex:cosmological-limits}
Consider the link $M = S^3$ with the standard contact form
$\theta$ (Example~\ref{ex:ale} with $\Gamma = \{1\}$) and let
$A$ be the standard instanton connection on the Hopf bundle
$P = S^7 \to S^4$ restricted to $S^3$.  The Levi form $L_\theta$
has signature $(1,0)$ with $\lvert L_\theta \rvert = 1$ everywhere.

As the instanton curvature concentrates near a point $p_0 \in S^3$
(a ``collapsing instanton''),
$\lvert \langle L_\theta, F_A^+ \rangle \rvert$ diverges at $p_0$ and
tends to zero away from $p_0$.  This produces two regimes simultaneously:
\begin{itemize}
  \item \emph{Near $p_0$}: $q(p) \to \infty$, $\hbar_q \to 1$,
        $[\hat{X},\hat{P}]_q \to J$ --- a ``big bang'' singularity.
  \item \emph{Away from $p_0$}: $q(p) \to 1$, $\hbar_q \to 0$,
        $[\hat{X},\hat{P}]_q \to 0$ --- a ``heat death'' region.
\end{itemize}
The substrate field $q$ thus captures the phase transition between
quantum and classical regimes as a geometric phenomenon.
\end{example}

\begin{remark}[$J^1E$ as the space of local states]
\label{rem:j1e-local-states}
The jet bundle $J^1 E$ serves as the space of local states of the
$\mathbb{C}$~quantum observable universe.  The classical gauge dynamics,
encoded by the \emph{instanton connection} $A$
(Lemma~\ref{lem:irreducible-connection}), are weighted by the substrate
parameter $q$ (Definition~\ref{def:substrate-parameter}) to define the
quantum field algebra $\mathrm{SU}_q(2)$
(Definition~\ref{def:quantum-gauge-group}).
\end{remark}

\begin{proposition}[$d_q^2$ does not vanish]
\label{prop:dq-squared}
Let $d_q$ be the $q$-deformed exterior derivative
(Definition~\ref{def:q-ext-deriv}).  For constant $q > 1$, the operator
$d_q^2$ is given by
\[
  d_q^2\,\alpha
  \;=\; q^{-1}(1 - q^{-1})\,d\theta \wedge \mathcal{L}_\xi\,\alpha.
\]
In particular, $d_q^2 = 0$ if and only if $q = 1$ (classical case) or
$\alpha$ lies in the kernel of $\mathcal{L}_\xi$ (Reeb-invariant forms).
On horizontal $\mathcal{L}_\xi$-invariant forms, $d_q$ \emph{does}
define a cochain complex.
\end{proposition}

\begin{proof}
Write $a = q^{-1}$ and $b = 1 - q^{-1}$, so $a + b = 1$.  Then
\begin{align*}
  d_q^2\,\alpha
  &= d_q\bigl(a\,d\alpha + b\,\theta \wedge \mathcal{L}_\xi\alpha\bigr) \\
  &= a^2\,d^2\alpha
     + ab\,d(\theta \wedge \mathcal{L}_\xi\alpha)
     + ab\,\theta \wedge \mathcal{L}_\xi\,d\alpha
     + b^2\,\theta \wedge \mathcal{L}_\xi(\theta \wedge \mathcal{L}_\xi\alpha).
\end{align*}
The first term vanishes since $d^2 = 0$.  The second expands as
$ab\,d\theta \wedge \mathcal{L}_\xi\alpha
- ab\,\theta \wedge d\mathcal{L}_\xi\alpha$.
The third term cancels the second's negative part because
$[d,\, \mathcal{L}_\xi] = 0$ implies
$\theta \wedge \mathcal{L}_\xi\,d\alpha
= \theta \wedge d\,\mathcal{L}_\xi\alpha$.
For the fourth term, the Reeb field satisfies
$\mathcal{L}_\xi\theta = 0$ (since $\iota_\xi d\theta = 0$ and
$\theta(\xi)=1$), so
$\mathcal{L}_\xi(\theta \wedge \mathcal{L}_\xi\alpha)
= \theta \wedge \mathcal{L}_\xi^2\alpha$;
the factor $\theta \wedge \theta = 0$ in $b^2\,\theta \wedge \theta
\wedge \mathcal{L}_\xi^2\alpha$ kills this term.  Hence
$d_q^2\alpha = ab\,d\theta \wedge \mathcal{L}_\xi\alpha$.
\end{proof}

\begin{proposition}[Jacobiator of ${{[\,\cdot\,,\cdot\,]}}_q$]
\label{prop:q-jacobiator}
Let $U, V, W \in \Gamma(H)$ be horizontal vector fields (i.e.\
$\theta(U) = \theta(V) = \theta(W) = 0$).  The \emph{Jacobiator}
\[
  \mathrm{Jac}_q(U,V,W)
  \;:=\; [[U,V]_q, W]_q + [[V,W]_q, U]_q + [[W,U]_q, V]_q
\]
is given by
\[
  \mathrm{Jac}_q(U,V,W)
  \;=\; -q^{-1}(1-q^{-1})\!\sum_{\mathrm{cyc}}
        d\theta(V,W)\,\mathcal{L}_\xi U,
\]
where the sum is over cyclic permutations of $(U,V,W)$.  In particular,
$\mathrm{Jac}_q = 0$ if and only if $q = 1$ (classical Jacobi identity)
or $d\theta|_{H \times H} = 0$ (integrable horizontal distribution).
Since $d\theta|_{H \times H}$ is the Levi form $L_\theta$, which is
non-degenerate by hypothesis, $[\,\cdot\,,\cdot\,]_q$ violates the Jacobi
identity for all $q \neq 1$.  Thus $(\mathfrak{X}(M), [\,\cdot\,,\cdot\,]_q)$
is a \emph{Leibniz algebra} (or a pre-Lie algebra), not a Lie algebra.
\end{proposition}

\begin{proof}
For horizontal $V, W$, the bracket reduces to $[V,W]_q = q^{-1}[V,W]$.
The mixed term $\theta([V,W]_q) = -q^{-1}\,d\theta(V,W)$ is in general
non-zero.  Substituting into the next bracket:
\begin{align*}
  [[V,W]_q,\, U]_q
  &= q^{-2}\,[[V,W],U]
    - q^{-1}(1-q^{-1})\,d\theta(V,W)\,\mathcal{L}_\xi U,
\end{align*}
using $\theta(U)=0$ and $\theta([V,W]_q) = -q^{-1}d\theta(V,W)$.
Summing cyclically, the $q^{-2}$ terms assemble into $q^{-2}\,\mathrm{Jac}(U,V,W) = 0$
by the classical Jacobi identity, leaving only the $d\theta$ terms.
\end{proof}

\begin{remark}[Classical recovery of $q$-deformed operators]
\label{rem:classical-recovery}
All four $q$-deformed operators interpolate between the classical ($q = 1$)
and deformed ($q > 1$) regimes via the uniform pattern
$q^{-1} \cdot (\text{classical}) + (1 - q^{-1}) \cdot (\text{Reeb correction})$.
At $q = 1$ each recovers its classical counterpart:
$d_q = d$ (Definition~\ref{def:q-ext-deriv}),
$[\,\cdot\,,\cdot\,]_q = [\,\cdot\,,\cdot\,]$
(Definition~\ref{def:q-lie-bracket}),
$A_q = A$ (Definition~\ref{def:q-connection}), and
$\xi_q = \xi$ (Definition~\ref{def:q-reeb-vector-field}).
The connection $A_q$ interpolates between the full instanton connection and
its Reeb component $\iota_\xi A\;\theta$.
\end{remark}

\begin{remark}[$q$-Torsion leaves as energy minima]
\label{rem:q-torsion-leaf-minimum}
Let $A_q$ denote the $q$-deformed connection
(Definition~\ref{def:q-connection}) and define the $q$-curvature
\[
  \Omega_q \;=\; dA_q + \tfrac{1}{2}\,[A_q,\, A_q]_q,
\]
where $[\,\cdot\,,\cdot\,]_q$ is the $q$-deformed Lie bracket
(Definition~\ref{def:q-lie-bracket}) extended to $\mathrm{Ad}(P)$-valued
forms.  A $q$-torsion leaf $L$ (Definition~\ref{def:q-torsion-leaf}) is a
local minimum of $\lVert\Omega_q\rVert$ restricted to the horizontal
distribution: the invariance of $L$ under $\phi_t^q$ forces the first
variation of $\lVert\Omega_q|_L\rVert$ to vanish, so $L$ sits in a
$q$-curvature well.  Physically, a $q$-torsion leaf represents a region of
the CR manifold where the gauge curvature is minimised and the $q$-deformed
dynamics are in equilibrium.
\end{remark}

\begin{remark}[Physical interpretation of $q$-harmonicity]
\label{rem:q-harmonic-physical}
The $q$-Laplacian $\Delta_q$ encodes the $q$-deformed viscosity of the
CR geometry.  A $q$-harmonic form is a zero-mode of this operator and thus
represents a stable excitation --- a particle --- of the quantum observable
universe.  The deformation parameter $q$ interpolates between the classical
regime ($q \to 1$, where $\Delta_q$ reduces to the Hodge Laplacian) and the
deep-quantum regime ($q \to \infty$), so the notion of ``particle'' is itself
$q$-dependent.
\end{remark}

\begin{remark}[$q$-Deformed Schr\"odinger equation]
\label{rem:schrodinger-jet-coaction}
The $q$-deformed Schr\"odinger equation governs the evolution of states
in the linearised state bundle.  Throughout this paper the symbol~$i$ in
the Schr\"odinger equation is identified with the complex structure
operator $J \colon H \to H$ on the Levi distribution ($J^2 = -\mathrm{Id}$):
it acts as a geometric twisting operator on horizontal planes rather than
as an algebraic imaginary unit.

\medskip\noindent\textit{Linearisation convention.}\;
The fibre $\mathbb{CP}^1$ is a non-linear space, so one cannot write a
linear PDE for sections of $E$ directly.  We work instead with the
\emph{linearised} model: identify a neighbourhood of a base-point in
$\mathbb{CP}^1$ with $\mathbb{C}$ via stereographic projection, and
interpret $\psi$ as a section of the line bundle $E_{\mathrm{lin}}$
obtained by linearising $E$ at the chosen section.  The covariant
derivative $D_A = d + A$ then acts on sections of this line bundle in the
standard way.

For a section $\psi \in \Gamma(E_{\mathrm{lin}})$ and the $q$-deformed Reeb
vector field $\xi_q$ (Definition~\ref{def:q-reeb-vector-field}), the
Schr\"odinger dynamics take the form
\[
  J\,\hbar_q(p)\,\iota_{\xi_q} D_A \psi \;=\; \hat{H}_q\, \psi,
\]
where $\hat{H}_q = -\tfrac{1}{2}\hbar_q\,\Delta_q + V_q$ is the
\emph{$q$-deformed Hamiltonian} (here $\Delta_q$ is the $q$-Laplacian of
Definition~\ref{def:q-harmonic-form} acting on sections of
$E_{\mathrm{lin}}$, and $V_q$ is a potential depending on the curvature
$\Omega_q$), and:
\begin{itemize}
  \item $\hbar_q(p) = 1 - q(p)^{-1}$: the quantum Planck constant
        (Definition~\ref{def:quantum-planck-constant}) is a local field;
  \item $D_A = d + A$: the instanton connection on $P$ acts as the potential;
  \item $\iota_{\xi_q} D_A \psi$: the covariant derivative $D_A\psi$ is an
        $E_{\mathrm{lin}}$-valued $1$-form; contracting it with $\xi_q$
        yields a section of $E_{\mathrm{lin}}$, matching the right-hand side.
\end{itemize}
As $q(p) \to 1$ (classicality), the commutator
$[\hat{X}, \hat{P}]_q \to 0$ (Remark~\ref{rem:heisenberg-commutator}), the
uncertainty vanishes, and the wave-like nature of the section freezes into a
classical trajectory.

\medskip\noindent\textit{Categorical role.}\;
Although this remark is stated as a PDE rather than in categorical language,
it is essential to the paper's structure: the Schr\"odinger equation
provides the bottom row of the main commutative diagram
(Proposition~\ref{prop:main-commutative-diagram}).  Categorically, the
$q$-deformed Hamiltonian $\hat{H}_q$ defines an endomorphism of the
module category $\mathrm{Mod}_{E_{\mathrm{lin}}}$ over the jet algebra
$\Gamma(J^1 E)$, and the Schr\"odinger equation is the condition that
$\psi$ lies in the kernel of $\hat{H}_q - J\hbar_q\,\iota_{\xi_q}D_A$.
In the base-field-independent setting of
Definitions~\ref{def:quantum-universe}
and~\ref{def:quantum-observable-universe}, this extends to arbitrary base
fields~$k$ by replacing $\Gamma(E_{\mathrm{lin}})$ with the appropriate
$k$-linear category of sections.
\end{remark}

 \chapter{Path Integrals and Braiding}
\label{chap:path-integrals-braiding}

In the QOU framework, traditional path integrals are replaced by finite sums
over the discrete fibre structure.  Since the base space $B$ and each fibre are
finite, the continuum measure $\mathcal{D}\phi$ collapses to a combinatorial
sum, eliminating the need for renormalisation of divergent integrals.

The braid group acts naturally on the fibre degrees of freedom.  What we
observe as particle interactions corresponds to the action of braid operators
exchanging strands in the fibre bundle.  The algebraic content of these
interactions is therefore encoded in braid-group representations rather than in
Feynman vertices.

The exchange operations carry a $q$-deformation parameter inherited from the
quantum group structure of the fibres.  Concretely, the $R$-matrix of
$\mathcal{U}_q(\mathfrak{su}(2))$ furnishes a $q$-deformed braid operator
whose eigenvalues govern the phase acquired under particle exchange.

Under this correspondence, specific knot types map to particle identities.  The
trefoil knot, being the simplest non-trivial torus knot, is identified with the
quark; the muon arises as a lift of the trefoil to a higher fibre level.  The
unknot, carrying trivial braiding, corresponds to the neutrino.
Experimental evidence supporting this $q$-deformed braiding picture is
discussed in Chapter~\ref{chap:experimental}.



\section{Braid group and symmetric group}
\label{sec:braid-symmetric}

\begin{definition}[Braid group on three strands]
\label{def:braid-group}
The \emph{braid group on three strands} is the group
\begin{equation}\label{eq:braid-presentation}
  B_3
  =\bigl\langle\sigma_1,\sigma_2
    \;\big|\;\sigma_1\sigma_2\sigma_1=\sigma_2\sigma_1\sigma_2\bigr\rangle.
\end{equation}
This is the fundamental group of the configuration space of three
unordered points in the plane, and is isomorphic to the fundamental group
of the complement of the trefoil knot
(Remark~\ref{rem:trefoil-knot-group}).
\end{definition}

\begin{definition}[Symmetric group $S_3$]
\label{def:s3-group}
Imposing the involution relations $\sigma_i^2=\mathrm{id}$ on the braid
group $B_3$ (Definition~\ref{def:braid-group}) yields the
\emph{symmetric group}
\begin{equation}\label{eq:s3-presentation}
  S_3
  =\bigl\langle s_1,s_2
    \;\big|\;s_1^2=s_2^2=(s_1 s_2)^3=\mathrm{id}\bigr\rangle.
\end{equation}
\end{definition}



\section{$\mathrm{SU}(2)$ and $\mathrm{SO}(3)$ framework}
\label{sec:su2-so3}

\begin{definition}[$\mathfrak{su}(2)$ basis]
\label{def:su2-basis}
Write $\tau_1,\tau_2,\tau_3$ for the Pauli matrices and set
\begin{equation}\label{eq:su2-basis}
  e_a=-\tfrac{i}{2}\,\tau_a,\qquad a=1,2,3,
\end{equation}
so that the Lie bracket is
\begin{equation}\label{eq:su2-commutator}
  [e_a,\,e_b]=\epsilon_{abc}\,e_c.
\end{equation}
\end{definition}

\begin{remark}[First-order approximation]
\label{rem:first-order-taylor}
The Lie bracket~\eqref{eq:su2-commutator} captures the first-order
behaviour of the group operation: for elements $A,B$ near the identity,
\[
  e^A e^B \;=\; e^{A+B+\tfrac{1}{2}[A,B]+\cdots},
\]
so the commutator $[e_a,e_b]$ is the first-order Taylor polynomial of
the non-commutativity.  Higher-order corrections (the full
Baker--Campbell--Hausdorff series) encode the complete nonlinear group
law and reappear naturally in the universal enveloping algebra
(Section~\ref{sec:enveloping-algebras}).
\end{remark}

\begin{definition}[$\mathrm{SU}(2)$ generators of $\mathrm{Dic}_3$]
\label{def:su2-generators}
Every finite subgroup of $\mathrm{SO}(3)$ lifts, via the double cover
$\mathrm{SU}(2)\twoheadrightarrow\mathrm{SO}(3)$, to a binary subgroup of
$\mathrm{SU}(2)$.  For $S_3\cong D_3$ the lift is the binary dihedral group
$\mathrm{Dic}_3\subset\mathrm{SU}(2)$ of order~$12$, generated by
\begin{align}
  \alpha &= \exp\!\bigl(\tfrac{2\pi}{3}\,e_3\bigr)
  = \begin{pmatrix}e^{i\pi/3}&0\\[2pt]0&e^{-i\pi/3}\end{pmatrix},
  \label{eq:su2-3cycle}\\[4pt]
  \beta &= \exp(\pi\,e_1)
  = \begin{pmatrix}0&i\\[2pt]i&0\end{pmatrix},
  \label{eq:su2-2cycle}
\end{align}
with $\alpha^3=\beta^2=-I$, dihedral relation
$\beta\alpha\beta^{-1}=\alpha^{-1}$, and commutator
$[\beta,\alpha]=\alpha^{-2}$.
\end{definition}

\begin{definition}[$\mathrm{SO}(3)$ generators]
\label{def:so3-generators}
Under the adjoint map to $\mathrm{SO}(3)$ the generators of
$\mathrm{Dic}_3$ (Definition~\ref{def:su2-generators}) have images
\begin{equation}\label{eq:so3-generators}
  R_\alpha
  =\begin{pmatrix}
    -\frac{1}{2}&-\frac{\sqrt{3}}{2}&0\\[3pt]
    \frac{\sqrt{3}}{2}&-\frac{1}{2}&0\\[3pt]
    0&0&1
  \end{pmatrix},\qquad
  R_\beta
  =\begin{pmatrix}
    1&0&0\\[2pt]
    0&-1&0\\[2pt]
    0&0&-1
  \end{pmatrix}:
\end{equation}
$R_\alpha$ is a rotation by $2\pi/3$ about the $e_3$-axis (3-cycle) and
$R_\beta$ is a $\pi$-rotation about~$e_1$ (2-cycle).
\end{definition}



\section{Braid involutions and holonomy}
\label{sec:braid-involutions}

\begin{definition}[Braid involutions: twist and flip]
\label{def:braid-twist-flip}
Let $\mathbf{w}=(w_1,w_2,w_3,w_4)$ be a spectral resolvent tuple.
Define the two involutions
\begin{align}
  \sigma_1(\mathbf{w}) &= (w_2,w_1,w_4,w_3)
    &&\text{(twist),}
    \label{eq:braid-twist}\\[4pt]
  \sigma_2(\mathbf{w}) &= (w_4,w_3,w_2,w_1)
    &&\text{(flip).}
    \label{eq:braid-flip}
\end{align}
As $4\times 4$ permutation matrices these are
\begin{equation}\label{eq:perm-matrices}
  \sigma_1=\begin{pmatrix}
    0&1&0&0\\1&0&0&0\\0&0&0&1\\0&0&1&0
  \end{pmatrix},\quad
  \sigma_2=\begin{pmatrix}
    0&0&0&1\\0&0&1&0\\0&1&0&0\\1&0&0&0
  \end{pmatrix}.
\end{equation}
On the 4-tuple these involutions generate the Klein four-group
$V_4=\{\mathrm{id},\sigma_1,\sigma_2,\sigma_1\sigma_2\}$, with
$\sigma_1\sigma_2\colon\mathbf{w}\mapsto(w_3,w_4,w_1,w_2)$.
\end{definition}

\begin{remark}[Holonomy generators and the trefoil]
\label{rem:holonomy-trefoil}
Given a branched cover $\mathcal{S}\to\mathbb{CP}^1$ with singular fibres
at the critical values, choose two loops
$\gamma_1,\gamma_2$ in $\mathbb{CP}^1$ encircling distinct branch
points.  Parallel-transporting the spectral modes~$w_k$ around these
loops yields holonomy operators
\begin{equation}\label{eq:holonomy-generators}
  H_1=\sigma_1,\qquad H_2=\sigma_2,
\end{equation}
where $\sigma_1,\sigma_2$ are the braid
generators~\eqref{eq:braid-twist},~\eqref{eq:braid-flip}.  In the
$\mathrm{SU}(2)$ embedding (Definition~\ref{def:su2-generators}) these
become
\[
  H_1=\exp(T\theta\,e_1),\qquad H_2=\exp(T\theta\,e_2),
\]
with $T$ the geometric torsion scalar and $\theta$ a branch angle, so
the holonomy group is a representation of $B_3$ inside
$\mathrm{SU}(2)$.

The fundamental group of the trefoil-knot complement admits the
presentation $\langle x,y\mid xyx=yxy\rangle$, which is exactly
$B_3$~\eqref{eq:braid-presentation}.  The assignment
$x\mapsto H_1$, $y\mapsto H_2$ is therefore an isomorphism
$\pi_1(S^3\setminus K_{3,2})\;\xrightarrow{\;\sim\;}\;
\langle H_1,H_2\rangle$: the monodromy group of the branched cover
\emph{is} the trefoil-knot group.

The composite holonomy
$H_1 H_2\colon\mathbf{w}\mapsto(w_3,w_4,w_1,w_2)$ is the
$3$-cycle~$\alpha$~\eqref{eq:su2-3cycle} and satisfies
$(H_1 H_2)^3=\mathrm{id}$ on $\mathrm{SO}(3)$.
\end{remark}



\section{The trefoil knot}
\label{sec:trefoil-knot}

\begin{remark}[Trefoil knot group]
\label{rem:trefoil-knot-group}
The trefoil knot $K_{3,2}$ is the simplest non-trivial torus knot.  Its
complement has fundamental group
\[
  \pi_1(S^3\setminus K_{3,2})
  \;=\;\langle x,y\mid xyx=yxy\rangle
  \;\cong\;B_3,
\]
so the trefoil-knot group is isomorphic to the braid group on three
strands (Definition~\ref{def:braid-group}).  This relation $xyx=yxy$ is
the \emph{braid relation}: it encodes the fact that the three-strand
braid may be closed into the trefoil by identifying the top and bottom of
the braid.
\end{remark}

\begin{example}[Trefoil braid relation and operator algebra]
\label{ex:trefoil-operator-algebra}
In the $\mathrm{SU}(2)$ representation
(Definition~\ref{def:su2-generators}), the trefoil braid relation
$xyx=yxy$ becomes the operator identity
\[
  \alpha\,\beta\,\alpha
  \;=\;\beta\,\alpha\,\beta
\]
for the generators $\alpha$~\eqref{eq:su2-3cycle} and
$\beta$~\eqref{eq:su2-2cycle}.  Explicitly, using
\[
  \alpha=\begin{pmatrix}e^{i\pi/3}&0\\0&e^{-i\pi/3}\end{pmatrix},
  \qquad
  \beta=\begin{pmatrix}0&i\\i&0\end{pmatrix},
\]
both products evaluate to
\[
  \alpha\beta\alpha
  = \begin{pmatrix}0&ie^{2i\pi/3}\\ie^{-2i\pi/3}&0\end{pmatrix}
  = \beta\alpha\beta.
\]
The eigenvalues of the composite $\alpha\beta$ govern the phases
acquired when quarks exchange colour charges along the trefoil path.
\end{example}



\section{Knot fundamental groups}
\label{sec:knot-fundamental-groups}

\begin{remark}[Fundamental group of the knot complement]
\label{rem:knot-complement-group}
For a knot $K\subset S^3$ (or, more generally,
$K\subset\mathbb{CP}^1$), the fundamental group
$\pi_1(\mathbb{CP}^1\setminus K)$ encodes the topology of the
complement.  A representation
\[
  \rho\colon\pi_1(\mathbb{CP}^1\setminus K)
  \;\longrightarrow\;\mathrm{PSL}(2,\mathbb{C})
\]
maps each meridian generator to a M\"obius transformation of the
Riemann sphere.  The image of $\rho$ characterises the geometric
structure on the knot complement.
\end{remark}

\begin{remark}[Non-trivial representations of knot groups
  {\normalfont(Kronheimer--Mrowka)}]
\label{rem:kronheimer-mrowka}
Kronheimer and Mrowka proved that the fundamental group of every
non-trivial knot admits a non-trivial representation into
$\mathrm{SU}(2)$, and hence into
$\mathrm{PSL}(2,\mathbb{C})$~\cite{KronheimerMrowka2004}.  In the
QOU context, this guarantees that every non-trivial knot type
contributing to the fibre bundle carries a non-abelian gauge holonomy.
The restriction to traceless elements of $\mathrm{SU}(2)$ aligns with
the category-theoretic setup of Section~\ref{sec:enveloping-algebras},
where the representation factors through the adjoint form.
\end{remark}

\begin{remark}[Reducible representations and fixed points]
\label{rem:reducible-reps}
When a representation
$\rho\colon\pi_1(\mathbb{CP}^1\setminus K)\to\mathrm{PSL}(2,\mathbb{C})$
is reducible, all generators fix the same point on $\mathbb{CP}^1$ ---
they are simultaneously conjugate to upper-triangular matrices, hence
all rotations around a common axis.  In the category-theoretic
formulation this corresponds to the abelianisation of the knot group.
Further generalisations via the Poincar\'e--Birkhoff--Witt theorem and
unital Gr\"obner bases (Section~\ref{sec:enveloping-algebras}) describe
the passage from the reducible to the irreducible sector.
\end{remark}



\section{Universal enveloping algebras and quantum groups}
\label{sec:enveloping-algebras}

\begin{remark}[Universal enveloping algebra]
\label{rem:enveloping-algebra}
The universal enveloping algebra
$\mathcal{U}(\mathfrak{su}(2))$ is the associative algebra generated by
$e_1,e_2,e_3$ (Definition~\ref{def:su2-basis}) subject to the
commutation relations~\eqref{eq:su2-commutator}.  It provides the
algebraic scaffolding in which every finite-dimensional representation of
$\mathfrak{su}(2)$ lives and is the classical limit ($q\to 1$) of the
quantum group $\mathcal{U}_q(\mathfrak{su}(2))$.
\end{remark}

\begin{definition}[Quantum group $\mathcal{U}_q(\mathfrak{su}(2))$]
\label{def:quantum-su2}
Let $q\in\mathbb{C}^{*}$ with $0<\lvert q\rvert\le 1$.  The
\emph{quantum group} $\mathcal{U}_q(\mathfrak{su}(2))$ is the
associative $\mathbb{C}$-algebra generated by $E,F,K,K^{-1}$ subject to
\begin{align}
  KK^{-1}&=K^{-1}K=1,
    \label{eq:quantum-KK}\\[3pt]
  KEK^{-1}&=q^2\,E,\quad KFK^{-1}=q^{-2}\,F,
    \label{eq:quantum-KE}\\[3pt]
  [E,F]&=\frac{K-K^{-1}}{q-q^{-1}}.
    \label{eq:quantum-EF}
\end{align}
In the classical limit $q\to 1$, the generator $K$ reduces to
$\exp(h\,e_3)$ and the relations recover~\eqref{eq:su2-commutator}.
\end{definition}

\begin{remark}[$R$-matrix and braiding]
\label{rem:r-matrix-braiding}
The $R$-matrix of $\mathcal{U}_q(\mathfrak{su}(2))$
(Definition~\ref{def:quantum-su2}) is the universal element
\[
  \mathcal{R}\;\in\;
  \mathcal{U}_q(\mathfrak{su}(2))\;\hat{\otimes}\;
  \mathcal{U}_q(\mathfrak{su}(2))
\]
satisfying the Yang--Baxter equation
$\mathcal{R}_{12}\mathcal{R}_{13}\mathcal{R}_{23}
=\mathcal{R}_{23}\mathcal{R}_{13}\mathcal{R}_{12}$.
Its action on a tensor product of representations gives the braid
operator $\check{R}=P\circ\mathcal{R}$ (where $P$ is the swap),
furnishing a $q$-deformed braid group representation that replaces
the classical twist/flip involutions
(Definition~\ref{def:braid-twist-flip}).
\end{remark}



\section{Varying Euler characteristic}
\label{sec:varying-euler}





\subsection{The genus-$2$ case}
\label{subsec:genus-2}





\section{Descent from $E_1$ to $E_0$ and the mirror bundle $E_{-1}$}
\label{sec:descent-mirror}

\begin{remark}[$E_{-1}$ as skew-mirror bundle]
\label{rem:e-minus-one}
In the descent from $E_1$ (full torsion bundle) to $E_0$ (trivial
bundle), the self-duality of spacetime and chirality admits a flip: the
orientation of the chiral fibre can reverse relative to the spacetime
base.  This produces a bundle $E_{-1}$ that acts as a skew or mirror
image of~$E_1$.  From the perspective of ascending the bundle hierarchy,
$E_{-1}$ appears to be matter; from the perspective of descending, it
appears to be antimatter.  The distinction is purely one of orientation
in the descent tower.
\end{remark}
 \chapter{Lifting of Quantum Torsion}
\label{chap:torsion}



\begin{proposition}[Jet Prolongation]
\label{prop:jet-prolongation}
Let $(M, G, P, E)$ be a $\mathbb{C}$ quantum observable universe
(Example~\ref{def:c-quantum-observable-universe}) with instanton connection
$A$ on $P$ and metric $g$ on $E$.  Write $\widetilde{M} = E$ for the total
space of $E$, viewed as a manifold, and define $P'$ as the frame bundle and
$E'$ as the bundle of orthonormal frames of the first jet bundle $J^1 E$.
Then $A$ and $g$ uniquely determine a connection $A'$ on $P'$ such that
$(\widetilde{M}, G, P', E')$ is a $\mathbb{C}$ quantum observable universe.
\end{proposition}

\begin{proof}
The proof proceeds in the rigid monoidal category underlying the quantum
observable universe.  Let $A$ denote the Frobenius-Hopf object
representing the gauge-bundle data of $(M,G,P,E)$.

\emph{Step 1: Categorical jet object.}\;
The first jet bundle $J^1 E$ is represented categorically by the object
$A \otimes \Omega^1(A)$, where
$\Omega^1(A) = \ker(m\colon A\otimes A \to A)$
is the kernel of the multiplication (the bimodule of K\"ahler
differentials, i.e.\ the universal first-order differential calculus
over~$A$).
Because $A$ is Frobenius, the non-degenerate pairing
$\varepsilon \circ m \colon A \otimes A \to \mathbf{1}$ exhibits
$\Omega^1(A)$ as a direct summand of $A \otimes A$, hence
$J^1(A) := A \otimes \Omega^1(A)$ is again an object of the ambient
category.

\emph{Step 2: Inherited Frobenius-Hopf structure.}\;
The Frobenius relation
$(\mathrm{Id} \otimes m) \circ (\Delta \otimes \mathrm{Id})
 = \Delta \circ m
 = (m \otimes \mathrm{Id}) \circ (\mathrm{Id} \otimes \Delta)$
ensures that the comultiplication $\Delta\colon A \to A \otimes A$ restricts
to a well-defined comultiplication
$\Delta'\colon J^1(A) \to J^1(A) \otimes J^1(A)$, and the counit
$\varepsilon\colon A \to \mathbf{1}$ induces a counit $\varepsilon'$ on
$J^1(A)$.  The Hopf antipode $S$ extends to $J^1(A)$ by functoriality of the
K\"ahler differential construction.  Thus $J^1(A)$ is a Frobenius-Hopf object.

\emph{Step 3: Non-degeneracy and vanishing.}\;
The non-degeneracy of the Frobenius form on $A$ (the pairing
$\varepsilon \circ m$ is a perfect pairing) implies non-degeneracy of the
induced form on $J^1(A)$, since $\Omega^1(A)$ is a direct summand.
For the vanishing condition, $\mathrm{Hom}(\mathbf{1}, \mathrm{Ad}(A)) = 0$
implies $\mathrm{Hom}(\mathbf{1}, \mathrm{Ad}(J^1(A))) = 0$: the adjoint
representation of $J^1(A)$ contains $\mathrm{Ad}(A)$ as a subobject, and any
invariant in $\mathrm{Ad}(J^1(A))$ would project to a non-zero invariant in
$\mathrm{Ad}(A)$, contradicting the hypothesis.

\emph{Step 4: Geometric realisation.}\;
Applying the fibre functor $\omega\colon \mathcal{C} \to \mathrm{Vect}_k$
(where $\mathcal{C}$ is the rigid monoidal category of the quantum
observable universe)
recovers the geometric data: $\omega(J^1(A))$ is the space of sections of the
first jet bundle $J^1 E$, the frame bundle $P'$ and orthonormal frame bundle
$E'$ are the geometric counterparts of the Frobenius-Hopf structure on
$J^1(A)$, and the connection $A'$ is the image under $\omega$ of the
comultiplication $\Delta'$.  Its uniqueness follows from the irreducibility
of $A'$, which is guaranteed by the inherited vanishing condition of
Step~3.
\end{proof}

\begin{remark}[Prolongation from the Frobenius-Hopf structure]
\label{rem:prolongation-categorical}
The proof above shows that jet prolongation is an \emph{intrinsic}
consequence of the Frobenius-Hopf structure in a rigid monoidal category:
no additional geometric input beyond the categorical data
$(\mathcal{C}, A, \omega)$ is required.  Concretely:
\begin{enumerate}
  \item The Frobenius structure provides the K\"ahler differentials
        $\Omega^1(A)$ as a direct summand, from which the jet object
        $J^1(A)$ is built.
  \item The Frobenius relation ensures that the comultiplication and
        multiplication are compatible, so the prolonged object inherits
        the full Hopf algebra structure.
  \item The non-degeneracy and vanishing conditions propagate from $A$ to
        $J^1(A)$ by functoriality.
\end{enumerate}
In this sense, the recursive brane tower
(Definition~\ref{def:recursive-brane-tower}) is generated entirely by
iterating the categorical jet construction.  The tower, together with its
associated descent (Subsection~\ref{ssec:descent}), is a formal consequence
of the axioms of a quantum observable universe and does not require any
case-by-case verification.
\end{remark}



\begin{definition}[Recursive Brane Tower]
\label{def:recursive-brane-tower}
Let $(M, G, P, E)$ be a $\mathbb{C}$ quantum observable universe
(Example~\ref{def:c-quantum-observable-universe}).  Set $M_0 = M$ and
$E_1 = E$.  \emph{Inductively}, define $E_{n+1}$ as the jet prolongation
(Proposition~\ref{prop:jet-prolongation}) of $E_n$.

The \emph{brane relation} is the following recursive structure:
\begin{enumerate}
  \item Each total space $E_n$ is treated as the base manifold $M_{n+1}$ for
        the next level; $E_n$ is the \emph{brane} of $E_{n+1}$.
  \item At each level, $A_n \in \mathcal{A}(E_{n-1})$ is the $n$-th level
        instanton connection satisfying the Hermitian Yang--Mills equations on
        $E_n$.
\end{enumerate}
\end{definition}

\begin{example}[First two levels]
\label{ex:brane-tower-levels}
Starting from a $\mathbb{C}$ quantum observable universe $(M, G, P, E_1)$, the
first jet prolongation yields $E_2$ with base manifold $M_1 = E_1$.  The
connection $A_1$ on $E_1$ satisfies the Hermitian Yang--Mills equations, and
$A_2$ on $E_2$ is uniquely determined by
Proposition~\ref{prop:jet-prolongation}.
\end{example}



\begin{definition}[Quantum Characteristic Classes]
\label{def:quantum-characteristic-classes}
Let $\{E_n\}_{n \geq 1}$ be the recursive brane tower
(Definition~\ref{def:recursive-brane-tower}).  The \emph{quantum
characteristic classes} are the cohomology classes
\[
  \kappa_n \;\in\; H^*\!\bigl(E_{n+1},\, \mathbb{Z}\bigr)
\]
that label the monodromy fixed points of the holonomy representation as one
transitions from $E_n$ to $E_{n+1}$.
\end{definition}



\begin{remark}[Betti numbers of the brane tower]
\label{rem:betti-numbers-brane}
Let $b_k(E_n) = \dim H^k(E_n, \mathbb{R})$ denote the $k$-th Betti number of
$E_n$.  Since $E_{n+1}$ is the bundle of orthonormal frames of $J^1 E_n$, the
Betti numbers of $E_{n+1}$ are determined by those of $E_n$ via the Leray--Serre
spectral sequence applied to the fibration
$E_{n+1} \to E_n$.  In particular, the quantum characteristic classes
(Definition~\ref{def:quantum-characteristic-classes}) obstruct the collapse of
this spectral sequence and control the growth of $b_k(E_{n+1})$ in terms of
$b_k(E_n)$.
\end{remark}



\begin{remark}[$q$-Harmonic basis of the brane tower]
\label{rem:q-harmonic-basis-brane}
For each level $n$, the $q$-harmonic $1$-forms $\alpha_{n,i}$ belong to the
space $\mathcal{H}^1(E_n,\, \mathrm{ad}\, E_n)$ of harmonic $1$-forms on
$E_n$ with values in the adjoint bundle $\mathrm{ad}\, E_n$.  These are
characterised by the self-duality condition relative to the instanton
connection $A_n$
(Definition~\ref{def:recursive-brane-tower}):
\[
  \alpha_{n,i} \in \mathcal{H}^1
  \iff
  d_{A_n}\,\alpha = 0
  \quad\text{and}\quad
  d_{A_n}^*\,\alpha = 0,
\]
where the Hodge-star operator is induced by the Hermitian metric that the
Levi form $L$ (Definition~\ref{def:substrate-parameter}) defines on the
horizontal distribution.
The index $i$ runs over $1 \leq i \leq b_1(E_n)$
(Remark~\ref{rem:betti-numbers-brane}).

At the base level $n = 0$, with Euler characteristic $\chi(M) = 3$, there are
three primary basis vectors $\{\,\phi_1,\, \phi_2,\, \phi_3\,\}$.  At
level $n = 1$ the basis lifts: the new basis includes the original modes
together with their $J^1$ prolongations
(Proposition~\ref{prop:jet-prolongation}), which are the first-order
resolutions of the base flows.
\end{remark}



\begin{remark}[Weight-space decomposition of the harmonic basis]
\label{rem:weight-space-decomposition}
The adjoint bundle $\mathrm{ad}\, E_n$ has fibre $\mathfrak{su}(2)$.  Via the
weight-space decomposition of $\mathrm{SU}(2)$, the harmonic basis of
$\mathcal{H}^1(E_n,\, \mathrm{ad}\, E_n)$ splits into two families.

\emph{Longitudinal modes.}
These forms constitute the $\mathrm{U}(1)$ backbone of the decomposition.
They take values in the Cartan subalgebra of $\mathfrak{su}(2)$ and have the
geometric form $\omega^0 \otimes H$, where $H$ is the diagonal generator.
They represent steady-state laminar flows aligned with the $\mathrm{U}(1)$
subgroup.  If the bundle splits at a monodromy fixed point, these are the only
modes that survive; they correspond to the magnetic flux of the stable
vortices.

\emph{Ladder modes.}
These are the raising and lowering vortices of the decomposition.  They take
values in the off-diagonal generators $E_\pm$ of $\mathfrak{su}(2)$ and have
the geometric form $\omega^\pm \otimes E_\pm$.  They describe the entangling
flows that twist the $\mathbb{CP}^1$ fibre as one moves along the brane.
These modes are responsible for the non-commutativity in the
$\mathrm{SU}_q(2)$ algebra (Definition~\ref{def:direct-limit-algebra}).

In the $1$-jet bundle $J^1 E_n$
(Remark~\ref{rem:j1e-local-states}), each basis element $\alpha$ is a pair
consisting of a \emph{value} (the direction of the $\mathrm{SU}(2)$ flow in
the fibre) and a \emph{momentum} (the gradient of that flow relative to the
complex structure of the brane).
\end{remark}



\begin{remark}[The algebraic filter]
\label{rem:algebraic-filter}
Let $\alpha$ be a harmonic $1$-form in the basis of
Remark~\ref{rem:q-harmonic-basis-brane}.  Its \emph{signature} is determined
by its commutator with the curvature $F_A$
(Definition~\ref{def:substrate-parameter}).

\emph{Fermionic type (matter).}
$[F_A,\, \alpha] \neq 0$.  These forms are entangled with the curvature: they
see the $\mathrm{SO}(3)$ holonomy and the $-6$ index fixed points.  Their
geodesics are trapped or spiralling around the three fundamental vortex
filaments.

\emph{Bosonic type (radiation).}
$[F_A,\, \alpha] = 0$.  These forms commute with the curvature.  They are the
longitudinal modes (Remark~\ref{rem:weight-space-decomposition}) that flow
through the bulk without being deflected by the monodromy fixed points.  Their
geodesics are straight relative to the gauge-fluid flow.
\end{remark}



\begin{remark}[Distributional connection at the limit]
\label{rem:distributional-connection}
At the limit $n \to \infty$ one no longer looks at a single connection $A$,
but at a \emph{distributional connection} represented by the series
\[
  A_\infty
  \;=\;
  \sum_{n=0}^{\infty}\;\sum_{i=1}^{b_1(E_n)}
  c_{n,i}\,\alpha_{n,i}\,,
\]
where $c_{n,i} \in \mathbb{C}$ and $\{\alpha_{n,i}\}$ is the $q$-harmonic
basis of Remark~\ref{rem:q-harmonic-basis-brane}.  This turns the gauge field
into a harmonic expansion across the infinite jet space $J^\infty E$
(Remark~\ref{rem:jet-shell}).
\end{remark}



\begin{definition}[$\mathbb{C}$ Quantum Observable Universe as Direct Limit]
\label{def:direct-limit-algebra}
Let $\{E_n\}_{n \geq 1}$ be the recursive brane tower
(Definition~\ref{def:recursive-brane-tower}) with substrate parameters $q_n$
(Definition~\ref{def:substrate-parameter}) at each level.  The
\emph{$\mathbb{C}$ quantum observable universe} is the direct limit algebra
\[
  \mathcal{A}_\infty
  \;=\;
  \varinjlim\,
  \Gamma\!\bigl(J^1 E_n,\, \mathrm{SU}_{q_n}(2)\bigr).
\]
\end{definition}

\begin{remark}[Direct limit algebra and the Lusztig integral form]
\label{rem:direct-limit-lusztig}
The algebra $\mathcal{A}_\infty$
(Definition~\ref{def:direct-limit-algebra}) is explicitly realised as the
Lusztig integral form of the quantised enveloping algebra
$\mathcal{U}_q(\mathfrak{su}(2))$~\cite{lusztig1993}.  At each level~$n$
the fibre algebra $\mathrm{SU}_{q_n}(2)$ is obtained by specialising the
$\mathbb{Z}[q,q^{-1}]$-form
$\mathcal{U}^{\mathbb{Z}}_q(\mathfrak{su}(2))$
(Definition~\ref{def:substrate-parameter}) along
$q \mapsto q_n \in k^{\times}$, and the direct limit inherits the divided
power basis $(E^{(m)}, F^{(m)}, K^{\pm 1})$ from the integral form.  This
ensures that the quaternionic jet embedding
(Definition~\ref{def:quaternionic-spectral-resolvents}) remains robust
under the arithmetic descent to~$M_0$, with the $q$-deformed scaling
function $\exp_q(r) - 1$ (equation~\eqref{eq:q-exp}) regulating the
transition between the divided power basis of the space-like and time-like
resolvents.
\end{remark}

\begin{remark}[Jet Shell of the universe]
\label{rem:jet-shell}
In the limit $n \to \infty$, the infinite jet bundle
$J^\infty E = \varprojlim\, J^k E$ becomes the \emph{jet shell} of the
universe.  All finite-level brane data are encoded in this pro-finite
structure, and the algebra $\mathcal{A}_\infty$
(Definition~\ref{def:direct-limit-algebra}) acts naturally on sections of the
jet shell.
\end{remark}



\begin{proposition}[Existence and Uniqueness of the Vacuum State]
\label{prop:vacuum-state}
Let $\mathcal{A}_\infty$ be the direct limit algebra
(Definition~\ref{def:direct-limit-algebra}) equipped with the global
$\mathrm{SU}(2)$ action and the recursive $q$-deformations.  Then there exists
a unique state $\omega \colon \mathcal{A}_\infty \to \mathbb{C}$ that is
$\mathrm{SU}(2)$-invariant:
\[
  \omega(g \cdot a) = \omega(a)
  \quad\text{for all } g \in \mathrm{SU}(2),\;
  a \in \mathcal{A}_\infty.
\]
\end{proposition}

\begin{proof}
\end{proof}



\begin{remark}[Mass in the brane tower]
\label{rem:mass-geometric-drag}
In the recursive brane tower (Definition~\ref{def:recursive-brane-tower}),
mass is not an intrinsic property of a point-particle but a measure of the
resistance encountered by a stable vortex as it propagates through the
$1$-jet bundle $J^1 E_n$ (Proposition~\ref{prop:jet-prolongation}).  The
substrate parameter~$q$ (Definition~\ref{def:substrate-parameter}) varies from
fibre to fibre; translation of the vortex along the base therefore requires a
corresponding variation of the Frobenius--Hopf structure on each fibre, and
the mass quantifies this variation.
\end{remark}

\begin{remark}[Non-commutativity and geometric drag]
\label{rem:mass-noncommutative-drag}
When the substrate parameter satisfies $q(p) \neq 1$ at every point of the recursive brane tower, the comultiplication on the Frobenius--Hopf object~$A$
(Definition~\ref{def:quantum-universe}) is non-cocommutative, so the
algebra $\mathcal{A}_\infty$ (Definition~\ref{def:direct-limit-algebra})
does not reduce to a commutative direct limit.  Moving a vortex through
$J^1 E_n$ then requires a continuous rearrangement of the
$\mathrm{SU}_q(2)$ fibre data, which is the source of the inertial
resistance described in Remark~\ref{rem:mass-geometric-drag}.
\end{remark}

\begin{proposition}[Stability of the Euler characteristic]
\label{prop:euler-char-stability}
Let $(M, G, P, E)$ be a $\mathbb{C}$ quantum observable universe and let
$\{E_n\}_{n \geq 1}$ be the recursive brane tower
(Definition~\ref{def:recursive-brane-tower}).  Then the Euler characteristic
$\chi(E)$ is stable under descent and lifting: for each level $n$,
\[
  \chi(E_{n+1}) \;=\; \chi(E_n).
\]
In particular, if the base level satisfies $\chi(E_1) = c$, then
$\chi(E_n) = c$ for all $n \geq 1$.
\end{proposition}

\begin{proof}
\end{proof}

\begin{definition}[Mass Identity]
\label{def:mass-identity}
Let $(M, G, P, E)$ be a $\mathbb{C}$ quantum observable universe with substrate
parameter $q$ (Definition~\ref{def:substrate-parameter}), connection $A$, Levi
form $L$, and self-dual curvature $F^+$.  Let $\gamma$ be a closed cycle
encircling the vortex filament.  The \emph{inertial mass} $m$ of a stable
vortex is
\[
  m \;\propto\; \oint_\gamma \lVert \nabla_A\, q \rVert
               \cdot \lVert L \wedge F^+ \rVert\,.
\]
\end{definition}



\subsection{Descent}
\label{ssec:descent}

\begin{definition}[Descent]
\label{def:descent}
Let $\{E_n\}_{n \geq 1}$ be the recursive brane tower
(Definition~\ref{def:recursive-brane-tower}).  The \emph{descent} is the
reverse traversal of the tower: starting from level~$n$ one passes to level
$n-1$ by restricting the connection $A_n$ and the bundle $E_n$ to the base
manifold $M_{n-1}$.  The descent terminates when the base manifold $M_k$
reaches $\dim M_k = 0$, at which point $M_k$ is a discrete set of points.
\end{definition}

\begin{remark}[Reduction at dimension~$1$]
\label{rem:descent-reduction}
When $\dim M_k = 1$ the geometric data of the $\mathbb{C}$ quantum observable
universe collapse as follows.
\begin{enumerate}
  \item \emph{Connection.}  The $\mathrm{SU}(2)$ connection $A_k$ restricts to
        a single dimension; its holonomy reduces to a $\mathrm{U}(1)$ phase
        around $S^1$.
  \item \emph{Bundle.}  The principal $\mathrm{SU}(2)$ bundle $E$ descends to a
        bundle over the $1$-dimensional base whose fibre dynamics are locked
        into a single periodic flow.
  \item \emph{Reeb direction.}  The inverse of the Reeb vector field
        (Definition~\ref{def:q-reeb-vector-field}) becomes the only direction
        of motion; the horizontal distribution is trivial and no bosonic modes
        (Remark~\ref{rem:algebraic-filter}) survive.
\end{enumerate}
The descent need not terminate here; a further step carries $M_k$ to a
zero-dimensional base (Remark~\ref{rem:descent-reduction-zero}).
\end{remark}

\begin{remark}[Physical implications of a $1$-dimensional base]
\label{rem:descent-physical}
At dimension~$1$ the physical content simplifies drastically.
\begin{enumerate}
  \item \emph{Primordial vortex.}  The vortex filaments of the mass identity
        (Definition~\ref{def:mass-identity}) are no longer embedded in a larger
        space; the entire manifold \emph{is} the filament.
  \item \emph{Mass as pure drag.}  The mass identity reduces to an integral
        over the whole of $M_k$:
        \[
          m \;\propto\; \oint_{M_k} \lVert \nabla_A\, q \rVert
                       \cdot \lVert L \wedge F^+ \rVert\,,
        \]
        so that mass represents the fundamental geometric resistance of the
        $1$-dimensional flow against the $\mathrm{SU}(2)$ rotation.
  \item \emph{Lepton generations.}  Since the number of lepton generations is
        bounded by the Euler characteristic and vertical degree
        (Remark~\ref{rem:betti-numbers-brane}), a $1$-dimensional base supports
        only the most fundamental fermionic state; there is no room for
        higher-degree vertical harmonic forms.
\end{enumerate}
\end{remark}

\begin{remark}[Reduction at dimension~$0$]
\label{rem:descent-reduction-zero}
When $\dim M_k = 0$ the base manifold is a discrete set of points and
the geometric data of the $\mathbb{C}$ quantum observable universe
reduce to purely algebraic invariants at each point $p\in M_k$.
\begin{enumerate}
  \item \emph{Connection.}  The $\mathrm{SU}(2)$ connection $A_k$ has no
        directions along which to parallel-transport; its curvature
        $F_A$ vanishes identically and the holonomy is trivial.
  \item \emph{Bundle.}  The principal $\mathrm{SU}(2)$ bundle $E$
        restricts to its fibre $E_p$ over the point~$p$; the bundle
        is determined entirely by this single fibre together with
        the group action.
  \item \emph{Substrate parameter.}  The substrate parameter
        $q(p) = 1 + \lVert L\wedge F_A^{+}\rVert(p)$
        (Remark~\ref{rem:c-substrate-via-levi}) assigns to each point a
        definite quantum constant $q\in[1,\infty)$, encoding the full
        content of the quantum observable universe at that point.
  \item \emph{Reeb direction.}  There is no tangent space, so the Reeb
        vector field (Definition~\ref{def:q-reeb-vector-field}) and the
        horizontal distribution are absent; only the internal
        $\mathrm{SU}(2)$ fibre data persist.
\end{enumerate}
In this terminal stage every observable is determined by the fibre
$E_p$ and the value~$q(p)$.  When $M_k$ consists of several points the
descent data decompose into independent summands, one per point; when
$M_k$ is connected---i.e.\ a single point---the entire quantum
observable universe is captured by the pair $(E_p, q(p))$.
\end{remark}

\begin{remark}[Physical implications of a $0$-dimensional base]
\label{rem:descent-physical-zero}
At dimension~$0$ all spatial extent has been removed and the physical
content is purely algebraic.
\begin{enumerate}
  \item \emph{Primordial state.}  The vortex filaments, flows and spatial
        modes of higher-dimensional strata are absent; what remains is
        the initial datum of the fermionic content encoded in the
        fibre~$E_p$.
  \item \emph{Mass.}  With no base manifold over which to integrate, the
        mass identity (Definition~\ref{def:mass-identity}) collapses to a
        single evaluation: $m \propto \lVert L \wedge F^+\rVert(p)$.
        Mass at dimension~$0$ is an instantaneous, pointwise quantity
        rather than a distributed one.
  \item \emph{Lepton generations.}  All vertical harmonic forms
        (Remark~\ref{rem:betti-numbers-brane}) are concentrated at the
        point; the Euler characteristic $\chi(E_p)$ retains the count of
        fermionic degrees of freedom without any spatial multiplicity.
\end{enumerate}
\end{remark}

\begin{remark}[Fixed-point mechanism: ascent and descent]
\label{rem:descent-fixed-point}
Each step down the brane tower removes a space-like dimension through the
$\mathrm{SU}(2)$ rotation and its associated fixed-point set.
\begin{enumerate}
  \item \emph{Ascent (generation).}  Moving from $M_k$ ($\dim = 0$) upward
        lifts the connection $A$ to higher categorical structures via the jet
        prolongation (Proposition~\ref{prop:jet-prolongation}).
  \item \emph{Descent (filtering).}  Moving back down acts as an algebraic
        filter (Remark~\ref{rem:algebraic-filter}).  The ladder modes
        (Remark~\ref{rem:weight-space-decomposition}) are stripped away:
        at dimension~$1$ only the longitudinal
        $\mathrm{U}(1)$ backbone remains
        (Remark~\ref{rem:descent-reduction}), and at dimension~$0$ even
        that is collapsed to a single fibre datum
        (Remark~\ref{rem:descent-reduction-zero}).
\end{enumerate}
\end{remark}

\begin{conjecture}[Topological Decomposition of Mass]
Let $(M, \theta, G, P, E)$ be a $\mathbb{C}$ quantum observable universe with substrate parameter $q$, and let $\gamma$ be a stable vortex filament. Suppose the topology of $\gamma$ is isomorphic to a knot generated by $N$ consecutive topological blowups applied to an unknot core $\gamma_0$. The continuous geometric mass of the vortex along $\gamma$ admits the discrete topological decomposition:
\begin{equation}
m(\gamma) \propto \oint_{\gamma_0} ||\nabla_A q|| \cdot ||L \wedge F^+|| + \sum_{i=1}^{N} \lambda_i \left| \Phi_q(z_i) \right|
\end{equation}
where the line integral evaluates the baseline geometric drag of the trivial unknot core, $\Phi_q(z_i)$ is the quantum dilogarithm evaluated at the spectral coordinate $z_i$ of the $i$-th exceptional divisor $E_i$, and $\lambda_i$ are local coupling weights.
\end{conjecture}

\begin{corollary}[Quantum Dilogarithm as Localized Torsion Volume]
The localized topological mass contribution $\Phi_q(z_i)$ at the $i$-th exceptional divisor $E_i$ is exactly the geometric integration of the 1-dimensional $q$-deformed expansion flow over the resolved defect. Integrating the spectral uncertainty flow $ln_q$ along the radial coordinate of the Briggs-Klein affine chart yields the quantum dilogarithm:
\begin{equation}
\Phi_q(z_i) = \int_0^{z_i} \frac{ln_q(1+t)}{t} dt = \sum_{n=1}^\infty (-1)^{n+1} \frac{z_i^n}{n [n]_q}
\end{equation}
where $ln_q(1+t) = \sum_{n=1}^\infty (-1)^{n+1} \frac{t^n}{[n]_q}$ governs the continuous Reeb flow[cite: 8]. The integration introduces the factor of $n$ in the denominator, elevating the 1-dimensional path propagation to a geometric volume measure trapping the topological charge.
\end{corollary}


\chapter{Bring's \DIFdelbegin \DIFdel{Surface}\DIFdelend \DIFaddbegin \DIFadd{Hierarchy}\DIFaddend }
\DIFaddbegin \DIFadd{\label{chap:brings-hierarchy}
}\DIFaddend \label{chap:brings-surface}

\DIFaddbegin \DIFadd{The \emph{Bring hierarchy} is the tower of algebraic curves
}\[
  \DIFadd{B_0 \;\hookrightarrow\; B_1 \;\hookrightarrow\; B_2 \;\hookrightarrow\; B_3,
}\]
\DIFadd{ordered by increasing genus $0, 0, 2, 4$ and connected by degree-$2$
branched coverings.  $B_3$ is Bring's curve proper; $B_2$ is its
hyperelliptic quotient; $B_1 = \mathbb{P}^1$ is the spacetime base;
and $B_0$ is the underlying ground field (a point in the Descartes
descent).  The hierarchy encodes a filtration of the QOU by topological
complexity and provides the geometric scaffolding for the particle
ansatz of Chapter~\ref{chap:path-integrals-braiding}.
}



\section{\DIFadd{General Definition and the Hierarchy}}
\DIFadd{\label{sec:hierarchy-general}
}

\begin{definition}[Bring hierarchy]
\DIFadd{\label{def:bring-hierarchy}
The \emph{Bring hierarchy} is the sequence of smooth projective curves
}\[
  \DIFadd{B_0,\quad B_1,\quad B_2,\quad B_3
}\]
\DIFadd{over $\mathbb{C}$ (or any algebraically closed field $k$ of
characteristic $\neq 5$) with genera $g(B_n) = 2(n-1)$ for $n\ge 1$
and $B_0$ a point, together with degree-$2$ holomorphic maps
}\[
  \DIFadd{B_3 \xrightarrow{\;\pi_{32}\;} B_2
  \xrightarrow{\;\pi_{21}\;} B_1
  \xrightarrow{\;\pi_{10}\;} B_0,
}\]
\DIFadd{characterised as follows.  The covering $\pi_{21}\colon B_2\to B_1$ is
the hyperelliptic involution of $B_2$; the covering
$\pi_{32}\colon B_3\to B_2$ is the restriction of the degree-$4$ map
$\pi\colon B_3\to B_1$ to the intermediate quotient $B_2$.
}\end{definition}

\begin{remark}[$B_0$: the point]
\DIFadd{\label{rem:B0-point}
$B_0 = \mathrm{Spec}(k)$ is a single $k$-rational point.  In the QOU
descent picture (Remark~\ref{rem:descent-to-zero}), passing to $B_0$
means projecting out all spatial and chiral degrees of freedom; the
resulting object retains only the substrate parameter $q \in [1,\infty)$
as data.  Topologically, $B_0$ has genus $g = -1$ if one follows the
recursion $g(B_n)=2(n-1)$; we set $g(B_0)=\varnothing$ (undefined) and
regard $B_0$ as the degenerate base of the tower.
}\end{remark}

\begin{remark}[$B_1 = \mathbb{P}^1(k)$, genus $0$]
\DIFadd{\label{rem:B1-projective-line}
$B_1 = \mathbb{P}^1(k)$ is the projective line over $k$, the unique
smooth projective curve of genus $0$ with a $k$-rational point.  In
the QOU framework it plays the role of the \emph{spacetime base}: the
coordinate $z = x_0/x_1 \in \mathbb{P}^1$ parametrises the
one-dimensional spacetime direction inherited from the Descartes
descent.  The Euler characteristic is $\chi(B_1) = 2$.
}

\DIFadd{The map $\pi_{10}\colon B_1 \to B_0$ is the structure morphism; it
records the fact that $\mathbb{P}^1$ has a unique rational point over the
ground field, namely the class $[B_1]$ itself.
}\end{remark}

\begin{remark}[$B_2$, genus $2$]
\DIFadd{\label{rem:B2-genus2}
$B_2$ is the unique (up to isomorphism over the moduli space) smooth
projective curve of genus $2$ that arises as the \emph{discriminant cover}
of $B_1$ branched at the six singular fibres of the elliptic fibration
of $B_3$ over $B_1$.  Its Euler characteristic is
$\chi(B_2) = 2 - 2\cdot 2 = -2$.
}

\DIFadd{Concretely, if $\Delta_{B_3}(z)$ denotes the discriminant of the chiral
elliptic cubic of $B_3$
(equation~\eqref{eq:weierstrass-depressed}, Remark~\ref{rem:t-duality-elliptic}),
then $B_2$ is the hyperelliptic curve
}\begin{equation}\DIFadd{\label{eq:B2-hyperelliptic}
  B_2\colon\quad w^2 = \Delta_{B_3}(z),
}\end{equation}
\DIFadd{where $\Delta_{B_3}(z) = -4\,p(z)^3 - 27\,q(z)^2$, with
$p(z) = \tfrac{1}{3}(2z^2+z+2)$ and
$q(z) = \tfrac{5}{27}(4z^3+3z^2+3z+4)$, is a degree-$6$ polynomial
with $6$ simple roots (the branch locus of $\pi_{21}$).
The automorphism group of a generic $B_2$ is $\mathbb{Z}/2\mathbb{Z}$
(the hyperelliptic involution $\iota_2\colon(z,w)\mapsto(z,-w)$);
the $B_2$ arising here inherits additional automorphisms from the
$S_5$-symmetry of $B_3$.
}\end{remark}

\begin{remark}[$B_3$: Bring's surface, genus $4$]
\DIFadd{\label{rem:B3-brings-surface}
$B_3$ is \emph{Bring's curve}, the unique smooth projective curve of
genus $g = 4$ and Euler characteristic $\chi = -6$ defined by the
simultaneous vanishing of the first three power sums on $\mathbb{P}^4$
(Definition~\ref{def:brings-surface}, Section~\ref{sec:brings-surface}).
Its automorphism group $\mathrm{Aut}(B_3) \cong S_5$ has order $120$,
the maximum possible for a genus-$4$ curve by the Hurwitz bound.
}

\DIFadd{The degree-$4$ map $\pi\colon B_3 \to B_1$ is a branched covering with
$6$ branch points at the roots of $\Delta_{B_3}(z)$; it factors as
$\pi = \pi_{21} \circ \pi_{32}$, where $\pi_{32}$ is a degree-$2$ map
to the intermediate curve $B_2$.
}\end{remark}

\begin{proposition}[Riemann--Hurwitz for the hierarchy]
\DIFadd{\label{prop:rh-hierarchy}
The three branched covers in the Bring hierarchy satisfy:
}\begin{alignat}{\DIFadd{3}}
  &\DIFadd{\pi_{32}\colon B_3\to B_2,\quad\deg=2,
  }&&\DIFadd{\quad 2\cdot 4 - 2 = 2(2\cdot 2-2)+2,
  }&&\DIFadd{\quad\text{$2$ simple ramification points,}
  \label{eq:rh-32}}\\[4pt]
  &\DIFadd{\pi_{21}\colon B_2\to B_1,\quad\deg=2,
  }&&\DIFadd{\quad 2\cdot 2 - 2 = 2(2\cdot 0-2)+6,
  }&&\DIFadd{\quad\text{$6$ simple branch points,}
  \label{eq:rh-21}}\\[4pt]
  &\DIFadd{\pi_{10}\colon B_1\to B_0,\quad\deg=1,
  }&&\DIFadd{\quad \text{unramified.}
  }&&
  \DIFadd{\label{eq:rh-10}
}\end{alignat}
\end{proposition}
\begin{proof}
\DIFadd{Equation~\eqref{eq:rh-21}: $\chi(B_2) = 2\chi(B_1) - R$ with $R=6$
ramification points gives $-2 = 4 - 6$ }\checkmark\DIFadd{.
Equation~\eqref{eq:rh-32}: $\chi(B_3) = 2\chi(B_2) - R$ with $R=2$
gives $-6 = -4 - 2$ }\checkmark\DIFadd{.  Equation~\eqref{eq:rh-10} is immediate.
}\end{proof}

\begin{remark}[Genera and Euler characteristics]
\DIFadd{\label{rem:hierarchy-genera}
The Bring hierarchy levels have the following topological invariants:
}\begin{center}
\renewcommand{\arraystretch}{1.3}
\begin{tabular}{ccccc}
\hline
\DIFadd{Level }& \DIFadd{Curve }& \DIFadd{Genus $g$ }& \DIFadd{$\chi = 2-2g$ }& \DIFadd{Automorphisms }\\
\hline
\DIFadd{$B_0$ }& \DIFadd{point }& \DIFadd{--- }& \DIFadd{--- }& \DIFadd{trivial }\\
\DIFadd{$B_1$ }& \DIFadd{$\mathbb{P}^1(k)$ }& \DIFadd{$0$ }& \DIFadd{$2$ }& \DIFadd{$\mathrm{PGL}(2,k)$ }\\
\DIFadd{$B_2$ }& \DIFadd{hyperelliptic }& \DIFadd{$2$ }& \DIFadd{$-2$ }& \DIFadd{$\mathbb{Z}/2\mathbb{Z}$ (generic) }\\
\DIFadd{$B_3$ }& \DIFadd{Bring's curve }& \DIFadd{$4$ }& \DIFadd{$-6$ }& \DIFadd{$S_5$ (order $120$) }\\
\hline
\end{tabular}
\end{center}
\end{remark}



\section{\DIFadd{Coordinate Development for $B_2$}}
\DIFadd{\label{sec:B2-coordinates}
}

\DIFadd{We develop the coordinate structure of $B_2$ in parallel with the
treatment of $B_3$ in Section~\ref{sec:brings-surface}.  The key
simplification is that $B_2$ is a hyperelliptic curve: all its
geometry is determined by the pair $(z, w)$ satisfying
equation~\eqref{eq:B2-hyperelliptic}, plus the involution
$\iota_2\colon w \mapsto -w$.
}

\begin{definition}[$B_2$ as discriminant cover]
\DIFadd{\label{def:B2-discriminant-cover}
Let $p(z) = \tfrac{1}{3}(2z^2+z+2)$ and
$q(z) = \tfrac{5}{27}(4z^3+3z^2+3z+4)$ be the
Weierstrass coefficients of the $B_3$ chiral elliptic fibration
(equation~\eqref{eq:weierstrass-depressed}).  The
\emph{discriminant curve} $B_2$ is the smooth projective model of the
affine curve
}\begin{equation}\DIFadd{\label{eq:B2-affine}
  w^2 \;=\; \Delta(z)
  \;:=\; -4\,p(z)^3 - 27\,q(z)^2
  \;\in\; \mathbb{Z}[z],
}\end{equation}
\DIFadd{a degree-$6$ polynomial.  The six simple roots $z_1,\dots,z_6$
of $\Delta(z)$ are the \emph{branch points} of $\pi_{21}\colon B_2\to B_1$
and coincide with the six singular fibres of the $B_3$ elliptic
fibration.
}\end{definition}

\begin{remark}[QOU $4$-tuple for $B_2$]
\DIFadd{\label{rem:B2-qou-tuple}
The $\mathbb{C}$ quantum observable universe
(Example~\ref{def:c-quantum-observable-universe}) for $B_2$
is the $4$-tuple $(M_{B_2}, G, P_{B_2}, E_{B_2})$ constructed as follows.
}

\medskip
\noindent\DIFadd{\textbf{1.~CR manifold $M_{B_2}$.}\;
The unit circle bundle $M_{B_2} = \{(p,\eta)\in T^*B_2 : |\eta|_{g}=1\}$
over $B_2$ is a $3$-dimensional strictly pseudoconvex CR manifold
with contact form $\theta_{B_2} = \eta\,dz|_{M_{B_2}}$.
}

\medskip
\noindent\DIFadd{\textbf{2.~Principal bundle $P_{B_2}$.}\;
The degree-$2$ covering $\pi_{21}\colon B_2\to\mathbb{P}^1$ branched
at the six roots $z_1,\dots,z_6$ lifts to a principal
$\mathrm{SU}(2)$-bundle $P_{B_2}\to M_{B_2}$.  The transition functions
are the $\mathrm{SU}(2)$ holonomy matrices around the six branch cuts;
the deck transformation $(z,w)\mapsto(z,-w)$ acts on $P_{B_2}$ by the
central element $-I\in\mathrm{SU}(2)$.
}

\medskip
\noindent\DIFadd{\textbf{3.~State bundle $E_{B_2}$.}\;
The associated state bundle $E_{B_2} = P_{B_2}\times_G\mathbb{CP}^1$
has fibre $\mathbb{CP}^1$ over each point of $M_{B_2}$.  The two sheets
of $\pi_{21}$ give the two hemispheres of each $\mathbb{CP}^1$ fibre,
glued by the monodromy involution.
}

\medskip
\noindent\DIFadd{\textbf{4.~Substrate parameter.}\;
The substrate parameter
$q(p) = 1 + |\langle L_{\theta_{B_2}}, F_{A}^+\rangle|(p)$
is maximal at the six branch points $z_i$ and tends to $1$ away from
them; the deep-quantum regime is therefore localised at
$\Delta(z)=0$.
}\end{remark}

\begin{remark}[Holomorphic differentials on $B_2$]
\DIFadd{\label{rem:B2-holomorphic-differentials}
A smooth genus-$2$ curve has a $2$-dimensional space of holomorphic
differentials $H^0(B_2, \Omega^1_{B_2})$.  For the hyperelliptic model
$w^2 = \Delta(z)$ the two basis differentials are
}\begin{equation}\DIFadd{\label{eq:B2-differentials}
  \omega_1 \;=\; \frac{dz}{w}
  \;=\; \frac{dz}{\sqrt{\Delta(z)}},
  \qquad
  \omega_2 \;=\; \frac{z\,dz}{w}
  \;=\; \frac{z\,dz}{\sqrt{\Delta(z)}}.
}\end{equation}
\DIFadd{These replace the single holomorphic differential of an elliptic curve;
the extra factor of $z$ in $\omega_2$ reflects the additional handle
(genus $2$ vs.\ genus $1$).  The period matrix
$\tau = (\oint_{A_j}\omega_i)^{-1}(\oint_{B_j}\omega_i)$ encodes the
complex structure of $B_2$; its imaginary part is positive definite
(Riemann's bilinear relations).
}\end{remark}

\begin{remark}[Deck transformation and $\mathbb{Z}/2\mathbb{Z}$ spectral action]
\DIFadd{\label{rem:B2-deck-transformation}
The hyperelliptic involution
$\iota_2\colon (z,w) \mapsto (z,-w)$
is the deck transformation of $\pi_{21}\colon B_2\to B_1$.  It acts on
the holomorphic differentials by
}\[
  \DIFadd{\iota_2^*\omega_1 = \omega_1,
  \qquad
  \iota_2^*\omega_2 = \omega_2,
}\]
\DIFadd{preserving the entire space $H^0(B_2,\Omega^1)$ (which is expected:
the deck transformation of the hyperelliptic map always preserves
the canonical bundle).  In terms of the DFT analogy with $B_3$:
whereas $B_3$ carries a $C_5$ spectral action
$w_k\mapsto\zeta_5^{-k}w_k$ on its Lagrange resolvents, $B_2$ carries a
$C_2 = \mathbb{Z}/2\mathbb{Z}$ action $w\mapsto -w$ corresponding to
the fourth root of unity $\zeta_4^2 = -1$.  The role of the DFT basis
$\{v_k\}_{k=0}^{3}$ with $v_k = \tfrac{1}{2}\sum_{j=0}^{3}i^{kj}y_j$
is played by the pair $(z, w)$ after the covering identification.
}\end{remark}

\begin{remark}[Spectral forms for $B_2$]
\DIFadd{\label{rem:B2-spectral-forms}
Parallel to the Lagrange resolvent spectral forms of $B_3$
(Remark~\ref{rem:lagrange-resolvents}), the $B_2$ model carries two
spectral constraints.  Setting $v_0 = \tfrac{1}{2}(y_0+y_1+y_2+y_3)$,
$v_1 = \tfrac{1}{2}(y_0+iy_1-y_2-iy_3)$,
$v_2 = \tfrac{1}{2}(y_0-y_1+y_2-y_3)$,
$v_3 = \tfrac{1}{2}(y_0-iy_1-y_2+iy_3)$ for homogeneous coordinates
$y_0,\dots,y_3$ on $\mathbb{P}^3$, the power-sum constraints
$\sum y_j = 0$ and $\sum y_j^2 = 0$ become:
}\begin{align}
  \DIFadd{v_0 }&\DIFadd{= 0
    }&&\DIFadd{\text{(linear projection),}
    \label{eq:B2-resolvent-linear}}\\[4pt]
  \DIFadd{2v_1 v_3 + v_2^2 }&\DIFadd{= 0
    }&&\DIFadd{\text{(quadric constraint).}
    \label{eq:B2-resolvent-quadric}
}\end{align}
\DIFadd{Together these cut out a smooth conic in $\mathbb{P}^2$,
isomorphic to $\mathbb{P}^1$, which is the spacetime base $B_1$.  The
full $B_2$ arises by adjoining the square root $w = \sqrt{\Delta(z)}$,
passing to the branched double cover of this $\mathbb{P}^1$.
The $C_4$ cyclic action $y_j\mapsto y_{j+1\bmod 4}$ descends to
the phase rotation $v_k\mapsto i^{-k}v_k$ in the spectral basis,
paralleling the $C_5$ rotation $w_k\mapsto\zeta_5^{-k}w_k$ of $B_3$.
}\end{remark}

\begin{remark}[Color tube coordinates for $B_2$]
\DIFadd{\label{rem:B2-color-tube}
The chiral fiber of $B_3$ has three roots $(x_2,x_3,x_4)$ giving three
color tube coordinates $(r,g,b)$
(Definition~\ref{def:color-tube-coordinates}).  For $B_2$ the analog is
a \emph{two-component} chiral fiber: at each $z\in B_1$, the two
branches of $\pi_{21}$ furnish two color coordinates
}\begin{equation}\DIFadd{\label{eq:B2-color-coords}
  \mathbf{r}^{(2)}_z = z + \sqrt{\Delta(z)},
  \qquad
  \mathbf{g}^{(2)}_z = z - \sqrt{\Delta(z)}.
}\end{equation}
\DIFadd{Their sum and product are
}\begin{align}
  \DIFadd{\mathbf{r}^{(2)}_z + \mathbf{g}^{(2)}_z }&\DIFadd{= 2z,
    \label{eq:B2-vieta-sum}}\\[4pt]
  \DIFadd{\mathbf{r}^{(2)}_z \cdot \mathbf{g}^{(2)}_z }&\DIFadd{= z^2 - \Delta(z).
    \label{eq:B2-vieta-product}
}\end{align}
\DIFadd{The $\mathbb{Z}/2\mathbb{Z}$ deck transformation swaps the two branches:
$\iota_2\colon (\mathbf{r}^{(2)}_z, \mathbf{g}^{(2)}_z) \mapsto (\mathbf{g}^{(2)}_z, \mathbf{r}^{(2)}_z)$.
This is the genus-$2$ analog of the $S_3$ permutation of three color
charges in $B_3$; the reduction from $S_3$ to $\mathbb{Z}/2\mathbb{Z}$
reflects the passage from three quark colors to two.
}\end{remark}

\begin{proposition}[Chiral fiber of $B_2$ as $\mathbb{P}^1$]
\DIFadd{\label{prop:B2-chiral-fiber}
For each $z \in B_1\setminus\{z_1,\dots,z_6\}$, the fiber
$(\pi_{21})^{-1}(z)$ consists of exactly two points, making the chiral
fiber of $B_2$ a genus-$0$ curve ($\mathbb{P}^1$).  At the six branch
points $z = z_i$ the fiber degenerates to a single (ramified) point.
}\end{proposition}
\begin{proof}
\DIFadd{Immediate from the definition of $\pi_{21}$ as the hyperelliptic cover
$w^2=\Delta(z)$: for $\Delta(z)\neq 0$ the equation has two solutions
$w=\pm\sqrt{\Delta(z)}$, while for $\Delta(z_i)=0$ the unique solution
is $w=0$.
}\end{proof}

\begin{remark}[Braid group action on $B_2$]
\DIFadd{\label{rem:B2-braid-action}
The mapping class group of the base $B_1 = \mathbb{P}^1$ punctured at
the six branch points $z_1,\dots,z_6$ is the braid group $B_6$ on six
strands.  Its action on $\pi_1$ of the punctured base lifts to the
\emph{symplectic representation}
}\[
  \DIFadd{\rho\colon B_6 \;\longrightarrow\; \mathrm{Sp}(4,\mathbb{Z}),
}\]
\DIFadd{which is the standard representation of the mapping class group
$\mathrm{Mod}_{0,6}$ on the homology $H_1(B_2,\mathbb{Z})\cong\mathbb{Z}^4$.
The monodromy around each branch point $z_i$ is a Dehn twist on $B_2$,
representing a simple transposition in $B_6$.  The full monodromy group
$\Gamma_2 = \mathrm{Sp}(4,\mathbb{Z})$ is the Siegel modular group, whose
orbifold quotient parametrises the principally polarised abelian surfaces
$J(B_2)$ (the Jacobians of $B_2$).  This replaces the $S_5$ automorphism
group of $B_3$ by a lattice group, reflecting the transition from finite
to infinite monodromy as we descend the hierarchy.
}\end{remark}

\begin{remark}[Seiberg--Witten differential for $B_2$]
\DIFadd{\label{rem:B2-SW-differential}
Parallel to the generation-dependent Seiberg--Witten differentials for
$B_3$ (Remark~\ref{rem:seiberg-witten}), the relevant
meromorphic $1$-form on $B_2$ is the \emph{$B_2$ Seiberg--Witten
differential}
}\begin{equation}\DIFadd{\label{eq:B2-SW-diff}
  \lambda_{B_2} \;=\; z\,\omega_1 \;=\; \frac{z\,dz}{w}
  \;=\; \frac{z\,dz}{\sqrt{\Delta(z)}}.
}\end{equation}
\DIFadd{Its $A$-cycle period
$a_{B_2} = \oint_{A_1}\lambda_{B_2}$
and dual $B$-cycle period
$a_{B_2}^D = \oint_{B_1}\lambda_{B_2}$
determine the effective prepotential of the genus-$2$ theory.
Because $B_2$ has genus $2$ rather than $1$, the prepotential carries
two independent moduli $(a_{B_2,1}, a_{B_2,2})$ corresponding to the
two handles of $B_2$; these are identified with the two lightest
generations in the QOU mass hierarchy.
}\end{remark}





\section{\DIFadd{Bring's Surface ($B_3$)}}
\DIFadd{\label{sec:brings-surface}
}

\DIFaddend \begin{remark}[Descent to dimension zero]
\label{rem:descent-to-zero}
By the descent theorem (Theorem~\ref{thm:descartes}), it is sufficient
to study dimension-zero manifolds equipped with their bundles $P$ and~$E$.
At a point $p\in M_0$, the substrate parameter
$q(p)= 1 + \lvert\langle L_\theta,\, F_A^{+}\rangle\rvert(p)$
(Definition~\ref{def:substrate-parameter}) assigns a specific quantum
constant $q\in[1,\infty)$ determined by the principal bundle~$P$ and the
state bundle~$E$.  Here $E$ represents the totality of fermionic matter
packed into spacetime at the Big~Bang.  The trivial bundle $E=E_0$
corresponds to empty space, in which every point's location is exactly
known; a non-trivial bundle~$E$ begins to fill spacetime with fermions,
entering the regime of the quantum Bring surface whose underlying geometry
is non-commutative (Section~\ref{sec:quantum-brings}).
\end{remark}

\begin{definition}[Bring's surface]
\label{def:brings-surface}
Let $x_j$, $j\in\{0,\dots,4\}$, be homogeneous coordinates on
$\mathbb{P}^4$.  \emph{Bring's surface} is the maximally symmetric CR
manifold of genus~$4$ and Euler characteristic $\chi=-6$, defined by the
vanishing of the first three power sums:
\[
  \sum_{j=0}^{4}x_j^{\,k}=0,
  \qquad k=1,2,3.
\]
\end{definition}

\begin{remark}[QOU $4$-tuple for Bring's surface]
\label{rem:brings-qou-4tuple}
We construct the $\mathbb{C}$ quantum observable universe
(Example~\ref{def:c-quantum-observable-universe}) explicitly for Bring's
surface.  The five homogeneous coordinates $x_j$, $j\in\{0,\dots,4\}$,
split into \emph{spacetime} and \emph{chiral} degrees of freedom:
\[
  \underbrace{x_0,\, x_1}_{\text{spacetime}}
  \qquad
  \underbrace{x_2,\, x_3,\, x_4}_{\text{chiral}}.
\]
The projective ratio $z = x_0/x_1 \in \mathbb{CP}^1$ parametrises the
spacetime base.

\medskip
\noindent\textbf{1.\;CR manifold $M_S$.}\;
Let $S$ denote Bring's surface (Definition~\ref{def:brings-surface}).
Because $S$ is a Riemann surface (real dimension~$2$), it does not itself
carry a contact form.  The CR manifold is the unit circle bundle
$M_S = \{(p,\eta) \in T^*S : \lvert\eta\rvert_{g_S} = 1\}$,
where $g_S$ is the Fubini--Study metric restricted to~$S$.  The Liouville
$1$-form $\lambda = \eta\,dz$ on $T^*S$ restricts to a contact form
$\theta = \lambda\big|_{M_S}$, making $M_S$ a $3$-dimensional strictly
pseudoconvex CR manifold.

\medskip
\noindent\textbf{2.\;Principal bundle $P$.}\;
The branched covering $\pi\colon S \to \mathbb{CP}^1$ (degree $4$, six
branch points) lifts to a principal $\mathrm{SU}(2)$-bundle
$P \to M_S$.  In the spacetime coordinate~$z$, the transition functions
of~$P$ are the $\mathrm{SU}(2)$ holonomy matrices
$H_1,H_2$~\eqref{eq:holonomy-generators} around the branch cuts; in the
chiral coordinates $(x_2,x_3,x_4)$, the $S_3$ permutation of the three
sheets gives the $\mathrm{SO}(3)$ holonomy representation
(Definitions~\ref{def:su2-generators}--\ref{def:so3-generators}).

\medskip
\noindent\textbf{3.\;State bundle $E$.}\;
The associated state bundle is $E = P \times_G \mathbb{CP}^1$,
whose fibre at each point of $M_S$ is the chiral
$\mathbb{CP}^1$ parametrised by the three roots of the elliptic
cubic~\eqref{eq:chiral-elliptic}.

\medskip
\noindent\textbf{4.\;Connection, Levi form, and substrate parameter.}\;
The instanton connection $A$ on $P$ is expressed in the resolvent basis
$(w_1,\dots,w_4)$ (Remark~\ref{rem:lagrange-resolvents}) as the flat
$B_3$-monodromy connection whose holonomy around the six branch points is
generated by the trefoil holonomy matrices.  The Levi form on $M_S$ is
\[
  L \;=\; d\theta\big|_{H \times H},
  \qquad H = \ker\theta \subset TM_S,
\]
where $\theta = \eta\,dz\big|_{M_S}$ and the horizontal distribution $H$ is
spanned by the lifts of the chiral differentials $dx_2, dx_3, dx_4$ to
$M_S$.  The substrate parameter
(Definition~\ref{def:substrate-parameter}) then reads
\[
  q(p) \;=\; 1 + \lvert\langle L_\theta,\, F_A^+ \rangle\rvert(p),
  \qquad p \in M_S,
\]
coupling the chiral curvature to the contact geometry.  At the six branch
points (where chiral sheets merge), $\lVert F_A^+\rVert$ is maximal and
$q \to \infty$ (deep-quantum regime); away from the defects,
$q \to 1$ (classicality).

\medskip
\noindent\textbf{Summary.}\;
The $4$-tuple $(M_S,\, G,\, P,\, E)$, with $G = \mathrm{SU}(2)$,
constitutes a $\mathbb{C}$ quantum observable universe
(Example~\ref{ex:brings-surface}), and every QOU structure---contact
form, Reeb field, Levi form, substrate parameter, $q$-deformed
operators---is determined by the spacetime coordinate~$z = x_0/x_1$ and
the chiral coordinates $(x_2, x_3, x_4)$.
\end{remark}

\begin{remark}[Lagrange resolvents and spectral form]
\label{rem:lagrange-resolvents}
Let $\zeta=e^{2\pi i/5}$.  Apply the discrete Fourier transform to the
coordinates $x_j$ of Bring's surface
(Definition~\ref{def:brings-surface}):
\[
  w_k \;=\; \frac{1}{\sqrt{5}}\sum_{j=0}^{4}\zeta^{kj}\,x_j,
  \qquad k\in\{0,1,2,3,4\}.
\]
Under this change of basis the three defining power-sum equations become the
following \emph{spectral forms}:
\begin{align}
  w_0 &= 0
    &&\text{(linear projection),}
    \label{eq:resolvent-linear}\\[4pt]
  w_1\,w_4+w_2\,w_3 &= 0
    &&\text{(Hermitian quadric),}
    \label{eq:resolvent-quadric}\\[4pt]
  w_1^2\,w_3+w_4^2\,w_2+w_2^2\,w_1+w_3^2\,w_4 &= 0
    &&\text{(symmetric cubic).}
    \label{eq:resolvent-cubic}
\end{align}
The cyclic permutation $\sigma\colon x_j\mapsto x_{j+1}$ acts by
$\sigma(w_k)=\zeta^{-k}w_k$, so the $C_5$ symmetry is manifest as a phase
rotation $w_k\mapsto\zeta^k\,w_k$.
\end{remark}

\begin{remark}[Braid-group action on the resolvents]
\label{rem:braid-group-action}
The symmetric group $S_5$ acts on Bring's surface
(Remark~\ref{rem:brings-curve}) by permuting
$(x_0,\dots,x_4)$; this is the maximal automorphism group
(order~$120$) for genus~$4$.  Three distinguished families of $S_3$
subgroups govern the interplay between the base coordinates~$x_j$ and the
spectral resolvents~$w_k$
(Remark~\ref{rem:lagrange-resolvents}).  We describe each in terms of the
braid group (Definition~\ref{def:braid-group}), the resolvent coordinates,
and the $\mathrm{SU}(2)$/$\mathrm{SO}(3)$ matrices of
Definitions~\ref{def:su2-generators}--\ref{def:so3-generators}.

\medskip
\noindent\textbf{(i)~Spacetime action ($\mathbb{CP}^1$ base).}\;
Fix the chiral-like coordinates $x_2,x_3,x_4$ as internal fiber
degrees of freedom.  The constraint
$w_0=0$~\eqref{eq:resolvent-linear} forces $x_0+x_1=C$~(constant), so
the spacetime degrees of freedom collapse projectively to the ratio
$z=x_0/x_1\in\mathbb{CP}^1$.  The transposition
$x_0\leftrightarrow x_1$ acts as the M\"obius involution
$z\mapsto 1/z$, which is a $\pi$-rotation on $\mathbb{CP}^1$ and embeds
into $\mathrm{SU}(2)$ as $\beta=\exp(\pi\,e_1)$~\eqref{eq:su2-2cycle}.
In the resolvent basis this acts as
\begin{equation}\label{eq:spacetime-shift}
  w_k\;\longmapsto\;
  w_k+\tfrac{1}{\sqrt{5}}\,(1-\zeta^k)(x_1-x_0),
  \qquad k=1,\dots,4,
\end{equation}
an affine shift that fixes the chiral modes and translates the
$\mathbb{CP}^1$ fibre.

\medskip
\noindent\textbf{(ii)~Spectral $S_3$ (trefoil action).}\;
The frequency-coupled pairs $(w_1,w_4)$ and $(w_2,w_3)$ are permuted by
the twist and flip involutions
(Definition~\ref{def:braid-twist-flip}).
The full
$S_3$ structure emerges on the projective ratios
\[
  \eta_1=\frac{w_1}{w_2},\qquad\eta_2=\frac{w_3}{w_4}:
\]
$\sigma_1$ inverts both ratios $(\eta_1\mapsto\eta_1^{-1},\;
\eta_2\mapsto\eta_2^{-1})$, while $\sigma_2$ exchanges them
$(\eta_1\mapsto\eta_2^{-1},\;\eta_2\mapsto\eta_1^{-1})$.  The
$2$-cycle (reflection) swaps the pairs
$(w_1,w_4)\leftrightarrow(w_2,w_3)$; the $3$-cycle (rotation) sends
$(\eta_1,\eta_2)\mapsto(\eta_1,\,\zeta^j\eta_2)$ for $j\in\{1,2,3\}$,
permuting the three complex branches of the cubic
constraint~\eqref{eq:resolvent-cubic}.  The
$\mathrm{SU}(2)$ representation assigns the $3$-cycle to
$\alpha$~\eqref{eq:su2-3cycle} and the reflection to
$\beta$~\eqref{eq:su2-2cycle}; the $\mathrm{SO}(3)$
images~\eqref{eq:so3-generators} act as the holonomy on the spectral
energy distribution, with $R_\alpha$ rotating the time-like resolvent
pair~$(\omega_2,\omega_3)$ and $R_\beta$ reflecting the space-like
axis~$\omega_1$.

\medskip
\noindent\textbf{(iii)~Coordinate $S_3$.}\;
Fixing any two coordinates (say $x_0,x_1$) as a spacetime base and
permuting the remaining three $(x_2,x_3,x_4)$ gives a stabiliser
$\mathrm{Stab}(\{x_0,x_1\})\cong S_3\subset S_5$; there are
$\binom{5}{2}=10$ such subgroups.  In the resolvent basis a
permutation $\sigma\in S_3$ of $(x_2,x_3,x_4)$ acts as
\begin{equation}\label{eq:coord-s3-shift}
  w_k\;\longmapsto\;
  w_k+\tfrac{1}{\sqrt{5}}\sum_{j=2}^{4}
    \bigl(\zeta^{k\,\sigma(j)}-\zeta^{kj}\bigr)\,x_j,
\end{equation}
so the $w_k$ transform as vector fields over the $(x_0,x_1)$ plane and
the $S_3$ permutes how these fields couple to the base.  The same
$\mathrm{SO}(3)$ representation~\eqref{eq:so3-generators} realises this
action, now playing the r\^ole of the holonomy group on the internal
fibre.

\medskip
\noindent\textbf{Geometric $S_3$.}\;
When Bring's curve is viewed as a branched cover
$\mathcal{S}\to\mathbb{P}^1$, a third family of $S_3$ subgroups arises:
it permutes the three critical values of the covering map.  These three
families --- spectral, coordinate, and geometric --- correspond to the
three copies of $S_3$ inside $S_5$ identified in
Remark~\ref{rem:chirality-trefoil}.
\end{remark}

\begin{definition}[Color tube coordinates]
\label{def:color-tube-coordinates}
  Let $z = x_0/x_1 \in \mathbb{CP}^1$ be the spacetime ratio
  (Remark~\ref{rem:braid-group-action}\textup{(i)}) and let $w_k$,
  $k \in \{0,\dots,4\}$, be the Lagrange resolvents of Bring's surface
  (Remark~\ref{rem:lagrange-resolvents}).  The \emph{color tube coordinates}
  are the shifted spectral coordinates
  \begin{equation}\label{eq:color-tube-coords}
    r = w_2 - z, \quad g = w_3 - z, \quad b = w_4 - z,
  \end{equation}
  together with the space-like shift $\sigma = w_1 - z$.  The
  \emph{color vector} is $\vec{c} = (r, g, b)$.
\end{definition}

\begin{remark}[Aggregate color quantities]
\label{rem:aggregate-color-quantities}
  Define the following aggregate quantities in color tube coordinates,
  with $\zeta = e^{2\pi i/5}$:
  \begin{align}
    \widetilde{S} &\;=\; \sigma + r + g + b
      &&\text{(total color deviation),}
      \label{eq:agg-S-tilde}\\[4pt]
    \widetilde{Q} &\;=\; \sigma b + r g
      &&\text{(paired color product),}
      \label{eq:agg-Q-tilde}\\[4pt]
    \widetilde{D} &\;=\; \zeta^{-1}\sigma + \zeta^{-2}r + \zeta^{-3}g + \zeta^{-4}b
      &&\text{(weighted color deviation),}
      \label{eq:agg-D-tilde}\\[4pt]
    \widetilde{C} &\;=\; \sigma^2 g + b^2 r + r^2\sigma + g^2 b
      &&\text{(intrinsic color cubic).}
      \label{eq:agg-C-tilde}
  \end{align}
  These are related to the unshifted aggregate quantities
  $S = \sum_{k=1}^{4}w_k$ and $D = \sum_{k=1}^{4}\zeta^{-k}w_k$ by
  \begin{equation}\label{eq:aggregate-shift}
    S \;=\; \widetilde{S} + 4z, \qquad D \;=\; \widetilde{D} - z,
  \end{equation}
  where the second identity uses $\sum_{k=1}^{4}\zeta^{-k} = -1$.
  The inverse Lagrange resolvent transform
  (Remark~\ref{rem:lagrange-resolvents}) gives
  \begin{equation}\label{eq:z-from-SD}
    z \;=\; \frac{x_0}{x_1} \;=\; \frac{S}{D}.
  \end{equation}
\end{remark}

\begin{remark}[Constraint equations in color tube coordinates]
\label{rem:color-tube-constraints}
  The Bring surface constraints~\eqref{eq:resolvent-linear}--\eqref{eq:resolvent-cubic}
  take the following form in color tube coordinates.

  \medskip
  \noindent\textbf{Spacetime recovery.}\;
  Substituting~\eqref{eq:aggregate-shift} into $z = S/D$
  (equation~\eqref{eq:z-from-SD}) yields $z(\widetilde{D} - z) = \widetilde{S} + 4z$,
  which rearranges to
  \begin{equation}\label{eq:spacetime-recovery}
    z^2 - (\widetilde{D} - 4)\,z + \widetilde{S} \;=\; 0.
  \end{equation}
  Thus the spacetime coordinate $z$ is determined implicitly by the
  shifted data $(\widetilde{S}, \widetilde{D})$.

  \medskip
  \noindent\textbf{Shifted quadric.}\;
  Substituting $w_k = c_k + z$ into the quadric~\eqref{eq:resolvent-quadric}
  gives
  \begin{equation}\label{eq:shifted-quadric}
    \widetilde{Q} + z\,\widetilde{S} + 2z^2 \;=\; 0,
  \end{equation}
  or equivalently $\sigma b + rg = -z(\sigma + r + g + b) - 2z^2$.

  \medskip
  \noindent\textbf{Shifted cubic.}\;
  Substituting $w_k = c_k + z$ into the cubic~\eqref{eq:resolvent-cubic}
  gives
  \begin{equation}\label{eq:shifted-cubic}
    \widetilde{C} + z(\widetilde{S}^2 - 2\widetilde{Q})
    + 3z^2\,\widetilde{S} + 4z^3 \;=\; 0.
  \end{equation}
  Eliminating $\widetilde{Q}$ via~\eqref{eq:shifted-quadric} yields the
  \emph{reduced cubic}:
  \begin{equation}\label{eq:reduced-cubic}
    \widetilde{C} + z\,\widetilde{S}^2 + 5z^2\,\widetilde{S} + 8z^3 \;=\; 0.
  \end{equation}
  All three equations are polynomial over $\mathbb{Z}$ in $z, r, g, b, \sigma$;
  they therefore define affine varieties over $\mathbb{Z}$ and descend
  to $\mathbb{F}_\pp$ for every prime $\pp \ne 5$.
\end{remark}

\begin{remark}[Interpretation of the color tube]
\label{rem:color-tube-interpretation}
  The terminology \emph{color tube} reflects two features.

  \medskip
  \noindent\textbf{Color.}\;
  The coordinates $r, g, b$ label the three chiral degrees of freedom
  that correspond, via the $S_3$ subgroups of $S_5$
  (Remark~\ref{rem:braid-group-action}\textup{(iii)}), to the three
  colour charges of quantum chromodynamics.  The chiral coordinates
  $x_2, x_3, x_4$ of the QOU $4$-tuple
  (Remark~\ref{rem:brings-qou-4tuple}) descend to $r = w_2 - z$,
  $g = w_3 - z$, $b = w_4 - z$ after shifting by the spacetime ratio~$z$.

  \medskip
  \noindent\textbf{Tube.}\;
  For fixed $z \in \mathbb{CP}^1$, the color vector $\vec{c} = (r, g, b)$
  is constrained by the shifted quadric~\eqref{eq:shifted-quadric} and
  cubic~\eqref{eq:shifted-cubic} to lie on a curve inside the
  three-dimensional color space $\mathbb{A}^3$.  As $z$ varies over the
  spacetime base, this curve sweeps out a tube --- the total space of
  the chiral fibration over $\mathbb{CP}^1$.

  \medskip
  \noindent\textbf{Space-like coordinate.}\;
  The shift $\sigma = w_1 - z$ is distinguished: it measures the deviation
  of the $C_5$-fundamental resolvent $w_1$ from the spacetime base,
  anchoring the $U(1)$ longitudinal backbone
  (Remark~\ref{rem:spectral-dynamics}) in the shifted frame.
\end{remark}

\begin{remark}[$C_5$ action in color tube coordinates]
\label{rem:c5-color-tube}
  The cyclic symmetry $w_k \mapsto \zeta^{-k} w_k$
  (Remark~\ref{rem:lagrange-resolvents}) becomes, in color tube coordinates:
  \begin{align}
    \sigma &\;\longmapsto\; \zeta^{-1}\sigma + (\zeta^{-1}-1)\,z,
      \label{eq:c5-sigma}\\[4pt]
    r &\;\longmapsto\; \zeta^{-2}r + (\zeta^{-2}-1)\,z,
      \label{eq:c5-r}\\[4pt]
    g &\;\longmapsto\; \zeta^{-3}g + (\zeta^{-3}-1)\,z,
      \label{eq:c5-g}\\[4pt]
    b &\;\longmapsto\; \zeta^{-4}b + (\zeta^{-4}-1)\,z.
      \label{eq:c5-b}
  \end{align}
  The $C_5$ action is no longer diagonal: each coordinate acquires an
  additive correction proportional to~$z$, reflecting the cost of
  breaking projective homogeneity by subtracting the affine
  quantity~$z$ from each resolvent.
\end{remark}

\begin{remark}[Intrinsic color cubic: pairing structure]
\label{rem:intrinsic-color-cubic-pairing}
  The intrinsic color cubic $\widetilde{C}$~\eqref{eq:agg-C-tilde}
  admits a natural decomposition into space-like and chiral sectors:
  \begin{equation}\label{eq:color-cubic-decomp}
    \widetilde{C}
    \;=\;
    \underbrace{\sigma(\sigma g + r^2)}_{\text{space--colour coupling}}
    \;+\;
    \underbrace{b(br + g^2)}_{\text{pure colour}}.
  \end{equation}
  The four terms pair the resolvent indices as $(1,3)$, $(2,1)$, $(4,2)$,
  $(3,4)$: each index is squared and contracted with the index two steps
  ahead in the cyclic order $1 \to 3 \to 4 \to 2 \to 1$.  This
  skip-one pairing respects the conjugation $w_k \leftrightarrow w_{5-k}$
  that pairs $(w_1, w_4)$ and $(w_2, w_3)$.
\end{remark}

\begin{proposition}[Chiral fibre as an elliptic curve]
\label{prop:chiral-elliptic}
Let $z=x_0/x_1\in\mathbb{CP}^1$ be the projective spacetime ratio
(Remark~\ref{rem:braid-group-action}\textup{(i)}).  Substituting the
constraint $x_0+x_1=C$ and eliminating $x_0,x_1$ in favour of~$z$,
the quadratic and cubic power-sum
conditions~\eqref{eq:resolvent-quadric}--\eqref{eq:resolvent-cubic} on
the chiral variable $X=(x_2,x_3,x_4)$ reduce at each fixed~$z$ to the
cubic
\begin{equation}\label{eq:chiral-elliptic}
  X^3+(z+1)\,X^2+(z^2+z+1)\,X+(z^3+z^2+z+1)=0.
\end{equation}
For generic~$z$ this defines an elliptic curve $E_z$ in the chiral
fibre; its $j$-invariant varies as a rational function of~$z$.  The
monodromy of~$E_z$ around the singular fibres (where the discriminant
of~\eqref{eq:chiral-elliptic} vanishes) recovers the holonomy
generators $H_1,H_2$~\eqref{eq:holonomy-generators} of
Remark~\ref{rem:holonomy-trefoil}.
\end{proposition}
\begin{proof}
Write $x_0=z\,x_1$ and use $x_0+x_1=C$ to eliminate both spacetime
variables.  The power-sum conditions $\sum x_j^k=0$ for $k=2,3$ then
become polynomial in~$z$ and $X=x_2+x_3+x_4$, $X_2=\sum_{i<j}x_ix_j$,
etc.  Reducing by the Newton identities yields the
cubic~\eqref{eq:chiral-elliptic}; its discriminant is a polynomial
in~$z$ whose roots give the singular fibres.
\end{proof}

\begin{remark}[Self-duality of the chiral elliptic curve]
\label{rem:t-duality-elliptic}
The cubic~\eqref{eq:chiral-elliptic} is symmetric under the exchange
$X\leftrightarrow z$: writing
\[
  F(X,z)\;=\;X^3+(z+1)\,X^2+(z^2+z+1)\,X+(z^3+z^2+z+1),
\]
one verifies $F(X,z)=F(z,X)$ term by term.  This involution
$\iota\colon(X,z)\mapsto(z,X)$ is a \emph{T-duality}: it identifies
every chiral fibre~$E_z$ with its dual spacetime wrapping~$E_X$, so
the internal and external degrees of freedom of the Bring surface are
algebraically indistinguishable.

\medskip\noindent\textbf{Quotient space.}\;
Passing to the symmetric coordinates $u=X+z$, $v=Xz$, the curve
becomes
\begin{equation}\label{eq:symmetric-quotient}
  v \;=\; \frac{(u+1)(u^2+1)}{2u+1}\,,
\end{equation}
which is rational in~$u$.  Hence the quotient
$E_z/\langle\iota\rangle\cong\mathbb{CP}^1$: folding the torus along
its $X=z$ diagonal collapses it to a genus-$0$ Riemann sphere.

\medskip\noindent\textbf{Weierstrass normal form.}\;
Depressing the cubic via $X=t-\tfrac{z+1}{3}$ eliminates the quadratic
term:
\begin{equation}\label{eq:weierstrass-depressed}
  t^{3}+p(z)\,t+q(z)=0,\qquad
  p=\frac{2z^2+z+2}{3},\quad
  q=\frac{5(4z^3+3z^2+3z+4)}{27}.
\end{equation}
The discriminant $\Delta(z)=-4\,p(z)^3-27\,q(z)^2$ is a polynomial in~$z$
whose roots determine the singular fibres of the elliptic
fibration $\mathcal{S}\to\mathbb{CP}^1$.  At those roots the
$j$-invariant $j(E_z)=-1728\cdot 64\,p(z)^3/\Delta(z)$ degenerates.
\end{remark}

\begin{remark}[Color tube cubic]
\label{rem:color-tube-cubic}
  At fixed $z \in \mathbb{CP}^1$, the three chiral roots of the
  elliptic cubic~\eqref{eq:chiral-elliptic} each satisfy the cubic in
  the color tube coordinate $\mathbf{c} \in \{\mathbf{r}_z, \mathbf{g}_z, \mathbf{b}_z\}$
  obtained by substituting $X = \mathbf{c} + z$
  into equation~\eqref{eq:chiral-elliptic}.  Expanding gives the
  \emph{color tube cubic}:
  \begin{equation}\label{eq:color-tube-cubic}
    \mathbf{c}^3 + (4z+1)\,\mathbf{c}^2 + (6z^2+3z+1)\,\mathbf{c}
    + (4z^3+3z^2+2z+1) \;=\; 0.
  \end{equation}
  This polynomial lies in $\mathbb{Z}[z][\mathbf{c}]$ and defines a
  variety over $\mathbb{Z}$.  The leading $z$-coefficients $4, 6, 4$
  equal $\binom{4}{1}$, $\binom{4}{2}$, $\binom{4}{3}$, and the
  constant term satisfies
  \[
    4z^3 + 3z^2 + 2z + 1 \;=\; \frac{d}{dz}(z^4 + z^3 + z^2 + z).
  \]

  \medskip
  \noindent Vieta's formulas for the three roots give:
  \begin{align}
    \mathbf{r}_z + \mathbf{g}_z + \mathbf{b}_z
      &\;=\; -(4z+1),
      \label{eq:vieta-sum}\\[4pt]
    \mathbf{r}_z\,\mathbf{g}_z + \mathbf{r}_z\,\mathbf{b}_z
      + \mathbf{g}_z\,\mathbf{b}_z
      &\;=\; 6z^2+3z+1,
      \label{eq:vieta-pairs}\\[4pt]
    \mathbf{r}_z\,\mathbf{g}_z\,\mathbf{b}_z
      &\;=\; -(4z^3+3z^2+2z+1).
      \label{eq:vieta-product}
  \end{align}
\end{remark}

\begin{definition}[Color tube fibers]
\label{def:color-tube-fibers}
  Let $z \in \mathbb{CP}^1$.  The \emph{color tube} at~$z$ is the
  affine curve defined over $\mathbb{Z}[z]$ by the color tube
  cubic~\eqref{eq:color-tube-cubic}; we denote it $R_z$, $G_z$, or
  $B_z$ according to which color tube coordinate $\mathbf{r}_z$,
  $\mathbf{g}_z$, or $\mathbf{b}_z$ is used:
  \begin{equation}\label{eq:color-tube-fiber-def}
    R_z \;:=\;
    \bigl\{\,\mathbf{r}_z \in \mathbb{A}^1 :
      \mathbf{r}_z^{\,3} + (4z+1)\,\mathbf{r}_z^{\,2}
      + (6z^2+3z+1)\,\mathbf{r}_z
      + (4z^3+3z^2+2z+1) = 0\,\bigr\},
  \end{equation}
  and $G_z$, $B_z$ are defined by the same equation with coordinates
  $\mathbf{g}_z$ and $\mathbf{b}_z$, respectively.  The coordinates
  are related to the Lagrange resolvents by
  \[
    \mathbf{r}_z = w_2 - z, \qquad \mathbf{g}_z = w_3 - z, \qquad
    \mathbf{b}_z = w_4 - z.
  \]
  As varieties, $R_z \cong G_z \cong B_z$ is isomorphic to the chiral
  elliptic fiber~$E_z$ of Proposition~\ref{prop:chiral-elliptic} via
  the shift $X = \mathbf{c} + z$.
\end{definition}

\begin{definition}[Observed spacetime value]
\label{def:observed-value}
  Let $\mathbf{r}_z, \mathbf{g}_z, \mathbf{b}_z$ be the color tube
  coordinates at a point of Bring's surface
  (Definition~\ref{def:color-tube-fibers}).  The
  \emph{observed spacetime value} is the element of $\mathbb{A}^1$
  recovered from the color tube coordinates by the Vieta
  identity~\eqref{eq:vieta-sum}:
  \begin{equation}\label{eq:observed-value}
    \mathbf{r}_z + \mathbf{g}_z + \mathbf{b}_z + 4z + 1 \;=\; 0.
  \end{equation}
  This identity defines the projection
  $\pi_z \colon R_z \times G_z \times B_z \to \mathbb{A}^1$,
  $({\mathbf{r}_z, \mathbf{g}_z, \mathbf{b}_z}) \mapsto z$,
  which is a polynomial map over $\mathbb{Z}$.  At a branch
  point $z_c$ where chiral sheets merge
  (Remark~\ref{rem:brings-smoothness}), the exceptional divisor is
  coordinatised by the limiting color vector; the observed spacetime
  value reads off $z_c$ from the color data via~\eqref{eq:observed-value},
  providing a coordinate on the exceptional $\mathbb{CP}^1$.
\end{definition}

\begin{remark}[T-duality in color tube coordinates]
\label{rem:t-duality-color-tube}
  The T-duality involution $\iota \colon (X, z) \mapsto (z, X)$
  (Remark~\ref{rem:t-duality-elliptic}) acts on color tube
  coordinates as follows.  Writing $c = X - z$, the dual point has
  base coordinate $z' = X = c + z$ and chiral coordinate $X' = z$,
  giving dual color tube coordinate $c' = z - X = -c$.  Hence
  \begin{equation}\label{eq:t-duality-color-tube}
    \iota \colon (c,\, z) \;\longmapsto\; (-c,\; c + z).
  \end{equation}
  The symmetry $F(X,z) = F(z,X)$
  (Remark~\ref{rem:t-duality-elliptic}) translates to covariance of
  the color tube cubic: the cubic~\eqref{eq:color-tube-cubic} at
  parameter~$z$ with root~$c$ is carried by $\iota$ to the cubic at
  parameter $c + z$ with root~$-c$.  One verifies this directly from
  $F(c+z,\, z) = F(z,\, c+z)$, since $F(z, c+z)$ is the color tube
  cubic at parameter $c+z$ evaluated at~$-c$.
\end{remark}

\begin{remark}[Sublocus $\widetilde{C}=0$ and colour confinement]
\label{rem:color-tube-subloci}
  \textbf{Vanishing of the intrinsic color cubic.}\;
  If $\widetilde{C} = 0$, the reduced cubic~\eqref{eq:reduced-cubic}
  factors as
  \begin{equation}\label{eq:reduced-cubic-factored}
    z\bigl(\widetilde{S}^2 + 5z\,\widetilde{S} + 8z^2\bigr) \;=\; 0,
  \end{equation}
  so either $z = 0$ (the spacetime origin) or $\widetilde{S}$ satisfies
  the quadratic $\widetilde{S}^2 + 5z\,\widetilde{S} + 8z^2 = 0$.
  The discriminant of this quadratic (in $\widetilde{S}$) is
  $25z^2 - 32z^2 = -7z^2$, which is negative for real $z \ne 0$.
  Over $\mathbb{C}$, the two solutions are
  \[
    \widetilde{S} \;=\; z\,\frac{-5 \pm \sqrt{-7}}{2};
  \]
  these are non-real for $z \ne 0$, so the locus $\widetilde{C} = 0$
  meets the real locus only at $z = 0$.  Over $\mathbb{F}_\pp$ the
  existence of a solution depends on whether $-7$ is a quadratic
  residue mod~$\pp$.

  \medskip
  \noindent\textbf{Pure colour sublocus.}\;
  The pure colour term $b(br + g^2)$
  (equation~\eqref{eq:color-cubic-decomp}) vanishes on
  \[
    \{b = 0\} \;\cup\; \{br + g^2 = 0\}.
  \]
  The condition $b = 0$ means $w_4 = z$: the fourth resolvent
  coincides with the spacetime ratio, collapsing one chiral degree of
  freedom onto the base.  The condition $br + g^2 = 0$, i.e.\
  $(w_4 - z)(w_2 - z) + (w_3 - z)^2 = 0$, is a quadratic relation
  among the chiral resolvents and~$z$.  Either condition can be
  interpreted as a \emph{colour-confinement} constraint: one chiral
  mode is frozen in register with the spacetime coordinate, and the
  remaining two are coupled by a Hermitian-quadric relation in the
  shifted frame.
\end{remark}

\begin{remark}[Chiral space as spectral curve and spinor bundle]
\label{rem:chiral-bundles}
The spacetime--chiral decomposition of Bring's surface maps naturally onto
three classical bundle structures over the $\mathbb{CP}^1$ base $z=x_0/x_1$.

\medskip
\noindent\textbf{Cotangent bundle.}\;
The three-sheeted cubic~\eqref{eq:chiral-elliptic} realises the chiral
variables $X=(x_2,x_3,x_4)$ as a \emph{spectral curve}
$\Sigma\subset T^{*}\mathbb{CP}^1$: the base coordinate~$z$ parametrises
spacetime position, and the three roots $X_1,X_2,X_3$ of
Proposition~\ref{prop:chiral-elliptic} are eigenvalues of a Higgs-like
field $\Phi$ on the base.  The six singular fibres (where the discriminant
$\Delta(z)$ vanishes) are the loci at which eigenvalues coincide and the
differential $dX$ degenerates in the cotangent fibre.

\medskip
\noindent\textbf{Spinor bundle.}\;
The spectral modes $w_1,\dots,w_4$
(Remark~\ref{rem:lagrange-resolvents}) live in a rank-$2$ complex vector
bundle over $\mathbb{CP}^1$.  The chiral doublet
$\Psi_L=(w_1,w_2)^T$ and its conjugate $\Psi_R=(w_4,w_3)^T$ transform
under the $\mathrm{SU}(2)$ holonomy generators
$H_1,H_2$~\eqref{eq:holonomy-generators} of
Remark~\ref{rem:holonomy-trefoil}, making $\Psi_L$ a section of the
\emph{spinor bundle} $\mathcal{S}^{+}\to\mathbb{CP}^1$.  The $B_3$
monodromy equips this bundle with a flat connection whose holonomy is the
trefoil-knot group.

\medskip
\noindent\textbf{Adjoint bundle.}\;
The $\mathrm{SO}(3)$ images $R_\alpha,R_\beta$~\eqref{eq:so3-generators}
act on the triple $(\omega_1,\omega_2,\omega_3)$, which therefore spans a
section of the adjoint bundle $\mathrm{ad}\,P\to\mathbb{CP}^1$ associated
to the principal $\mathrm{SU}(2)$-bundle~$P$.  The T-duality involution
$\iota\colon X\leftrightarrow z$ (Remark~\ref{rem:t-duality-elliptic})
interchanges the fibre and base of the adjoint bundle, exchanging
spectral eigenvalues with spacetime coordinates.
\end{remark}

\begin{remark}[Smoothness of the Bring surface and resolution of spacetime defects]
\label{rem:brings-smoothness}
The Bring surface $\mathcal{S}$ (Definition~\ref{def:brings-surface}) is a
\emph{smooth} algebraic variety: the three power-sum constraints are
transverse in~$\mathbb{P}^4$, so their intersection carries no singular
points.  The six topological defects identified in
Remark~\ref{rem:chiral-bundles}---where the discriminant $\Delta(z)$
vanishes and two chiral sheets $X_i,X_j$ collide---are singularities of
the \emph{projection} $\pi\colon\mathcal{S}\to\mathbb{CP}^1$, not of the
surface itself.

\medskip
\noindent\textbf{Resolution in the $x_j$ coordinates.}\;
At a branch point $z=z_c$ where, say, $x_2=x_3$, define the local chiral
coordinate $\delta=x_2-x_3$.  In a neighbourhood of the defect the
spacetime ratio satisfies
\[
  z-z_c\;\sim\;\delta^2,
\]
so while the projection map develops a ramification cusp, $\delta$ passes
smoothly through zero, carrying the geometry from one sheet to the other.
The exceptional divisor of the blow-up at each branch point is therefore
already embedded in the chiral coordinates $x_2,x_3,x_4$: no external
resolution is required.

\medskip
\noindent\textbf{Resolution in the $w_k$ coordinates.}\;
The DFT mixes all five $x_j$ into every resolvent $w_k$
(Remark~\ref{rem:lagrange-resolvents}), so no single $w_k$ vanishes at a
branch point.  The Briggs--Klein chart
(Definition~\ref{def:briggs-klein-chart}) remains a smooth coordinate
system throughout $0<r<\infty$: the radial parameter~$r=\lvert w_1\rvert^2$
is a regular function on $\mathcal{S}$, and the $C_5$ orbits
$w_k\mapsto\zeta^k w_k$ foliate the surface without fixed points.  The
infinite torsion and eigenvalue crossings visible in the base
$\mathbb{CP}^1$ are artefacts of projecting away the internal degrees of
freedom that the $w_k$ retain.
\end{remark}

\begin{definition}[Briggs--Klein affine chart]
\label{def:briggs-klein-chart}
Let $r=\lvert w_1\rvert^2$ denote the squared modulus of the first resolvent
coordinate on Bring's surface
(Definition~\ref{def:brings-surface}).  The
\emph{Briggs--Klein affine chart} is the coordinate system $(w_1,w_2)$
restricted to the open set $0<r<\infty$, equipped with the Hermitian metric
inherited from the Fubini--Study metric on $\mathbb{P}^4$.  In this chart
the $C_5$ action is the phase rotation $w_k\mapsto\zeta^k\,w_k$, and the
radial coordinate $r$ is $C_5$-invariant.
\end{definition}

\begin{remark}[Radial projection of Bring's surface]
\label{rem:radial-projection}
Restricting a Briggs--Klein affine chart
(Definition~\ref{def:briggs-klein-chart}) to the interval $r\in(0,1)$ and
projecting Bring's surface onto this radial coordinate yields a fibration
over the base manifold $M$ with trivial bundle $E=E_0$.  At each
$r$, the coordinate $w_1$ is restricted to a circle of radius $\sqrt{r}$ in
the complex plane $\mathbb{C}_{w_1}$.  This circle is the orbit of the $C_5$
symmetry: the rotation $w_1\mapsto\zeta\,w_1$ moves the point along this
circle without changing~$r$.

\medskip
\begin{center}
\begin{tikzpicture}[>=stealth, scale=2.5]
\fill[brown!12]
    plot[domain=0.02:1,samples=50] (\x,{sqrt(\x)})
    -- plot[domain=1:0.02,samples=50] (\x,{-sqrt(\x)}) -- cycle;
\foreach \rv in {0.1,0.2,...,0.9}{
    \draw[brown!35,thin] (\rv,{-sqrt(\rv)}) -- (\rv,{sqrt(\rv)});
  }
\foreach \hv in {-0.9,-0.7,...,0.9}{
    \draw[brown!35,thin]
      plot[domain={\hv*\hv}:1,samples=30] (\x,{\hv});
  }
\draw[thick,domain=0.005:1.08,samples=80]
    plot (\x,{sqrt(\x)}) node[right]{$\sqrt{r}$};
\draw[thick,domain=0.005:1.08,samples=80]
    plot (\x,{-sqrt(\x)});
\draw[->] (-0.08,0) -- (1.22,0) node[right]{$r$};
  \draw[->] (0,-1.18) -- (0,1.18) node[above]{$\mathbb{C}_{w_1}$};
\draw (1,0.04) -- (1,-0.04) node[below]{\small$1$};
\draw[dashed,gray] (0.5,{-sqrt(0.5)}) -- (0.5,{sqrt(0.5)});
\draw[thick,->] ({0.5+sqrt(0.5)*cos(55)},{sqrt(0.5)*sin(55)})
    arc[start angle=55, end angle=125, radius={sqrt(0.5)}];
  \node at (0.5,{sqrt(0.5)+0.13}) {\small$\zeta$};
\node at (0.78,0.22) {$E=E_0$};
  \node[below] at (0.5,-0.08) {\small$M$};
\end{tikzpicture}
\end{center}
\noindent
The brown tessellation filling the brane illustrates the trivial bundle
$E=E_0$: every cell is flat, reflecting the absence of torsion on the total
space.
\end{remark}

\begin{remark}[Four-colourability of Bring's surface]
\label{rem:four-colourability}
By the Heawood bound, an orientable surface of genus~$g$ has chromatic number
at most $\lfloor(7+\sqrt{1+48g})/2\rfloor$.  For genus~$4$ this bound
equals~$10$.  Bring's surface, however, admits a proper colouring with only
four colours --- considerably fewer than the maximum permitted by its genus.
This low chromatic number reflects the high degree of symmetry: the $S_5$
automorphism group constrains the combinatorics of any cell decomposition,
forcing a more efficient colouring than a generic genus-$4$ surface would
require.
\end{remark}

\begin{remark}[The four-colour theorem and Bring's surface]
\label{rem:four-colour-theorem}
The classical four-colour theorem asserts that every planar graph (equivalently,
every graph embeddable in $S^2$) is four-colourable.  For higher-genus
surfaces the chromatic number can be much larger: the Heawood conjecture
(proved by Ringel--Youngs for $g\geq 1$) shows that the maximum chromatic
number grows as $O(\sqrt{g})$.  That Bring's surface attains the planar
bound of four --- despite living at genus~$4$ --- is a reflection of the
rigidity imposed by its $S_5$ symmetry.  In the language of the Descartes
universe, the trivial bundle $E=E_0$
(Remark~\ref{rem:radial-projection}) carries no torsion, and the resulting
flatness of the tessellation forces the same chromatic economy as the
sphere.
\end{remark}

\begin{remark}[Cyclotomic reduction and $C_{\mathfrak{p}}$ preservation]
\label{rem:cyclotomic-reduction}
Let $\mathfrak{p}$ be a prime (not to be confused with a point $p\in M$).
Working over the cyclotomic ring
\begin{equation}\label{eq:cyclotomic-ring}
  \mathbb{Z}_{\zeta_{\mathfrak{p}}}
  \;:=\;\mathbb{Z}[i,\,\zeta_{\mathfrak{p}}],
\end{equation}
where $\zeta_{\mathfrak{p}}=e^{2\pi i/\mathfrak{p}}$ is a primitive
$\mathfrak{p}$-th root of unity, consider the sublocus
\[
  S_{\mathfrak{p}}
  \;:=\;
  \Sigma_4\;\cap\;\mathbb{Z}_{\zeta_{\mathfrak{p}}}^{\,5},
\]
that is, the set of points on Bring's surface whose resolvent coordinates
lie in $\mathbb{Z}_{\zeta_{\mathfrak{p}}}$.  The cyclic subgroup
$C_{\mathfrak{p}}\subset S_5$ generated by the $\mathfrak{p}$-cycle preserves
$S_{\mathfrak{p}}$: for any $P\in S_{\mathfrak{p}}$, the image
$\sigma^k(P)$ remains in $\mathbb{Z}_{\zeta_{\mathfrak{p}}}^{\,5}$ because
$\sigma$ acts by multiplication by $\zeta_{\mathfrak{p}}^{-k}$ in the
resolvent basis (Remark~\ref{rem:lagrange-resolvents}), and
$\zeta_{\mathfrak{p}}^{-k}\in\mathbb{Z}_{\zeta_{\mathfrak{p}}}$.  This
arithmetic compatibility between the cyclic symmetry and the cyclotomic ring
underlies the descent of the $C_5$ action to finite characteristic.
\end{remark}

\begin{definition}[Quaternionic spectral resolvents]
\label{def:quaternionic-spectral-resolvents}
Let $\mathcal{S}$ denote the Bring's surface
(Definition~\ref{def:brings-surface}).  Let $q\in\mathbb{C}$ with
$0<\lvert q\rvert\le 1$, and set $r=-\ln\lvert q\rvert$ so that
$\lvert q\rvert=e^{-r}$.
Write $\mathbf{i},\mathbf{j},\mathbf{k}$ for the standard basis of
$\operatorname{Im}\mathbb{H}$.
Using the Lusztig $q$-exponential and $q$-logarithm,
\begin{equation}\label{eq:q-exp}
  \exp_q(x)\;=\;\sum_{n=0}^{\infty}\frac{x^n}{[n]_q!},
  \qquad
  \ln_q(1+x)\;=\;\sum_{n=1}^{\infty}(-1)^{n+1}\frac{x^n}{[n]_q},
\end{equation}
we define the \emph{quaternionic spectral resolvents}
$\omega_k\colon\mathcal{S}\to J^\infty E$ by
\begin{align}
  \omega_0 &\;=\;
    \Bigl(\tfrac{1}{\sqrt{5}}\!\sum_{j=0}^{4}x_j\Bigr)\,
    \ln_q(1+r)
    &&\text{(uncertainty),}
    \label{eq:omega-0}\\[4pt]
  \omega_1 &\;=\;
    \frac{\mathbf{i}}{\ln_q(1+r)}\,
    \Bigl(\tfrac{1}{\sqrt{5}}\!\sum_{j=0}^{4}\zeta^{j}\,x_j+i\Bigr)
    &&\text{(space-like),}
    \label{eq:omega-1}\\[4pt]
  \omega_2 &\;=\;
    \mathbf{j}\,
    \Bigl(\tfrac{1}{\sqrt{5}}\!\sum_{j=0}^{4}\zeta^{2j}\,x_j+i\Bigr)\,
    \bigl(\exp_q(r)-1\bigr)
    &&\text{(time-like),}
    \label{eq:omega-2}\\[4pt]
  \omega_3 &\;=\;
    \mathbf{k}\,
    \Bigl(\tfrac{1}{\sqrt{5}}\!\sum_{j=0}^{4}\zeta^{3j}\,x_j+i\Bigr)\,
    \bigl(\exp_q(r)-1\bigr)
    &&\text{(time-like).}
    \label{eq:omega-3}
\end{align}
By the first power sum~\eqref{eq:resolvent-linear}, $\omega_0=0$ in the
balanced state; the remaining coordinates embed $\mathcal{S}$ into the
infinite jet shell $J^\infty E$
(Remark~\ref{rem:jet-shell}).
\end{definition}

\begin{remark}[Algebraic constraints in the $\omega_k$ basis]
\label{rem:omega-constraints}
In the spectral resolvent basis
(Definition~\ref{def:quaternionic-spectral-resolvents}) the
Bring's surface equations take the following form:
\begin{itemize}
\item \emph{Barycentric:}\; $\omega_0=0$.
\item \emph{Quadric ($\mathbb{H}$ condition):}\;
  $\omega_1\bar{\omega}_1+\omega_2\bar{\omega}_3=0$.
  This preserves the horizontal self-duality $F_A|_H={*_H}F_A|_H$
  (Proposition~\ref{prop:self-dual-curvature}) and enforces a
  Minkowski-like balance between the space-like and time-like modes.
\item \emph{Cubic (phase locking):}\; a quintic resonance condition that
  locks the expansion of $\ln_q$ to $\exp_q$, preserving the $C_5$
  symmetry and the Descartes condition $\chi=-6$.
\end{itemize}
\end{remark}

\begin{remark}[Spectral dynamics and the $\mathrm{Sp}(1)$ action]
\label{rem:spectral-dynamics}
The growth profiles in
Definition~\ref{def:quaternionic-spectral-resolvents} separate the space-like
and time-like sectors of the Bring's surface:
\begin{itemize}
\item \emph{$\omega_1$ (space-like):}\; the factor $1/\ln_q(1+r)$ diverges
  as $r\to 0$, so $\omega_1$ achieves infinite density at the base stratum
  $M_0$, characterising the primordial vortex.  This mode anchors the
  $U(1)$ longitudinal backbone of the $\mathrm{SU}(2)$ bundle.
\item \emph{$\omega_2,\omega_3$ (time-like):}\; the factor
  $\exp_q(r)-1$ vanishes at $r=0$ and grows rapidly for small $r$,
  driven by the $q$-deformed Reeb flow
  (Definition~\ref{def:q-reeb-flow}).  These ladder modes expand at the big
  bang before tapering as the geometric viscosity
  (Remark~\ref{rem:q-harmonic-physical}) takes hold.
\end{itemize}
To realise the rotation as a function of $r$, embed the radial parameter
into $\mathrm{Sp}(1)$ via the unit quaternion
\begin{equation}\label{eq:sp1-quaternion}
  \mathbf{q}(r)
  \;=\;
  \cos_q\!\theta(r)+\mathbf{u}\,\sin_q\!\theta(r),
  \qquad
  \theta(r)=\tfrac{\pi}{2}\,e^{-r},
\end{equation}
where $\cos_q$ and $\sin_q$ are the $q$-deformed trigonometric functions
(built from $\exp_q$, equation~\eqref{eq:q-exp})
\[
  \cos_q(x)=\frac{\exp_q(ix)+\exp_q(-ix)}{2},
  \qquad
  \sin_q(x)=\frac{\exp_q(ix)-\exp_q(-ix)}{2i},
\]
and $\mathbf{u}\in\operatorname{Im}\mathbb{H}$ is a fixed unit imaginary
quaternion.  The conjugation
$\boldsymbol{\omega}'=\mathbf{q}(r)\,\boldsymbol{\omega}\,\mathbf{q}(r)^{*}$
gives maximum twist ($\theta=\pi/2$) in the deep-quantum regime $r\to 0$
and classical freeze ($\theta\to 0$) as $r\to\infty$.
\end{remark}

\begin{remark}[Primordial vortex and the deep-quantum limit]
\label{rem:primordial-vortex}
As $r\to 0$ the spectral resolvents converge toward the
\emph{primordial vortex}:
\begin{itemize}
\item The quantum Planck constant
  $\hbar_q\to 1$ (Definition~\ref{def:quantum-planck-constant}).
\item The $q$-torsion leaves
  (Definition~\ref{def:q-torsion-leaf}) collapse to the zero-dimensional
  stratum $M_0=\{\pm 1\}$.
\item The $\omega_0$ coordinate, although identically zero in the
  balanced state, formally blows up in the jet prolongation, representing
  the concentration of energy at a singular spatial point.
\end{itemize}
Space-like energy is trapped because the time-like rotation becomes infinite
in frequency but zero in displacement, effectively turning the
$\mathrm{SU}(2)$ holonomy into a localised topological charge.
\end{remark}

\begin{remark}[Geometric drag and the mass residue]
\label{rem:geometric-drag}
The Levi form $L$ (Definition~\ref{def:substrate-parameter}) must be
preserved under the $\mathrm{Sp}(1)$ conjugation.  Since
\[
  \lVert\boldsymbol{\omega}'\rVert^2
  =\lVert\boldsymbol{\omega}\rVert^2
  =\omega_1\bar{\omega}_1+\omega_2\bar{\omega}_2+\omega_3\bar{\omega}_3,
\]
the quadric constraint~\eqref{eq:resolvent-quadric} remains satisfied and
the horizontal self-duality $F_A|_H={*_H}F_A|_H$ is preserved.  The mass of a stable vortex is
recovered as the residue of the interaction between the growth profiles:
\begin{equation}\label{eq:mass-residue}
  m \;\propto\;
  \oint_{\gamma}
  \lVert\nabla_A q\rVert\cdot
  \lVert\Omega_1\wedge\Omega_{2,3}\rVert,
\end{equation}
where $\Omega_k$ denotes the curvature form restricted to the $\omega_k$
sector.  When $q=1$ (trivial bundle), the $\omega_k$ commute and $m=0$;
when $q\to 0$ (maximal twist), energy is trapped in the $\omega_1$ spatial
filament.
\end{remark}



\begin{remark}[Geometric resolution and the Seiberg--Witten differential]
\label{rem:seiberg-witten}
The infinite geometric torsion~$T$ and the six topological defects on the
$\mathbb{CP}^1$ base $z=x_0/x_1$ are strictly artefacts of
dimensional projection.  The total space, when lifted into the
complete coordinate framework, is a globally smooth complex manifold.

\medskip
\noindent\textbf{1.~The blow-up and the exceptional divisor.}\;
Near a spacetime defect $z_c$ where chiral sheets (e.g.\ $x_2$ and
$x_3$) merge, the singularity is resolved by blowing up the point into
an exceptional divisor parameterised by the local coordinate
$\delta=x_2-x_3$.  The local geometry satisfies
\[
  z-z_c\;\sim\;\delta^2,
\]
proving that the spacetime base ``bounces'' at the defect: the chiral
coordinate continuously glides through the exceptional divisor and
transitions between sheets without encountering a singularity
(cf.\ Remark~\ref{rem:brings-smoothness}).

In the spectral domain, the Jacobian determinant of the DFT
(Remark~\ref{rem:lagrange-resolvents}) is a constant
$\lvert\det(J)\rvert=1$, so the global volume form never collapses:
\[
  \Omega \;=\; dw_1\wedge dw_2\wedge dw_3\wedge dw_4
  \;=\; \det(J)\;dx_0\wedge\cdots\wedge dx_4\Big|_{w_0=0}.
\]
The geometric torsion~$T$ is thereby localised not as a point
singularity but as a trapped, finite magnetic flux integrated over
the surface of the exceptional $\mathbb{CP}^1$ sphere.

\medskip
\noindent\textbf{2.~The Seiberg--Witten differential.}\;
With the singularities geometrically resolved, the dynamics of the chiral
space are completely parameterised by a meromorphic $1$-form embedded in
the cotangent bundle $T^{*}\mathbb{CP}^1$, the \emph{Seiberg--Witten
differential}:
\begin{equation}\label{eq:seiberg-witten-differential}
  \lambda_{\mathrm{SW}} \;=\; X\,dz,
\end{equation}
where $X\in\{x_2,x_3,x_4\}$ is the internal momentum (Higgs-field
eigenvalue) defining the spectral curve
(Remark~\ref{rem:chiral-bundles}), and $z$ is the spacetime position.
The exact physical properties of the topological defects---specifically
the mass/energy spectrum generating the torsion---are rigorously
determined by integrating this differential along the fundamental
non-contractible loops ($A$ and $B$ cycles) of the resolved chiral
torus:
\begin{equation}\label{eq:seiberg-witten-periods}
  a \;=\; \oint_{A}\lambda_{\mathrm{SW}},
  \qquad
  a_D \;=\; \oint_{B}\lambda_{\mathrm{SW}}.
\end{equation}

By lifting the Bring surface into the full spatial and spectral
coordinate spaces, the degenerate branch points resolve completely into
smooth exceptional divisors.  The infinite local torsion is geometrically
captured as trapped flux, revealing that the discrete braid-group
($B_3$) holonomy (Remark~\ref{rem:holonomy-trefoil}) is simply the
continuous Kodaira monodromy of a globally smooth Seiberg--Witten gauge
theory.
\end{remark}



\begin{remark}[Reeb vector field on the exceptional divisor]
\label{rem:reeb-exceptional-divisor}
The Reeb vector field $\xi_q$
(Definition~\ref{def:q-reeb-vector-field}) is the dynamical
complement to the static topology of the Bring surface.  While the
surface itself provides the handles (genus~$4$), the Reeb vector
provides the ``wind'' blowing through them, and the quantum torsion~$q$
sets the strength of that wind.

\medskip
\noindent\textbf{The classical limit ($q\to 1$): cones and frozen geometry.}\;
If the substrate parameter $q\to 1$
(Definition~\ref{def:substrate-parameter}), the
torsion coupling vanishes and the phase rotation
$w_k\mapsto\zeta^k\,w_k$ loses its temporal component.  The fibre
over each base point ceases to twist as one moves along a path on
$\mathbb{CP}^1$: the chiral tubes are frozen in the configuration
determined by the power-sum equations.  Metrically this is a conical
limit---the curvature concentrates at the blow-up points but no
circulating torsion field smooths it out.  In terms of the Reeb
flow (Definition~\ref{def:q-reeb-flow}), $\xi_q\to\xi_1=\xi$ becomes
the classical Reeb field, and the Heisenberg commutator
$[\hat{X},\hat{P}]_q\to 0$
(Remark~\ref{rem:heisenberg-commutator}): the Bring variety
collapses to a non-rotating algebraic surface.

\medskip
\noindent\textbf{Exceptional divisor as mixing interface.}\;
Blowing up a spacetime defect $z_c$ where chiral sheets merge
(Remark~\ref{rem:brings-smoothness}) replaces the singular point with
an exceptional divisor $E_c\cong\mathbb{CP}^1\cong S^2$.  This
sphere acts as the boundary between the macroscopic coordinates
$(x_0,x_1)$ and the internal chiral modes $(w_1,w_2,w_3)$: the
self-intersections of the chiral fibres are resolved, and the
$C_5$ phase rotation now circulates \emph{around} the exceptional
divisor rather than passing through a singularity.

\medskip
\noindent\textbf{Reeb vector as torsion field.}\;
On the blown-up surface, the Reeb vector field~$\xi_q$ is the
infinitesimal generator of the $U(1)$ action that rotates the chiral
fibre above the exceptional divisor.  Concretely, the
$q$-deformed Reeb flow $\phi_t^q$
(Definition~\ref{def:q-reeb-flow}) restricted to $E_c$ generates
the Hopf fibration $S^3\to S^2\cong E_c$, and the substrate
parameter~$q$ controls the pitch of this rotation.  The vector
field~$\xi_q$ thereby provides the ``time-like'' direction within the
chiral space: it is the flow that allows the spectral modes~$w_k$ to
execute the spinning phase rotation
$w_k\mapsto\zeta^k\,w_k$.

\medskip
\noindent\textbf{Global picture.}\;
The Bring surface provides the genus-$4$ topology---four handles
through which the chiral tubes thread.  The Reeb vector field gives
the longitudinal flow along these handles.  The quantum torsion~$q$
modulates the pitch of this flow: when $q\to\infty$ (deep-quantum,
Big Bang) the spinning is maximal and the Heisenberg uncertainty
$[\hat{X},\hat{P}]_q\to J$; when $q\to 1$ (heat death) the spinning
ceases, the Reeb field trivialises, and the topology reduces to a
static cone.
\end{remark}


\begin{remark}[Generation-Dependent Seiberg-Witten Differentials]
The classical Seiberg-Witten integrals over the $A$ and $B$ cycles [cite: 954] fail to distinguish lepton generations. In the QOU framework, the vertical degree of the Reeb decomposition lifts the geometry from the base state bundle $E_1$ to the jet prolongation $J^1E_1$. Consequently, the meromorphic 1-form governing the chiral dynamics is strictly generation-dependent.

For the electron (vertical degree 1), the mass is generated by the baseline Seiberg-Witten differential smoothly integrating over the resolved exceptional divisor $E_c$:
$$a_D^{(e)} = \oint_{E_c} \lambda_{SW}^{(0)} = \oint_{E_c} X dz = \oint_{E_c} (c + z) dz$$

For the muon (vertical degree 2), the gauge field is lifted to the 1-jet bundle, forcing the integration of the momentum gradient. The prolonged differential is $\lambda_{SW}^{(1)} = dX$. Substituting the implicit differentiation of the color tube cubic yields the explicit dual mass integral:
$$a_D^{(\mu)} = \oint_{E_c} \lambda_{SW}^{(1)} = \oint_{E_c} \frac{4z^3+3z^2+2z+1}{c \left( 3c^2 + 2(4z+1)c + (6z^2+3z+1) \right)} dz$$

At the topological defect $z_c$ where chiral sheets merge, the discriminant vanishes and $\partial F / \partial c \to 0$. This induces a severe pole in the prolonged differential $\lambda_{SW}^{(1)}$. By Cauchy's Residue Theorem, the integral over the exceptional divisor captures the residue of this jet-bundle gradient. This geometrically inflates the trapped quantum dilogarithm volume, producing the mass hierarchy $m_\mu \gg m_e$ purely from the topology of the blowup, without requiring arbitrarily high crossing numbers.
\end{remark}


\section{Quantum Brings}
\label{sec:quantum-brings}





\section{Other Euler characteristics}
\label{sec:other-euler}

 


\DIFaddbegin \section{\DIFadd{Knots on Surfaces of Genus at Most Four}}
\DIFadd{\label{sec:knots-genus-leq-4}
}

\subsection{\DIFadd{Particle ansatz}}
\DIFadd{\label{subsec:particle-ansatz}
}

\begin{remark}[Knot--particle correspondence]
\DIFadd{\label{rem:knot-particle-correspondence}
The Bring hierarchy provides the geometric scaffolding for a knot--particle
correspondence: stable particle states of the QOU are identified with
isotopy classes of knots (or links) that can be drawn on a surface of
genus at most $4$ --- i.e.\ knots $K$ with Seifert genus
$g_s(K) \le 4$.  The embedding surface is one of the $B_n$:
}\begin{itemize}
  \item \DIFadd{$g_s(K)=0$: $K$ embeds on $B_1 = \mathbb{P}^1$; corresponds to
        massless/trivial states (neutrinos).
  }\item \DIFadd{$g_s(K)=1$: $K$ embeds on the elliptic fiber of $B_3$;
        corresponds to first-generation fermions (electron, up/down
        quark).
  }\item \DIFadd{$g_s(K)=2$: $K$ embeds on $B_2$; second-generation
        fermions (muon, strange/charm quark).
  }\item \DIFadd{$g_s(K)=3,4$: $K$ embeds on $B_3$; third-generation
        fermions (tau, bottom/top quark) and heavier states.
}\end{itemize}
\DIFadd{Composite particles (hadrons, nuclei) correspond to connected sums
$K_1 \mathbin{\#} K_2$ whose Seifert genus equals $g_s(K_1)+g_s(K_2)$.
}\end{remark}

\begin{remark}[Knots embeddable in Bring's surface]
\DIFadd{\label{rem:knots-embeddable}
A knot $K \subset S^3$ embeds on a Seifert surface of genus $g$ if and
only if $g_s(K) \le g$.  Since Bring's surface $B_3$ has genus $4$, the
knots relevant to the QOU are exactly those with $g_s(K) \le 4$.  Every
torus knot $T(p,q)$ with Seifert genus
$g_s(T(p,q)) = \tfrac{(p-1)(q-1)}{2} \le 4$ embeds.  The fundamental
knots for the QOU particle spectrum are:
}\begin{itemize}
  \item \DIFadd{$0_1$ (unknot), $g_s=0$: neutrino.
  }\item \DIFadd{$2^2_1$ (Hopf link), $g_s=0$: gauge boson (photon).
  }\item \DIFadd{$3_1 = T(2,3)$, $g_s=1$: quark (trefoil, first generation).
  }\item \DIFadd{$4_1$, $g_s=1$: amphichiral first-generation state.
  }\item \DIFadd{$5_1 = T(2,5)$, $g_s=2$: second-generation fermion.
  }\item \DIFadd{$5_2$, $g_s=1$: excited first-generation state.
  }\item \DIFadd{$T(3,4)$, $g_s=3$: third-generation fermion.
  }\item \DIFadd{$T(2,7)$, $g_s=3$: alternative third-generation state.
  }\item \DIFadd{$T(3,5)$, $g_s=4$: maximal-genus fermion on $B_3$.
  }\item \DIFadd{$T(2,9)$, $g_s=4$: tau-like state at the boundary.
}\end{itemize}
\DIFadd{Composite knots $3_1 \mathbin{\#} 3_1$ ($g_s=2$) and
$3_1 \mathbin{\#} 4_1$ ($g_s=2$) correspond to two-quark bound
states.
}\end{remark}

\subsection{\DIFadd{Knot invariant table}}
\DIFadd{\label{subsec:knot-invariant-table}
}

\DIFadd{The following table lists invariants for each knot.  Columns:
\textbf{Name}: standard Rolfsen/Conway notation.
\textbf{$\Delta(t)$}: Alexander polynomial (symmetric, $\Delta(1)=1$).
\textbf{$\nabla(z)$}: Conway polynomial ($\Delta(t)=\nabla(t^{1/2}-t^{-1/2})$).
\textbf{$V(t)$}: Jones polynomial.
\textbf{$\mathrm{vol}$}: hyperbolic volume of the complement ($0$ for
non-hyperbolic).
\textbf{A/C}: amphichiral~(A) or chiral~(C).
\textbf{$g_s$}: Seifert (minimal) genus.
\textbf{Tw}: twist knot (Y/N).
\textbf{Bl}: algebraic link / singularity blowup (Y/N).
}

\bigskip
{\renewcommand{\arraystretch}{1.4}
\noindent
\begin{tabular}{llllrcccl}
\hline
\DIFadd{\textbf{Name} }& \DIFadd{\textbf{$\Delta(t)$} }& \DIFadd{\textbf{$\nabla(z)$} }& \DIFadd{\textbf{$V(t)$}
  }& \DIFadd{\textbf{vol} }& \DIFadd{\textbf{A/C} }& \DIFadd{\textbf{$g_s$} }& \DIFadd{\textbf{Tw} }& \DIFadd{\textbf{Bl} }\\
\hline
\DIFadd{$0_1$ unknot
  }& \DIFadd{$1$
  }& \DIFadd{$1$
  }& \DIFadd{$1$
  }& \DIFadd{$0$
  }& \DIFadd{A }& \DIFadd{$0$ }& \DIFadd{Y }& \DIFadd{Y }\\[2pt]
\DIFadd{$2^2_1$ Hopf
  }& \DIFadd{$-1$ (link)
  }& \DIFadd{$z$
  }& \DIFadd{$-t^{-1/2}-t^{-5/2}$
  }& \DIFadd{$0$
  }& \DIFadd{C }& \DIFadd{$0$ }& \DIFadd{--- }& \DIFadd{Y }\\[2pt]
\DIFadd{$3_1$ trefoil
  }& \DIFadd{$t^{-1}-1+t$
  }& \DIFadd{$1+z^2$
  }& \DIFadd{$-t^{-4}+t^{-3}+t^{-1}$
  }& \DIFadd{$0$
  }& \DIFadd{C }& \DIFadd{$1$ }& \DIFadd{Y }& \DIFadd{Y }\\[2pt]
\DIFadd{$4_1$ fig.-eight
  }& \DIFadd{$-t^{-1}+3-t$
  }& \DIFadd{$1-z^2$
  }& \DIFadd{$t^2-t+1-t^{-1}+t^{-2}$
  }& \DIFadd{$2.029$
  }& \DIFadd{A }& \DIFadd{$1$ }& \DIFadd{Y }& \DIFadd{N }\\[2pt]
\DIFadd{$5_1 = T(2,5)$
  }& \DIFadd{$t^{-2}-t^{-1}+1-t+t^2$
  }& \DIFadd{$1+3z^2+z^4$
  }& \DIFadd{$-t^{-4}+t^{-3}-t^{-2}+t^{-1}-1+t$
  }& \DIFadd{$0$
  }& \DIFadd{C }& \DIFadd{$2$ }& \DIFadd{N }& \DIFadd{Y }\\[2pt]
\DIFadd{$5_2$
  }& \DIFadd{$2t^{-1}-3+2t$
  }& \DIFadd{$1+2z^2$
  }& \DIFadd{$-t^6+t^5+t^3$
  }& \DIFadd{$2.828$
  }& \DIFadd{C }& \DIFadd{$1$ }& \DIFadd{Y }& \DIFadd{N }\\[2pt]
\DIFadd{$T(3,4)$
  }& \DIFadd{$t^{-3}-t^{-2}+1-t^2+t^3$
  }& \DIFadd{$1+2z^2+z^4$
  }& \DIFadd{see~\eqref{eq:jones-T34}
  }& \DIFadd{$0$
  }& \DIFadd{C }& \DIFadd{$3$ }& \DIFadd{N }& \DIFadd{Y }\\[2pt]
\DIFadd{$T(2,7)$
  }& \DIFadd{$t^{-3}-t^{-2}+t^{-1}-1+t-t^2+t^3$
  }& \DIFadd{$1+4z^2+3z^4+z^6$
  }& \DIFadd{see~\eqref{eq:jones-T27}
  }& \DIFadd{$0$
  }& \DIFadd{C }& \DIFadd{$3$ }& \DIFadd{N }& \DIFadd{Y }\\[2pt]
\DIFadd{$T(3,5)$
  }& \DIFadd{$t^{-4}-t^{-3}+t^{-1}-1+t-t^3+t^4$
  }& \DIFadd{$1+3z^2+z^4+z^6$
  }& \DIFadd{see~\eqref{eq:jones-T35}
  }& \DIFadd{$0$
  }& \DIFadd{C }& \DIFadd{$4$ }& \DIFadd{N }& \DIFadd{Y }\\[2pt]
\DIFadd{$T(2,9)$
  }& \DIFadd{$\sum_{k=-4}^{4}(-1)^k t^k$
  }& \DIFadd{$1+5z^2+6z^4+3z^6+z^8$
  }& \DIFadd{see~\eqref{eq:jones-T29}
  }& \DIFadd{$0$
  }& \DIFadd{C }& \DIFadd{$4$ }& \DIFadd{N }& \DIFadd{Y }\\[2pt]
\DIFadd{$3_1 \mathbin{\#} 3_1$
  }& \DIFadd{$(t^{-1}-1+t)^2$
  }& \DIFadd{$(1+z^2)^2$
  }& \DIFadd{$(-t^{-4}+t^{-3}+t^{-1})^2$
  }& \DIFadd{$0$
  }& \DIFadd{C }& \DIFadd{$2$ }& \DIFadd{N }& \DIFadd{N }\\[2pt]
\DIFadd{$3_1 \mathbin{\#} 3_1^*$
  }& \DIFadd{$(t^{-1}-1+t)^2$
  }& \DIFadd{$(1+z^2)^2$
  }& \DIFadd{$(-t^{-4}+t^{-3}+t^{-1})(-t^4+t^3+t)$
  }& \DIFadd{$0$
  }& \DIFadd{A }& \DIFadd{$2$ }& \DIFadd{N }& \DIFadd{N }\\[2pt]
\DIFadd{$3_1 \mathbin{\#} 4_1$
  }& \DIFadd{$(t^{-1}-1+t)(-t^{-1}+3-t)$
  }& \DIFadd{$(1+z^2)(1-z^2)$
  }& \DIFadd{$(-t^{-4}+t^{-3}+t^{-1})(t^2-t+1-t^{-1}+t^{-2})$
  }& \DIFadd{$2.029$
  }& \DIFadd{C }& \DIFadd{$2$ }& \DIFadd{N }& \DIFadd{N }\\[2pt]
\hline
\end{tabular}}

\bigskip
\noindent\DIFadd{\textbf{Jones polynomials for higher torus knots.}\;
Using the Kauffman bracket / HOMFLY skein relation
$t^{-1}V(L_+) - tV(L_-) = (t^{1/2}-t^{-1/2})V(L_0)$:
}\begin{align}
  \DIFadd{V(T(3,4)) }&\DIFadd{= t^{-8}+t^{-6}-t^{-5}+t^{-4}-t^{-3}+t^{-2},
    \label{eq:jones-T34}}\\[4pt]
  \DIFadd{V(T(2,7)) }&\DIFadd{= -t^{-10}+t^{-9}-t^{-8}+t^{-7}+t^{-5},
    \label{eq:jones-T27}}\\[4pt]
  \DIFadd{V(T(3,5)) }&\DIFadd{= -t^{-16}+t^{-15}-t^{-14}+2t^{-13}-t^{-12}+t^{-11}-t^{-10}+t^{-9},
    \label{eq:jones-T35}}\\[4pt]
  \DIFadd{V(T(2,9)) }&\DIFadd{= -t^{-14}+t^{-13}-t^{-12}+t^{-11}+t^{-9}.
    \label{eq:jones-T29}
}\end{align}

\subsection{\DIFadd{Knot group presentations}}
\DIFadd{\label{subsec:knot-groups}
}

\begin{remark}[Groups of the fundamental knots]
\DIFadd{\label{rem:knot-groups}
The fundamental group $G_K = \pi_1(S^3 \setminus K)$ of each knot in the
particle ansatz has the following presentation.  For torus knots
$T(p,q)$ the group is the \emph{torus knot group}
$\langle a, b \mid a^p = b^q \rangle$.
}

\medskip
{\renewcommand{\arraystretch}{1.3}
\begin{tabular}{lll}
\hline
\DIFadd{\textbf{Knot} }& \DIFadd{\textbf{Generators} }& \DIFadd{\textbf{Relations} }\\
\hline
\DIFadd{$0_1$ }& \DIFadd{$\langle a \rangle$ }& \DIFadd{none }\\
\DIFadd{$2^2_1$ Hopf }& \DIFadd{$\langle a, b \rangle$ }& \DIFadd{$ab = ba$ }\\
\DIFadd{$3_1$ }& \DIFadd{$\langle a, b \rangle$ }& \DIFadd{$a^2 = b^3$ }\\
\DIFadd{$4_1$ }& \DIFadd{$\langle x, y \rangle$ }& \DIFadd{$xyx^{-1}yx = yxy^{-1}xy^{-1}$ }\\
\DIFadd{$5_1 = T(2,5)$ }& \DIFadd{$\langle a, b \rangle$ }& \DIFadd{$a^2 = b^5$ }\\
\DIFadd{$5_2$ }& \DIFadd{$\langle x, y \rangle$ }& \DIFadd{$xy^2x = yx^2y$ }\\
\DIFadd{$T(3,4)$ }& \DIFadd{$\langle a, b \rangle$ }& \DIFadd{$a^3 = b^4$ }\\
\DIFadd{$T(2,7)$ }& \DIFadd{$\langle a, b \rangle$ }& \DIFadd{$a^2 = b^7$ }\\
\DIFadd{$T(3,5)$ }& \DIFadd{$\langle a, b \rangle$ }& \DIFadd{$a^3 = b^5$ }\\
\DIFadd{$T(2,9)$ }& \DIFadd{$\langle a, b \rangle$ }& \DIFadd{$a^2 = b^9$ }\\
\DIFadd{$3_1\mathbin{\#}3_1$ }& \DIFadd{$\langle a,b,c,d \rangle$
  }& \DIFadd{$a^2=b^3,\;c^2=d^3,\;a=c_{\mathrm{merid}}$ }\\
\DIFadd{$3_1\mathbin{\#}3_1^*$ }& \DIFadd{$\langle a,b,c,d \rangle$
  }& \DIFadd{$a^2=b^3,\;c^2=d^{-3},\;a=c_{\mathrm{merid}}$ }\\
\hline
\end{tabular}}

\medskip\noindent
\DIFadd{The figure-eight group $G_{4_1}$ has the equivalent Fox presentation
$\langle x,t \mid txt^{-1}xt = xtx^{-1}tx^{-1}\rangle$ and is the only
hyperbolic knot group in the table, corresponding to the non-zero
hyperbolic volume.  For torus knot groups, the center is generated by
the element $a^p = b^q$ (the longitude), which maps to the Seifert fiber
in the Seifert fibration of the torus knot complement.
}\end{remark}



\section{\DIFadd{Knot Braiding and Skein Relations}}
\DIFadd{\label{sec:knot-braiding}
}

\subsection{\DIFadd{HOMFLY-PT skein relation}}
\DIFadd{\label{subsec:homfly-skein}
}

\begin{definition}[HOMFLY-PT polynomial]
\DIFadd{\label{def:homfly}
The \emph{HOMFLY-PT polynomial} $P(K; \ell, m) \in \mathbb{Z}[\ell^{\pm 1}, m^{\pm 1}]$
of an oriented knot or link $K$ is characterised by the skein relation
}\begin{equation}\DIFadd{\label{eq:homfly-skein}
  \ell\,P(L_+) + \ell^{-1}\,P(L_-) + m\,P(L_0) = 0,
}\end{equation}
\DIFadd{together with the normalisation $P(\text{unknot}) = 1$.  Here
$L_+$, $L_-$, $L_0$ are three diagrams identical except near one
crossing, where $L_+$ has a positive crossing, $L_-$ a negative
crossing, and $L_0$ the oriented smoothing.
}

\medskip\noindent
\DIFadd{Specialisations recover classical invariants:
}\begin{align}
  &\DIFadd{\text{Jones:} }&
    \DIFadd{V(K)(t) }&\DIFadd{= P(K;\, t^{-1},\, t^{1/2}-t^{-1/2}),
    \label{eq:jones-specialisation}}\\[4pt]
  &\DIFadd{\text{Alexander:} }&
    \DIFadd{\Delta(K)(t) }&\DIFadd{= P(K;\, i,\, i(t^{1/2}-t^{-1/2})),
    \label{eq:alexander-specialisation}}\\[4pt]
  &\DIFadd{\text{Conway:} }&
    \DIFadd{\nabla(K)(z) }&\DIFadd{= P(K;\, i,\, -iz).
    \label{eq:conway-specialisation}
}\end{align}
\end{definition}

\begin{definition}[Kauffman bracket skein relation]
\DIFadd{\label{def:kauffman-bracket}
The \emph{Kauffman bracket} $\langle D \rangle \in \mathbb{Z}[A^{\pm 1}]$ of
an unoriented link diagram $D$ satisfies
}\begin{equation}\DIFadd{\label{eq:kauffman-bracket}
  \langle L_+ \rangle
  = A\,\langle L_{0} \rangle
  + A^{-1}\langle L_{\infty} \rangle,
  \qquad
  \langle D \sqcup \bigcirc \rangle = (-A^2 - A^{-2})\langle D \rangle,
  \qquad
  \langle \bigcirc \rangle = 1,
}\end{equation}
\DIFadd{where $L_0$ and $L_\infty$ are the two smoothings of a positive crossing $L_+$
(the $A$-smoothing and $A^{-1}$-smoothing respectively).
The Jones polynomial is recovered by
$V(K)(t) = (-1)^{w(D)}A^{-3w(D)}\langle D\rangle\big|_{A=t^{-1/4}}$,
where $w(D)$ is the writhe of a diagram $D$ for $K$.
}\end{definition}

\subsection{\DIFadd{Skein triples and braid words}}
\DIFadd{\label{subsec:skein-triples}
}

\begin{remark}[Skein triples for fundamental knots]
\DIFadd{\label{rem:skein-triples}
Each knot in the particle ansatz fits into a skein triple
$(L_+, L_-, L_0)$ showing how to reduce it to simpler knots via the
HOMFLY-PT relation~\eqref{eq:homfly-skein}.  The table below lists the
skein triples and braid word presentations for each knot.  All knots
are oriented; the braid word gives the closure $\hat\beta$ realising the
knot as a closed braid.
}

\medskip
{\renewcommand{\arraystretch}{1.4}
\noindent\begin{tabular}{llll}
\hline
\DIFadd{\textbf{Knot} }& \DIFadd{\textbf{Braid word} }& \DIFadd{\textbf{$L_+$} }& \DIFadd{\textbf{$L_0$} }\\
\hline
\DIFadd{$0_1$ unknot
  }& \DIFadd{$1$ (empty)
  }& \DIFadd{---
  }& \DIFadd{--- }\\[2pt]
\DIFadd{$2^2_1$ Hopf
  }& \DIFadd{$\sigma_1^2$
  }& \DIFadd{$2^2_1$
  }& \DIFadd{$0_1 \sqcup 0_1$ }\\[2pt]
\DIFadd{$3_1$ trefoil
  }& \DIFadd{$\sigma_1^3$ ($B_2$-braid)
  }& \DIFadd{$3_1$
  }& \DIFadd{$2^2_1$ }\\[2pt]
\DIFadd{$4_1$ fig.-eight
  }& \DIFadd{$\sigma_1\sigma_2^{-1}\sigma_1\sigma_2^{-1}$ ($B_3$)
  }& \DIFadd{$3_1^*$
  }& \DIFadd{$2^2_1$ }\\[2pt]
\DIFadd{$5_1 = T(2,5)$
  }& \DIFadd{$\sigma_1^5$
  }& \DIFadd{$5_1$
  }& \DIFadd{$T(2,4)$ }\\[2pt]
\DIFadd{$5_2$
  }& \DIFadd{$\sigma_1^2\sigma_2^{-1}\sigma_1^2\sigma_2^{-1}$ ($B_3$)
  }& \DIFadd{$4_1$
  }& \DIFadd{$\text{link}$ }\\[2pt]
\DIFadd{$T(3,4)$
  }& \DIFadd{$(\sigma_1\sigma_2)^4$ ($B_3$)
  }& \DIFadd{$T(3,4)$
  }& \DIFadd{$T(3,3)$ }\\[2pt]
\DIFadd{$T(2,7)$
  }& \DIFadd{$\sigma_1^7$
  }& \DIFadd{$T(2,7)$
  }& \DIFadd{$T(2,6)$ }\\[2pt]
\DIFadd{$T(3,5)$
  }& \DIFadd{$(\sigma_1\sigma_2)^5$ ($B_3$)
  }& \DIFadd{$T(3,5)$
  }& \DIFadd{$T(3,4)$ }\\[2pt]
\DIFadd{$T(2,9)$
  }& \DIFadd{$\sigma_1^9$
  }& \DIFadd{$T(2,9)$
  }& \DIFadd{$T(2,8)$ }\\[2pt]
\hline
\end{tabular}}
\end{remark}

\begin{remark}[Explicit skein computations]
\DIFadd{\label{rem:skein-computations}
Applying~\eqref{eq:homfly-skein} at the listed crossing for each knot:
}

\medskip
\noindent\DIFadd{\textbf{Trefoil $3_1$ from unknot via Hopf link.}\;
The positive trefoil $3_1^+$ satisfies
}\begin{equation}\DIFadd{\label{eq:skein-trefoil}
  \ell\,P(3_1^+) + \ell^{-1}\,P(0_1) + m\,P(2^2_1) = 0.
}\end{equation}
\DIFadd{With $P(0_1)=1$ and $P(2^2_1) = -\ell(m^{-1}+m) - \ell^{-1}m^{-1}$
(two-component unlink skein), one recovers
}\[
  \DIFadd{P(3_1^+; \ell, m) = -\ell^{-4}+2\ell^{-2}+\ell^{-2}m^2.
}\]

\noindent\DIFadd{\textbf{Figure-eight $4_1$ as amphichiral fixed point.}\;
The skein relation at the central crossing gives
}\begin{equation}\DIFadd{\label{eq:skein-fig8}
  \ell\,P(L_+) + \ell^{-1}\,P(L_-) + m\,P(L_0) = 0,
}\end{equation}
\DIFadd{where $L_+ = 3_1^*$ (negative trefoil) and $L_- = 3_1^+$ (positive trefoil).
Since $P(4_1;\ell,m) = P(4_1;\ell^{-1},m)$ (amphichirality),
the equation constrains $P(4_1)$ to be symmetric in $\ell\leftrightarrow\ell^{-1}$.
}

\noindent\DIFadd{\textbf{Torus knot $T(2,2k+1)$ recursion.}\;
The skein relation on the last positive crossing of $T(2,2k+1)$ gives
}\begin{equation}\DIFadd{\label{eq:skein-torus-2n}
  \ell\,P(T(2,2k+1)) + \ell^{-1}\,P(T(2,2k-1)) + m\,P(T(2,2k)) = 0,
}\end{equation}
\DIFadd{with $P(T(2,1))=P(0_1)=1$ and $P(T(2,2))=P(2^2_1)$.
This determines $P(T(2,2k+1))$ inductively for all $k \ge 1$.
}

\noindent\DIFadd{\textbf{Torus knot $T(3,n)$ recursion.}\;
Writing $T(3,n)$ as the closure of $(\sigma_1\sigma_2)^n$ in $B_3$, the skein relation gives
}\begin{equation}\DIFadd{\label{eq:skein-torus-3n}
  \ell\,P(T(3,n)) + \ell^{-1}\,P(T(3,n-2)) + m\,P(L_n) = 0,
}\end{equation}
\DIFadd{where $L_n$ is an intermediate link whose HOMFLY polynomial is computed
from lower-crossing diagrams.
}\end{remark}

\begin{remark}[Jones skein relation for each knot]
\DIFadd{\label{rem:jones-skein-per-knot}
Specialising $\ell = t^{-1}$, $m = t^{1/2}-t^{-1/2}$ in
\eqref{eq:homfly-skein} gives
$t^{-1}V(L_+) - tV(L_-) = (t^{1/2}-t^{-1/2})V(L_0)$.
The table below shows the skein relation at one distinguished crossing
for each knot, together with the resulting Jones computation:
}

\medskip
{\renewcommand{\arraystretch}{1.4}
\noindent\begin{tabular}{lll}
\hline
\DIFadd{\textbf{Knot $K$} }& \DIFadd{\textbf{Jones skein identity} }& \DIFadd{\textbf{Result} }\\
\hline
\DIFadd{$0_1$
  }& \DIFadd{$V(0_1) = 1$ (base case)
  }& \DIFadd{$1$ }\\[2pt]
\DIFadd{$2^2_1$
  }& \DIFadd{$t^{-1}V(2^2_1) - t\,V(0_1\sqcup 0_1) = (t^{1/2}-t^{-1/2})V(0_1)$
  }& \DIFadd{$-t^{-1/2}-t^{-5/2}$ }\\[2pt]
\DIFadd{$3_1$
  }& \DIFadd{$t^{-1}V(3_1) - t\,V(0_1) = (t^{1/2}-t^{-1/2})V(2^2_1)$
  }& \DIFadd{$-t^{-4}+t^{-3}+t^{-1}$ }\\[2pt]
\DIFadd{$4_1$
  }& \DIFadd{$t^{-1}V(3_1^*) - t\,V(3_1) = (t^{1/2}-t^{-1/2})V(2^2_1)$
  }& \DIFadd{$t^2-t+1-t^{-1}+t^{-2}$ }\\[2pt]
\DIFadd{$5_1$
  }& \DIFadd{$t^{-1}V(5_1)-t\,V(3_1)=(t^{1/2}-t^{-1/2})V(T(2,4))$
  }& \DIFadd{see~\eqref{eq:jones-T35} }\\[2pt]
\DIFadd{$5_2$
  }& \DIFadd{$t^{-1}V(4_1)-t\,V(5_2^*)=(t^{1/2}-t^{-1/2})V(L)$
  }& \DIFadd{$-t^6+t^5+t^3$ }\\[2pt]
\DIFadd{$T(3,4)$
  }& \DIFadd{$t^{-1}V(T(3,4))-t\,V(T(3,2))=(t^{1/2}-t^{-1/2})V(T(3,3))$
  }& \DIFadd{see~\eqref{eq:jones-T34} }\\[2pt]
\DIFadd{$T(2,7)$
  }& \DIFadd{$t^{-1}V(T(2,7))-t\,V(T(2,5))=(t^{1/2}-t^{-1/2})V(T(2,6))$
  }& \DIFadd{see~\eqref{eq:jones-T27} }\\[2pt]
\DIFadd{$T(3,5)$
  }& \DIFadd{$t^{-1}V(T(3,5))-t\,V(T(3,3))=(t^{1/2}-t^{-1/2})V(T(3,4))$
  }& \DIFadd{see~\eqref{eq:jones-T35} }\\[2pt]
\DIFadd{$T(2,9)$
  }& \DIFadd{$t^{-1}V(T(2,9))-t\,V(T(2,7))=(t^{1/2}-t^{-1/2})V(T(2,8))$
  }& \DIFadd{see~\eqref{eq:jones-T29} }\\[2pt]
\hline
\end{tabular}}
\end{remark}

\begin{remark}[$q$-deformed skein relation in the QOU]
\DIFadd{\label{rem:q-deformed-skein}
In the QOU framework, the standard Jones skein relation is replaced by
a $q$-deformed version incorporating the substrate parameter
$q(p)$ (Definition~\ref{def:substrate-parameter}):
}\begin{equation}\DIFadd{\label{eq:q-skein}
  q^{-1}\,V_q(L_+) \;-\; q\,V_q(L_-) \;=\;
  (q^{1/2}-q^{-1/2})\,V_q(L_0),
}\end{equation}
\DIFadd{where $V_q(K)$ is the $q$-deformed Jones polynomial evaluated at the
substrate parameter $q(p)$ of the QOU point $p \in B_3$.  At
$q\to 1$ (classical limit, heat death) equation~\eqref{eq:q-skein}
reduces to the ordinary Jones skein relation.  At $q\to\infty$
(deep-quantum, Big Bang) the $q$-deformed $R$-matrix of
$\mathcal{U}_q(\mathfrak{su}(2))$
(Remark~\ref{rem:r-matrix}) furnishes the twisted exchange operator,
and the braiding becomes maximal.
}

\medskip\noindent
\DIFadd{The $q$-deformation parameter couples the knot invariant to the
topology of the Bring hierarchy: the substrate parameter $q(p)$ is
maximal at the six branch points (where $B_3$ and $B_2$ have their
most complex fiber structure), and the $q$-deformed skein relation
therefore encodes the transition between hierarchy levels as a
deformation of the classical braid relations.
}\end{remark}
 \DIFaddend \chapter{The Descartes Universe}
\label{chap:descartes}



\begin{definition}[Descartes condition]
\label{def:descartes-condition}
A $\mathbb{C}$ quantum observable universe $(M, G, P, E)$
(Example~\ref{def:c-quantum-observable-universe}) satisfies the
\emph{Descartes condition} if the Euler characteristic of the bulk bundle
is
\[
  \chi(E) \;=\; -6.
\]
A $\mathbb{C}$ quantum observable universe satisfying the Descartes condition
is called a \emph{Descartes universe}.

\medskip
\begin{center}
\itshape
``Je pense, donc je suis.''\\
--- Ren\'e Descartes
\end{center}
\end{definition}



\begin{remark}[The Descartes condition on Calabi--Yau manifolds]
\label{rem:descartes-calabi-yau}
Let $(M, G, P, E)$ be a Calabi--Yau quantum observable universe
(Example~\ref{ex:calabi-yau}).  The Descartes condition $\chi(E) = -6$
constrains the Hodge numbers of the Calabi--Yau manifold $M$: the
alternating sum of Betti numbers of the bulk bundle $E$ must equal $-6$.
In particular, only Calabi--Yau manifolds whose topology admits a principal
$\mathrm{SU}(2)$-bundle with this Euler characteristic can realise the
Descartes universe.
\end{remark}

\begin{remark}[The Descartes condition on ALE spaces]
\label{rem:descartes-ale}
For an ALE quantum observable universe (Example~\ref{ex:ale}), the Descartes
condition $\chi(E) = -6$ selects specific finite subgroups
$\Gamma \subset \mathrm{SU}(2)$ in the quotient singularity
$\mathbb{C}^2 / \Gamma$.  The Euler characteristic of the resolved space
and its principal bundle are tied to the order and representation theory of
$\Gamma$, so the Descartes condition imposes arithmetic constraints on the
McKay quiver of~$\Gamma$.
\end{remark}

\begin{remark}[The Descartes condition on affine space]
\label{rem:descartes-affine}
For the affine quantum observable universe $M = \mathbb{C}^n$,
the Descartes condition $\chi(E) = -6$ forces the
principal $\mathrm{SU}(2)$-bundle $E$ to carry non-trivial topology despite
the contractibility of the base.  In the two-dimensional case $M = \mathbb{C}^2$,
the $q$-harmonic forms on this bundle
(Definition~\ref{def:q-harmonic-form}) determine the spectrum of stable
excitations.  In the regime $q(p) \to 1$ (near-classicality), the zero-modes
of $\Delta_q$ reproduce the thermal fluctuation spectrum of the Cosmic
Microwave Background: the Hodge numbers of $E$ govern the multiplicity of
angular modes, and the negative Euler characteristic forces a non-trivial
distribution of these modes across the CR manifold $M$.
\end{remark}



\begin{remark}[Descent to one complex dimension]
\label{rem:descent-one-dim}
By the descent theorem (Proposition~\ref{prop:jet-prolongation}), a
two-complex-dimensional Descartes universe $(M, G, P, E)$ with $M = \mathbb{C}^2$
descends to a one-complex-dimensional variety $M_0$.  The stability of the
Euler characteristic under descent and lifting
(Proposition~\ref{prop:euler-char-stability}) ensures that the descended
universe retains $\chi(E) = -6$.  We may therefore study the Descartes
universe in the setting of one-complex-dimensional varieties.
The descent may continue further to dimension zero
(Remark~\ref{rem:descent-reduction-zero}), where the entire quantum
observable universe is captured by fibre data at a single point.
\end{remark}



\begin{proposition}[Genus of the one-complex-dimensional descent]
\label{prop:genus-maximal-descent}
Let $(M, G, P, E)$ be a Descartes universe that descends to a compact
one-complex-dimensional variety $M_0$.  Then $M_0$ is a Riemann surface of
genus $g = 4$.
\end{proposition}

\begin{proof}
For a compact Riemann surface of genus $g$, the Euler characteristic is
$\chi = 2 - 2g$.  The Descartes condition requires $\chi(E) = -6$.  On
the one-complex-dimensional descent $M_0$, the stability of the Euler characteristic
(Proposition~\ref{prop:euler-char-stability}) gives $\chi(M_0) = -6$,
whence $2 - 2g = -6$ and $g = 4$.
\end{proof}



\begin{remark}[Electrons, muons and the Reeb decomposition]
\label{rem:descartes-leptons}
The Reeb decomposition (Definition~\ref{def:reeb-decomposition}) applied to
a Descartes universe yields two classes of fermionic excitations
(Definition~\ref{def:fermion}).  The degree of the vertical component
$\alpha_V$ distinguishes them: forms with $\deg\alpha_V = 1$ correspond to
electrons, while forms with $\deg\alpha_V = 2$ correspond to muons.  More
generally, each additional unit of vertical degree is identified with a
successive lepton generation.  The constraint $\chi(E) = -6$ bounds the
number of such generations, since the Hodge decomposition of $E$ admits only
finitely many independent vertical harmonic forms.  This
generation-dependent structure leads to testable predictions, including a
$q$-correction to the proton charge radius that differs between muonic and
electronic hydrogen (Section~\ref{ssec:proton-radius}).
\end{remark}

\begin{remark}[Chirality as Reeb orientation]
\label{rem:descartes-chirality}
The Reeb vector field $\xi$ on the CR manifold $M$ carries a natural
orientation: the sign of the contact form $\theta(\xi) = 1$ singles out a
preferred direction along the vertical distribution.  Left-handed and
right-handed chiralities correspond to the two orientations of $\xi$ relative
to the complex structure of $M$.  For fermions
(Definition~\ref{def:fermion}), the vertical component $\alpha_V$ transforms
with a definite sign under reversal of $\xi$, so chirality is encoded in the
$\pm$-eigenspaces of the Reeb involution $\xi \mapsto -\xi$.  Bosons
(Definition~\ref{def:boson}), being purely horizontal, are chirality-blind ---
consistent with the observed parity invariance of the electromagnetic and
strong interactions.
\end{remark}



\begin{remark}[Light cone of an affine chart]
\label{rem:light-cone-affine}
Consider an open subset $U \subset M_0$ representing a portion of the
descended universe from the big bang to heat death.  An \emph{affine chart}
on $U$ is a local coordinate patch in which the metric on the base manifold
$M_0$ defines a causal structure.  The \emph{light cone} of the affine chart
is the boundary of the causal future in $U$: it is determined by the metric
on $M_0$ and has the property that no singularity of the ambient space is
visible from within the chart.  In this way, the affine chart provides a
non-singular window onto a portion of the Descartes universe.
\end{remark}



\begin{remark}[Chirality and trefoil subgroups]
\label{rem:chirality-trefoil}
The chirality of the observable universe --- manifest in the three colour
charges of quantum chromodynamics --- is encoded in the symmetry group of the
descended surface $M_0$.  Proposition~\ref{prop:genus-maximal-descent}
identifies $M_0$ as a genus-$4$ Riemann surface.  We require a genus-$4$
surface whose automorphism group contains three subgroups isomorphic to $S_3$
--- quotients of the trefoil knot group
$\pi_1(S^3 \setminus K) \cong \langle x, y \mid x^2 = y^3 \rangle$ ---
corresponding to the three colour charges (red, green, blue).  These $S_3$
subgroups generate a subgroup of the monodromy group, and the remaining
generators of the full automorphism group act by conjugation, permuting the
three $S_3$ subgroups among themselves.
\end{remark}

\begin{remark}[Bring's curve]
\label{rem:brings-curve}
Bring's curve is the unique Riemann surface of genus $4$ with the maximal
automorphism group $S_5$ (order~$120$) for its genus.  It is defined by
the intersection of three symmetric hypersurfaces in $\mathbb{P}^4$:
\[
  \sum_{i=1}^{5} z_i = 0,
  \qquad
  \sum_{i=1}^{5} z_i^2 = 0,
  \qquad
  \sum_{i=1}^{5} z_i^3 = 0.
\]
The symmetric group $S_5$ acts by permuting the five homogeneous coordinates,
and this action realises the monodromy structure described in
Remark~\ref{rem:chirality-trefoil}: the group $S_5$ contains three copies of
$S_3$ (quotients of the trefoil knot group) corresponding to the three
colour charges, and elements outside these subgroups act by conjugation,
permuting the three copies among themselves.
\end{remark}



\begin{theorem}[Descartes theorem]
\label{thm:descartes}
Let $\mathcal{C}_{-6}$ denote the category of $\mathbb{C}$ quantum
observable universes satisfying the Descartes condition $\chi(E) = -6$,
with morphisms given by QOU maps that preserve the Euler characteristic.
Then Bring's curve $\Sigma_4$ (Remark~\ref{rem:brings-curve}) is a terminal
object in $\mathcal{C}_{-6}$: for every Descartes universe $(M, G, P, E)$ in
$\mathcal{C}_{-6}$, there exists a unique morphism
$(M, G, P, E) \to (\Sigma_4, G, P_{\Sigma_4}, E_{\Sigma_4})$ in $\mathcal{C}_{-6}$.
\end{theorem}

\begin{proof}
\end{proof}

\begin{remark}[Universality for a given Euler characteristic]
\label{rem:universality-general}
The Descartes theorem (Theorem~\ref{thm:descartes}) is expected to be a
special case of a more general result: for any admissible Euler
characteristic $\chi$, the category $\mathcal{C}_{\chi}$ of $\mathbb{C}$
quantum observable universes with $\chi(E) = \chi$ should admit a terminal
object determined by the unique maximally symmetric Riemann surface of the
corresponding genus.  The Descartes case $\chi = -6$ (genus~$4$) is
distinguished by Bring's curve, but analogous universal objects are expected
at other values of~$\chi$.
\end{remark}

\begin{remark}[The notion of an observation]
\label{rem:observation}
Given a fermion (Definition~\ref{def:fermion}), an \emph{observation}
is a continuous open map from the spacetime base $\mathbb{CP}^1$ to
the bundle of endomorphisms
$\operatorname{End}(E)\to\mathbb{CP}^1$, encoding the coordinate
change in torsion energy of the bundle~$E$.  Because the torsion
energy is modified by the act of measurement, observing the universe
changes it: every observation perturbs the substrate parameter~$q(p)$
(Definition~\ref{def:substrate-parameter}) and hence shifts the
non-commutative geometry of the Bring surface
(Section~\ref{sec:quantum-brings}).

\smallskip
\begin{quote}
\itshape ``Observing the universe changes it.''
\end{quote}
\end{remark}

 \chapter{\texorpdfstring{$\pp$}{p}-Adic Quantum Observable Universe}
\label{chap:p-adic}



\section{Key Results in Categorical Language}
\label{sec:categorical-reformulation}

\begin{itemize}
  \item Rephrase the QOU definitions
        (Chapter~\ref{chap:qou}) in the language of monoidal categories
        and functors.
  \item Key theorems to establish in the categorical setting:
    \begin{itemize}
      \item Lifting (Proposition~\ref{prop:jet-prolongation}) as a
            categorical lift/extension
            (Remark~\ref{rem:prolongation-categorical}).
      \item Descent (Subsection~\ref{ssec:descent}) via fibred categories.
      \item Obstructions described in terms of cohomology of the category
            (sheaf cohomology, derived functors).
      \item TODO: show that the Atiyah exact sequence splits if and only
            if the bundle admits a linear connection.
    \end{itemize}
  \item Role of $\mathrm{SU}(2)$ and $\mathbb{CP}$ from the symmetry
        perspective; motivate categorical analogies
        (e.g.\ Tannaka--Krein duality).
\end{itemize}



\section{Universal Enveloping Algebras and Quantum Groups}
\label{sec:enveloping-algebras}

\begin{itemize}
  \item Lusztig integral form of
        $\mathcal{U}_q(\mathfrak{su}(2))$ and its triangular decomposition.
  \item Poincar\'{e}--Birkhoff--Witt (PBW) basis and its
        $q$-deformed analogue.
  \item Universal enveloping algebras and their role in lifting
        (cf.\ \href{https://arxiv.org/abs/math/0409565}{arXiv:math/0409565}).
  \item Heisenberg uncertainty as a $q$-commutator:
        $[\hat{X},\hat{P}]_q = i\hbar_q$.
\end{itemize}



\section{\texorpdfstring{$\pp$}{p}-Adic Number Fields and Universality}
\label{sec:p-adic-fields}

\begin{itemize}
  \item Describe $\mathbb{Q}$, $\mathbb{C}_\pp$, $\mathbb{Q}_\pp$, and
        $\mathbb{F}_\pp$ in terms of their category-theoretic universality
        relationships to each other.
  \item Ostrowski's theorem as a categorical coproduct characterisation.
  \item Connections between $\pp$-adic completions and inverse limits in
        the category of topological rings.
\end{itemize}



\section{\texorpdfstring{$\mathbb{C}_\pp$}{C\_p} Quantum Universe}
\label{sec:cp-que}

\begin{example}[$\mathbb{Q}_\pp$ realisation]
\label{ex:p-adic-realisation}
Let $k = \mathbb{Q}_\pp$ and let $\mathcal{C}$ be the category of
$F$-isocrystals over~$\mathbb{Q}_\pp$ (modules over the Dieudonn\'e ring
$\mathbb{Q}_\pp\langle F, V \rangle / (FV - \pp)$ that are finite-dimensional
over~$\mathbb{Q}_\pp$).
This is a $\mathbb{Q}_\pp$-linear rigid monoidal category.
The quantum group $\mathcal{O}_q(\mathrm{SU}(2))$ base-changed
to~$\mathbb{Q}_\pp$ provides a Frobenius--Hopf object~$A$, where the
Frobenius form~$\lambda$ arises from the $\pp$-adic Haar measure.
The Frobenius endomorphism~$\varphi$ on $\mathcal{C}$ plays the
r\^ole of the Levi-form contraction: the substrate parameter is
$q = \pp^{\,-v_\pp(\det\varphi)}$ (a $\pp$-adic analogue of
Definition~\ref{def:substrate-parameter}).  The fibre functor~$\omega$
is the forgetful functor to $\mathrm{Vect}_{\mathbb{Q}_\pp}$.
The observability condition
$\mathrm{Hom}(\mathbf{1}, \mathrm{Ad}(A)) = 0$ amounts to the
irreducibility of the associated $F$-crystal, generalising
Lemma~\ref{lem:irreducible-connection} to the $\pp$-adic setting.

\medskip\noindent\emph{Key structures.}
\begin{itemize}
  \item The classical Frobenius integrability theorem constrains
        $q$-torsion leaves to be Legendrian
        (Remark~\ref{rem:q-torsion-existence}).  Over~$\mathbb{Q}_\pp$,
        this is replaced by compatibility of the connection with the
        Frobenius lift: a $\pp$-adic instanton connection is
        ``integrable'' if and only if the associated $F$-crystal is
        unit-root.
  \item Crystal cohomology (de Rham--Witt complex) plays the r\^ole of
        de Rham cohomology.
  \item The base-field-independent
        Definitions~\ref{def:quantum-universe}
        and~\ref{def:quantum-observable-universe} provide the categorical
        framework via Tannakian reconstruction.
\end{itemize}
\end{example}



\section{Quotient to Finite Fields}
\label{sec:finite-fields}

\begin{example}[$\mathbb{F}_\pp$ realisation]
\label{ex:finite-field-realisation}
Let $k = \mathbb{F}_\pp$ and let $\mathcal{C} =
\mathrm{Rep}_{\mathbb{F}_\pp}(\mathrm{SU}_q(2))$, the category of
finite-dimensional representations of the quantum group
$\mathrm{SU}_q(2)$ over~$\mathbb{F}_\pp$, where $q \in
\mathbb{F}_\pp^{\times}$.  More generally, when $q$ is a primitive
$\ell$-th root of unity in $\overline{\mathbb{F}}_\pp$ (with $\ell$
coprime to~$\pp$), $\mathcal{C}$ extends to a fusion category over the
extension field $\mathbb{F}_\pp(q)$ with finitely many simple objects.
The Frobenius--Hopf object~$A$ is the regular representation, and
the fibre functor~$\omega$ is the forgetful functor to
$\mathrm{Vect}_{\mathbb{F}_\pp}$.  The non-degeneracy of the
Frobenius form~$\lambda$ follows from the semisimplicity of the fusion
category, and the observability condition
$\mathrm{Hom}(\mathbf{1}, \mathrm{Ad}(A)) = 0$ holds whenever $q \neq 1$
(since the trivial representation does not appear as a summand of the
adjoint representation for generic~$q$).
\end{example}

\begin{itemize}
  \item Quotienting from $\mathbb{Q}_\pp$ to work over $\mathbb{F}_\pp$:
        reduction modulo~$\pp$.
  \item Finite group representations in the QOU context.
  \item Betti numbers over finite fields and comparison with characteristic-zero
        values.
\end{itemize}



\section{Jones Polynomials}
\label{sec:jones-polynomials}

 \chapter{Gauge Field Fluid Dynamics}
\label{chap:gauge-fluid}



\section{Fluid Dynamics of the Brane}
\label{sec:brane-fluid}

\begin{itemize}
  \item Describe the fluid dynamics on the brane whose thickness is
        determined by the substrate parameter~$q$
        (Definition~\ref{def:substrate-parameter}).
  \item TODO: investigate whether the viscosity gradient can be encoded
        by the $q$-Laplacian $\Delta_q$.
  \item TODO: formulate a Navier--Stokes analogue in the $q$-deformed
        gauge-field setting.
\end{itemize}



\section{Geodesics as Stable Vortex Points}
\label{sec:geodesic-vortices}

\begin{itemize}
  \item Geodesics on $M$ as stable fixed points of vortex configurations
        in the brane fluid.
  \item TODO: perform stability analysis via the $q$-deformed energy
        functional.
\end{itemize}



\section{Braid Operators in the Fluid Picture}
\label{sec:braid-fluid}





\section{$q$-Parameter, Planck Limit, and Reynolds Number}
\label{sec:planck-reynolds}

 \chapter{Information-Theoretic Perspective}
\label{chap:information-theory}



\section{Emergent Gravity from Information Theory}
\label{sec:emergent-gravity-info}

\begin{itemize}
  \item Link to Verlinde's entropic gravity programme and related work on
        emergent gravity from an information-theoretic perspective.
  \item Connection to the entropic-gravity signatures discussed in
        Section~\ref{sec:entropic-gravity}.
\end{itemize}



\section{Local Planck Length as Minimal Information Unit}
\label{sec:planck-info-unit}

\begin{itemize}
  \item Local Planck length $\ell_P(p)$ as the minimal quantisable size of an
        information unit at a point $p \in M$.
  \item Relation to the substrate parameter: $\ell_P \sim \hbar_q$.
\end{itemize}



\section{Time as Information Transfer}
\label{sec:time-info}

\begin{itemize}
  \item Time as a measure of information transfer and change.
  \item Time slows at heat death ($q \to 1$) because information transfer
        rate decreases.
  \item Time runs fast at the Big Bang ($q \to \infty$, equivalently $\hbar_q \to 1$) because information
        transfer is high; torsion of $E$ provides the driving mechanism.
\end{itemize}



\section{Big Bang as Initial Information State}
\label{sec:big-bang-info}

\begin{itemize}
  \item At $t = 0$: no information is known about the state of the QOU;
        matter is concentrated with high torsion pressure.
  \item At $t > 0$: the bundle $E$ acts on fermions, information begins to
        propagate, and the brane expands.
\end{itemize}



\section{Planck Units and Computational Mechanics}
\label{sec:planck-computation}



 \chapter{Interpretations}
\label{chap:interpretations}
 \chapter{Experimental Evidence}
\label{chap:experimental}

The QOU framework (Chapter~\ref{chap:qou}) makes predictions that differ from
the standard $\Lambda$CDM cosmology and the Standard Model of particle physics.
This chapter catalogues observable signatures, organised by scale, and maps each
to existing data sets, planned observations, and potential new experiments.



\section{Large-Scale Observations}
\label{sec:large-scale}

\subsection{CMB Anisotropy and $q$-Torsion Imprint}
\label{sec:cmb-torsion}

\begin{remark}[CMB non-uniformity from $q$-torsion]
\label{rem:cmb-q-torsion}
In the Descartes universe (Example~\ref{ex:descartes}), the negative Euler
characteristic $\chi(E) = -6$ forces a non-trivial distribution of angular
modes across the CR manifold (Remark~\ref{rem:descartes-cmb}).  The
$q$-torsion leaves (Definition~\ref{def:q-torsion-leaf}; here ``leaves'' is the
geometric plural of \emph{leaf}) imprint a
characteristic pattern on the CMB that deviates from the statistically
isotropic Gaussian field predicted by $\Lambda$CDM\@.  Specifically, the
$q$-deformed Reeb flow (Definition~\ref{def:q-reeb-flow}) introduces
\emph{torsion} into the temperature and polarisation multipoles, producing
preferred spiral directions in the $B$-mode spectrum.
\end{remark}

\paragraph{Existing data sets.}
\begin{itemize}
  \item \emph{Planck} 2018 full-sky temperature and polarisation maps
        (PR3/PR4).
  \item \emph{WMAP} nine-year data release.
  \item ACT (Atacama Cosmology Telescope) and SPT (South Pole Telescope)
        high-$\ell$ power spectra.
\end{itemize}

\paragraph{Planned observations.}
\begin{itemize}
  \item \emph{LiteBIRD} (JAXA, launch $\sim$2032): full-sky $B$-mode
        polarisation survey sensitive to large-angular-scale torsion
        signatures.
  \item \emph{CMB-S4}: ground-based next-generation CMB survey with
        improved sensitivity to polarisation anomalies.
  \item \emph{Simons Observatory}: intermediate-scale polarimetry bridging
        current and CMB-S4 sensitivities.
\end{itemize}

\paragraph{Potential new experiments.}
\begin{itemize}
  \item A dedicated search for spiral (chiral) patterns in existing CMB
        $B$-mode data, testing for the preferred handedness predicted by the
        Reeb orientation (Remark~\ref{rem:descartes-chirality}).
  \item Cross-correlation of CMB anomalies (e.g.\ the Cold Spot, hemispherical
        asymmetry) with the angular structure of $q$-torsion leaves.
\end{itemize}

\subsection{Gravitational Signatures and Emergent Geometry}
\label{sec:gravity-emergent}

\begin{remark}[Gravity as emergent substrate geometry]
\label{rem:gravity-emergent}
In the QOU framework, general relativity emerges as the substrate geometry of
the CR manifold (Remark~\ref{rem:gr-substrate-geometry}).  The quantum-weighted
stress-energy tensor $T_{\mu\nu}^{(q)}$ introduces $q$-dependent corrections
to the Einstein field equations.  This is consistent with recent work treating
gravity as arising from an information-theoretic context~\cite{carney2025}.
\end{remark}

\paragraph{Existing data sets.}
\begin{itemize}
  \item LIGO/Virgo/KAGRA gravitational wave catalogues (GWTC-1 through
        GWTC-3).
  \item Pulsar Timing Array data sets (NANOGrav 15-year, EPTA, PPTA,
        IPTA).
  \item Solar System ephemeris tests of general relativity (Cassini, MESSENGER).
\end{itemize}

\paragraph{Planned observations.}
\begin{itemize}
  \item \emph{LISA} (ESA, launch $\sim$2035): space-based gravitational
        wave detector sensitive to low-frequency waveform corrections from the
        $q$-deformed field equations.
  \item \emph{Einstein Telescope}: next-generation ground-based detector
        with sensitivity to post-merger ringdown deviations.
\end{itemize}

\paragraph{Potential new experiments.}
\begin{itemize}
  \item Precision tests of the equivalence principle at quantum scales,
        probing whether the substrate parameter $q$
        (Definition~\ref{def:substrate-parameter}) introduces
        composition-dependent deviations.
  \item Tabletop gravity experiments at sub-millimetre scales to search for
        $q$-dependent modifications to the Newtonian potential.
\end{itemize}

\subsection{Large-Scale Structure and Topological Defects}
\label{sec:large-scale-structure}

\begin{remark}[Topological signatures in large-scale structure]
\label{rem:topological-defects}
The $\mathrm{SO}(3)$ holonomy of the bulk bundle
(Example~\ref{def:c-quantum-universe}) and the monodromy fixed points of the
quantum characteristic classes (Definition~\ref{def:quantum-characteristic-classes})
may leave topological imprints on the large-scale distribution of matter --- for
instance, preferred filamentary alignments or voids correlated with the brane
tower structure (Definition~\ref{def:recursive-brane-tower}).
\end{remark}

\paragraph{Existing data sets.}
\begin{itemize}
  \item SDSS (Sloan Digital Sky Survey) galaxy catalogues and BAO
        measurements.
  \item DES (Dark Energy Survey) weak-lensing maps.
  \item \emph{eROSITA} X-ray sky survey cluster catalogues.
\end{itemize}

\paragraph{Planned observations.}
\begin{itemize}
  \item \emph{Euclid} (ESA): wide-field survey of galaxy shapes and
        redshifts for weak-lensing tomography.
  \item \emph{Vera C.\ Rubin Observatory} (LSST): deep optical survey
        enabling statistical tests of filament chirality.
  \item \emph{Roman Space Telescope} (NASA): high-redshift galaxy survey.
\end{itemize}

\paragraph{Potential new experiments.}
\begin{itemize}
  \item Statistical analysis of cosmic-web filament handedness as a probe
        of global $q$-torsion.
  \item Void-shape analysis to detect signatures of monodromy fixed points
        in the brane tower.
\end{itemize}



\section{Small-Scale Observations}
\label{sec:small-scale}

\subsection{Mass Equation Predictions}
\label{sec:mass-predictions}

\begin{remark}[Theoretical mass ratios from the mass identity]
\label{rem:mass-ratio-predictions}
The mass identity (Definition~\ref{def:mass-identity}) expresses inertial mass
as a line integral of $\lVert\nabla_A q\rVert \cdot \lVert L \wedge F^+\rVert$
around a vortex filament.  For the Descartes universe with $\chi(E) = -6$, the
topological data of the bundle constrain the allowed vortex configurations and
hence the mass spectrum.  The electron--muon mass ratio, for instance, is
governed by the ratio of vertical degrees in the Reeb decomposition
(Remark~\ref{rem:descartes-leptons}).  Evaluating the mass identity at
successive levels of the brane tower
(Definition~\ref{def:recursive-brane-tower}) may yield theoretical predictions
for lepton and hadron mass ratios.
\end{remark}

\paragraph{Existing data sets.}
\begin{itemize}
  \item CODATA recommended values of fundamental constants (electron, muon,
        tau masses; fine-structure constant).
  \item Particle Data Group (PDG) Review of Particle Physics.
  \item Precision atomic spectroscopy measurements (hydrogen, muonium,
        positronium).
\end{itemize}

\paragraph{Planned observations.}
\begin{itemize}
  \item Muon $g-2$ experiment at Fermilab (E989): precision measurement of
        the anomalous magnetic moment.
  \item COMET and Mu2e experiments: searches for charged lepton flavour
        violation that would constrain inter-generation mixing in the Reeb
        decomposition.
\end{itemize}

\paragraph{Potential new experiments.}
\begin{itemize}
  \item High-precision comparisons of the electron--muon and muon--tau mass
        ratios against values computed from the mass identity with specific
        $q$ and $\chi(E)$.
  \item Lattice calculations of QOU vortex masses to compare with QCD
        lattice results.
\end{itemize}

\subsection{Quantum Hall Effect of Photons}
\label{sec:quantum-hall-photons}

\begin{remark}[Photonic quantum Hall effect and $q$-deformed braiding]
\label{rem:photonic-hall}
Recent experiments have demonstrated a quantized Hall drift in photonic Chern
insulators~\cite{chenier2026}.  In the QOU framework, bosons are purely
horizontal $q$-harmonic forms (Definition~\ref{def:boson}), and the Hall
conductance arises from the Chern number of the bulk bundle restricted to the
horizontal distribution.  The $q$-deformation modifies the quantisation
condition, predicting fractional corrections of order $1 - q^{-1}$ to the Hall
conductance.
\end{remark}

\paragraph{Existing data sets.}
\begin{itemize}
  \item Frequency-encoded photonic Chern insulator measurements of
        Ch\'enier et al.~\cite{chenier2026}.
  \item Many-body interference data in kagome lattices~\cite{guo2025}.
  \item Chiral phonon experiments in rare-earth halides~\cite{luo2023}.
\end{itemize}

\paragraph{Planned observations.}
\begin{itemize}
  \item Extension of photonic Chern insulator experiments to tunable
        $q$-regimes via lattice-geometry variation.
  \item Fractional quantum Hall experiments in synthetic photonic systems
        with controllable topology.
\end{itemize}

\paragraph{Potential new experiments.}
\begin{itemize}
  \item Measurement of Hall conductance in photonic lattices at varying
        effective $q$ parameters, seeking the $1 - q^{-1}$ correction predicted
        by the quantum Planck constant
        (Definition~\ref{def:quantum-planck-constant}).
  \item Cold-atom realisations of $\mathrm{SU}(2)$ gauge fields in optical
        lattices, testing the $q$-deformed braiding statistics of
        Chapter~\ref{chap:path-integrals-braiding}.
\end{itemize}

\subsection{Lepton Anomalies and Reeb-Decomposition Tests}
\label{sec:lepton-anomalies}

\begin{remark}[Electron--muon discrepancies]
\label{rem:lepton-discrepancies}
The QOU framework predicts that lepton generations are distinguished by the
vertical degree of the Reeb decomposition
(Remark~\ref{rem:descartes-leptons}).  Discrepancies between electron and muon
interactions --- such as the anomalous magnetic moment tensions and lepton
flavour universality violations observed at the $B$-factories --- may reflect
the distinct Reeb-vertical structure of each generation rather than new
heavy-particle loops.
\end{remark}

\paragraph{Existing data sets.}
\begin{itemize}
  \item Muon $g-2$ results from BNL E821 and Fermilab E989.
  \item LHCb $R_K$ and $R_{K^*}$ measurements of lepton flavour
        universality.
  \item Belle~II early data on $B$-meson decays.
\end{itemize}

\paragraph{Planned observations.}
\begin{itemize}
  \item Full Belle~II data set ($\sim$50\,ab$^{-1}$ integrated luminosity):
        precision tests of
        lepton universality in $\tau$ decays.
  \item HL-LHC (High-Luminosity LHC): extended reach for rare lepton
        flavour violating processes.
\end{itemize}

\paragraph{Potential new experiments.}
\begin{itemize}
  \item Electron--muon scattering experiments (e.g.\ MUonE at CERN) to
        directly test the $q$-dependent corrections to the electromagnetic
        vertex.
\end{itemize}

\subsection{Generation-Dependent Couplings and the Proton Charge Radius}
\label{ssec:proton-radius}

\begin{remark}[The Proton Radius Puzzle]
\label{rem:proton-radius-puzzle}
The ``Proton Radius Puzzle'' --- the discrepancy between the proton charge
radius measured in muonic hydrogen and in electronic hydrogen --- remains a
key area of study.  In the QOU framework, lepton generations are distinguished
by the vertical degree of the Reeb decomposition
(Remark~\ref{rem:descartes-leptons}); each generation couples to the substrate
parameter $q$ (Definition~\ref{def:substrate-parameter}) through a different
vertical harmonic form.  If this generation-dependent coupling is physical, the
effective electromagnetic interaction between the proton and the orbiting lepton
acquires a $q$-correction that depends on the lepton generation.  The measured
charge radius is therefore expected to differ between electronic and muonic
hydrogen by an amount of order $\hbar_q$
(Definition~\ref{def:quantum-planck-constant}), providing a direct test of
generation-dependent substrate couplings.
\end{remark}

\paragraph{Existing data sets.}
\begin{itemize}
  \item Muonic hydrogen Lamb shift measurements (CREMA collaboration, PSI).
  \item Electronic hydrogen spectroscopy (CODATA recommended value).
  \item Electron--proton scattering cross-section data (PRad at JLab).
\end{itemize}

\paragraph{Planned observations.}
\begin{itemize}
  \item MUSE (MUon proton Scattering Experiment, PSI): simultaneous
        electron--proton and muon--proton scattering under identical
        conditions.
  \item High-precision muonic deuterium and muonic helium spectroscopy
        to extend the puzzle to heavier nuclei.
\end{itemize}

\paragraph{Potential new experiments.}
\begin{itemize}
  \item Precision comparison of the proton charge radius extracted from
        muonic vs.\ electronic hydrogen, interpreted through the
        generation-dependent $q$-correction predicted by the Reeb
        decomposition.
  \item Tau-hydrogen spectroscopy (if feasible at future facilities) to
        test the third-generation coupling to the substrate parameter.
\end{itemize}

\subsection{Informational and Entropic Gravity Tests}
\label{sec:entropic-gravity}

\begin{remark}[Entropic forces and the substrate parameter]
\label{rem:entropic-gravity}
The substrate parameter $q$ (Definition~\ref{def:substrate-parameter}) encodes
the coupling between the Levi form and the gauge curvature.  In the
near-classical regime $q \to 1$, gravity emerges as a geometric back-reaction
(Remark~\ref{rem:gr-substrate-geometry}).  Recent work on the quantum mechanics
of entropic forces~\cite{carney2025} provides a complementary perspective: the
entropic interpretation may correspond to the thermodynamic limit of the
$q$-deformed partition function on the brane tower.
\end{remark}

\paragraph{Existing data sets.}
\begin{itemize}
  \item Cavendish-type torsion balance experiments (e.g.\ E\"ot-Wash group).
  \item Casimir force measurements at micrometre scales.
  \item Neutron interferometry (gravity resonance spectroscopy).
\end{itemize}

\paragraph{Planned observations.}
\begin{itemize}
  \item MICROSCOPE follow-up missions: improved equivalence principle
        tests in space.
  \item Atom interferometry experiments (MAGIS, AION, ZAIGA) probing
        gravitational decoherence.
\end{itemize}

\paragraph{Potential new experiments.}
\begin{itemize}
  \item Entanglement-assisted gravity experiments testing whether
        gravitational interactions preserve quantum coherence as predicted
        by the $q$-deformed field equations.
  \item High-precision measurements of Casimir forces in geometries
        designed to isolate contributions from $q$-deformed vacuum
        fluctuations.
\end{itemize}

\subsection{Crystal Lattice Resonance and Substrate Telescopes}
\label{sec:crystal-resonance}

\begin{remark}[Quantum resonance in crystal lattices]
\label{rem:crystal-resonance}
A periodic crystal lattice provides a macroscopic, coherent probe of the
substrate parameter $q$ (Definition~\ref{def:substrate-parameter}).  The
discrete translational symmetry of the lattice selects a set of reciprocal
vectors that couple to the horizontal distribution of the CR manifold
$M$.  Where the lattice spacing is commensurate with the characteristic scale
of a $q$-torsion leaf (Definition~\ref{def:q-torsion-leaf}), a resonance
condition is satisfied: the $q$-harmonic forms
(Definition~\ref{def:q-harmonic-form}) on the brane constructively interfere
with the phonon modes of the crystal, amplifying vacuum fluctuations at
specific frequencies.  The resonance frequency $\nu_{\mathrm{res}}$ is
predicted to depend on the local value of $q$ and the lattice geometry.
\end{remark}

\begin{remark}[Substrate Telescopes]
\label{rem:substrate-telescope}
The directional dependence of $q$-torsion suggests a new class of detector:
the \emph{Substrate Telescope}.  A large, high-purity single crystal is
slowly rotated while the variation of vacuum noise power or phonon resonance
efficiency is recorded as a function of orientation.  Because the $q$-torsion
leaves carry a preferred direction inherited from the $q$-deformed Reeb flow
(Definition~\ref{def:q-reeb-flow}), the noise spectrum should exhibit a
modulation correlated with the crystal's orientation relative to the bulk
bundle.  This is conceptually analogous to a Weber-bar gravitational wave
detector, but targets the anisotropy of the substrate parameter rather than
metric strain.
\end{remark}

\paragraph{Existing data sets.}
\begin{itemize}
  \item Many-body interference measurements in kagome
        crystals~\cite{guo2025}, which probe coherent lattice-scale quantum
        effects.
  \item Chiral phonon spectroscopy in rare-earth halide
        crystals~\cite{luo2023}, demonstrating orientation-dependent
        effective magnetic fields from lattice vibrations.
  \item High-$Q$ quartz bulk acoustic wave (BAW) oscillator noise
        measurements, which record excess phase noise beyond thermal limits.
  \item Precision Casimir-effect experiments using crystalline surfaces
        (e.g.\ silicon and gold single crystals).
  \item Resonant-mass gravitational wave detector data (e.g.\ AURIGA,
        NAUTILUS bar detectors), which recorded orientation-dependent noise
        floors in cryogenic crystalline bars.
\end{itemize}

\paragraph{Planned observations.}
\begin{itemize}
  \item Next-generation phononic crystal experiments with tunable lattice
        constants, enabling a systematic scan of the resonance condition
        across different effective $q$ regimes.
  \item Extension of kagome and pyrochlore quantum-interference experiments
        to larger single crystals with controlled rotational staging.
  \item Cryogenic SQUID-based noise spectroscopy on rotating high-purity
        sapphire or silicon crystals.
\end{itemize}

\paragraph{Potential new experiments.}
\begin{itemize}
  \item A dedicated Substrate Telescope prototype: a cryogenically cooled,
        slowly rotating single crystal (e.g.\ ultra-pure silicon or quartz)
        instrumented with SQUID magnetometers and phonon sensors, searching
        for orientation-dependent modulation of the vacuum noise floor.
  \item Comparative resonance measurements across crystals with different
        lattice symmetries (cubic, hexagonal, trigonal) to map the
        dependence of $\nu_{\mathrm{res}}$ on lattice geometry and thereby
        constrain the local value of $q$.
  \item Re-analysis of archival resonant-bar detector data (AURIGA,
        NAUTILUS, EXPLORER) for residual orientation-dependent signals
        consistent with $q$-torsion anisotropy rather than gravitational
        wave backgrounds.
\end{itemize}
 \chapter{Speculations}
\label{chap:speculations}


\bibliographystyle{plain}
\bibliography{references}

\end{document}
